\documentclass[twoside]{article}
\hfuzz=\maxdimen
\usepackage[
    a5paper,
    margin=0.5in,
    includehead,
    headsep=4pt,
    headheight=13pt
]{geometry}
\usepackage{gregosheet}
\usepackage{gregosheet-psalm}
% Font sizes
\newcommand{\musicfontsize}{20}
\newcommand{\lyricfontsize}{10}
\newcommand{\psalmfontsize}{10}

% Fonts
% Change this to any font command you want:
% \rmfamily (roman/serif), \sffamily (sans-serif), \ttfamily (monospace)
% or use \setmainfont{FontName} or \newfontfamily for specific fonts
\newfontfamily\lyricfont{Cambria}
\newfontfamily\psalmfont[BoldFont={Cambria Bold}]{Cambria}
\newfontfamily\headingfont[BoldFont={Cambria Bold}]{Cambria}
\setmainfont{Cambria}

% Spacing
\newcommand{\systemvskip}{10pt}
\newcommand{\lyricvskip}{2pt}
\newcommand{\blockvskip}{10pt}

% Page layout
\usepackage{fancyhdr}
\pagestyle{fancy}
\fancyhf{}
\fancyhead[CO]{\textcolor{red}{\MakeUppercase{\leftmark}}}
\fancyhead[RO]{\thepage}
\fancyhead[LE]{\thepage}
\fancyhead[CE]{\textcolor{red}{\MakeUppercase{\rightmark}}}
\renewcommand{\headrulewidth}{0pt}
\setlength{\parindent}{0pt}

% Title style
\makeatletter
\renewcommand{\maketitle}{
  \begin{titlepage}
    \centering
    \vspace*{\fill}
    {\headingfont\fontsize{32}{38}\selectfont\bfseries\MakeUppercase{\@title}\par}
    \vskip 2em
    {\headingfont\fontsize{16}{20}\selectfont\MakeUppercase{\@subtitle}\par}
    \vspace*{\fill}
  \end{titlepage}
}
\newcommand{\subtitle}[1]{\gdef\@subtitle{#1}}
\newcommand{\@subtitle}{}
\makeatother

% Section style
\usepackage{titlesec}
\titleformat{\section}
  {\centering\headingfont\fontsize{20}{24}\selectfont\bfseries}
  {}{0pt}{\MakeUppercase}
\titlespacing*{\section}{0pt}{\blockvskip}{\blockvskip}

% Subsection style
\titleformat{\subsection}
  {\centering\headingfont\fontsize{16}{20}\selectfont\bfseries\addfontfeature{LetterSpace=10.0}}
  {}{0pt}{\MakeUppercase}
\titlespacing*{\subsection}{0pt}{\blockvskip}{\blockvskip}

% Subsubsection style
\titleformat{\subsubsection}
  {\centering\headingfont\fontsize{12}{14}\selectfont\bfseries\scshape}
  {}{0pt}{}
\titlespacing*{\subsubsection}{0pt}{\blockvskip}{\blockvskip}

% Part title command
\newcommand{\parttitle}[2][]{
  \vskip\blockvskip
  \par\noindent
  \headingfont\fontsize{10}{12}\selectfont
  \textcolor{red}{\MakeUppercase{#2}}%
  \ifx&#1&\else\hfill\textcolor{red}{#1}\fi
  \nopagebreak
  \par
}

% Rubric command
\newcommand{\rubric}[1]{
  \par\noindent
  {\fontsize{8}{10}\selectfont\textcolor{red}{#1}}
  \par
}

% Red text command
\newcommand{\redtext}[1]{
  \par\noindent
  \textcolor{red}{#1}
  \par
}

\newcommand{\p}{\textcolor{red}{P.: }}

\newcommand{\h}{\textcolor{red}{H.: }}

\newcommand{\e}{\textcolor{red}{E.: }}

\newcommand{\lec}{\textcolor{red}{L.: }}

\newcommand{\diac}{\textcolor{red}{D.: }}

% Red dash command
\newcommand{\reddash}{\textcolor{red}{\textemdash{}}\xspace}
\usepackage{xspace}

% Lectio environment
\newenvironment{lectio}[1]{
  \par\noindent
  {\centering\textcolor{red}{\textit{#1}}\par}
  \setlength{\parindent}{0.2in}
  \everypar{\setlength{\parindent}{0.2in}}
}{
  \par
}

% Versicle-Response command
\newcommand{\versicle}[2]{
  \par\noindent
  \hangindent=1em\hangafter=1
  \emergencystretch=1em
  \makebox[1em][l]{\textcolor{red}{V.}}#1
  \par\noindent
  \hangindent=1em\hangafter=1
  \emergencystretch=1em
  \makebox[1em][l]{\textcolor{red}{F.}}#2
  \par
}

% Hymn stanzas environment
\usepackage{multicol}
\newenvironment{hymnstanzas}{
  \setlength{\leftmargin}{0pt}
  \setlength{\rightmargin}{0pt}
  \setlength{\columnsep}{4em}
  \begin{multicols}{2}
  \raggedright
}{
  \end{multicols}
}

\newcommand{\stanzabreak}{%
  \par\addvspace{0.8\baselineskip}%
}


% Single centered stanza (for odd last stanza)
\newcommand{\centerstanza}[2][0.4]{
  \par
  \begin{center}
  \begin{minipage}{#1\textwidth}
  \raggedright
  #2
  \end{minipage}
  \end{center}
  \par
}

% Structured text environment
\newenvironment{preces}[2]{
  \par\noindent
  #1
  \par\noindent
  \hskip0.2in\textit{#2}
  \par
  \setlength{\parindent}{0pt}
  \everypar{\hangindent=0.2in\hangafter=1}
}{
  \par
}

\newcommand{\printsheet}[1]{%
  \ifcsname sheet@#1\endcsname
    \csname sheet@#1\endcsname
  \else
    \textbf{[Missing sheet: #1]}%
  \fi
}

\usepackage{ragged2e}


% ============================================================================
% I. PROPRIUM DE TEMPORE
% ============================================================================

% ----------------------------------------------------------------------------
% I. Tempus Adventus
% ----------------------------------------------------------------------------

% ----------------------------------------------------------------------------
% a. Usque ad diem 16 decembris
% ----------------------------------------------------------------------------

% --- Pro cunctis diebus ---

% ··· Invitatorium ···

\defsheet{RegemVenturum}
  {<-3-3r-3----:-3--3--3--4---3--4t-5---,-5---3--4t---4-4---3-.}
  {  A Ki-rályt, az el-jö-ven-dő U--rat * jöj-je-tek, i-mád-juk!}

% --- Hebdomada 1 Adventus ---

% --- Dominica ---

% ··· I Vesperas ···

% Antiphona 2 (Adv43)
\defsheet[tone={5-c}]{EcceDominusVeniet}
  {<-7-5--7--4u-5--4--3----:-4----4--3---5---7i----7--7-,-5--7--7----8-6--7--tT4-:-3----4--4----4t---5--4t-3--3-?,}
  {  Í-me az Úr el-jö-vend * s_ve-le min-den szent-je-i,  és lé-szen a-ma na-pon   nagy fé-nyes-ség, al-le-lu-ja.}

% Magnificat (Adv40)
\defsheet[tone={1la-a}]{EcceNomenDomini}
  {<X-1-1---3--1--0---34t-5--:-4-6--5---4---4--3t-4-,-4--4--4-2--4---tT4-3-:-2--3---4--eE2-1---1-?,}
  {   Í-me, az Úr-nak ne--ve * a tá-vol-gól kö-ze-leg és az Ő ra-gyo-gá--sa  be-töl-ti a   föl-det.}

% ··· Laudes ···

% Antiphona 1 (Adv42)
\defsheet[tone={8szo-c}]{InIllaDie}
  {<-4-4--4t-4---:-£------------------5u---4-5---rR3-3--,-3r-5-7---4-:-4---4---5---4---eE2--1--:-2--3r-4--4-?,}
  {  A-ma na-pon * megcsordulnak_a_he-gyek é-des-ség-gel, és a dom-bok tej-jel-méz-zel foly-nak, al-le-lu-ja.}

% Antiphona 2 (Adv60)
\defsheet[tone={1la-a}]{MontesEtColles}
  {<X-1--3----4--3t--5---:-5-4--2---3---4----eE2-1--1--,-1--2--3---4---5---4---3r-4-:-5---7---4---5--5u--5--,-5----tT4-2---3--rR3-1--3-1---1--0-:-1-0---3--3---2--3---4--eE2-1--1-?,}
  {   He-gyek és hal-mok * é-ne-kel-nek majd az  Úr-nak, az er-dők-nek min-den fá-i   tap-sol-nak ke-zük-kel, mert el--jön az Úr, az U-ral-ko-dó  ö-rök-or-szá-gá-ba, al-le--lu-ja.}

% Antiphona 3 (Adv44)
\defsheet[tone={4-a}]{EcceVeniet}
  {<-1-2---4t-5--:-4t-6----tT4-5--5-,-2--5----4---Ï3-wW1-wW1-:-0--1--2---2--–3---4--3--2--2-?,}
  {  Í-me, el-jő * a  nagy pró-fé-ta, ki majd meg-új-ít--ja    Je-ru-zsá-le-met, al-le-lu-ja.}

% Benedictus (A) (Adv85)
\defsheet[tone={1re-a}]{DiesDominiSicutFur}
  {<-3--0--1---1tu-5-:-5----5-5u--5----:-[------------tT4¨5uj5T4-:-4--2---3---4-eE2--1--1---,-3----2-1----1-1--1---2---3--4t-4--3--:-2--3---4--eE2-1---0q-1-?,}
  {  Az Úr-nak nap-ja  mint a tol-vaj, * éjszaka_érke-zik,         le-gye-tek a-zért ké-szen, mert a-mely ó-rá-ban nem is vé-li-tek, el-jön az Em--ber-fi-a.}

% ··· II Vesperas ···

% Antiphona 1 (Adv43)
\defsheet[tone={8szo-c}]{Iucundare}
  {<-3--4--5----4uiI7U6¨7iI7i-iI7-:-uU6-5--uU6-tT4-4--,-4--4----4--eE2-1-:-3--3--3---4t--4--3r-rR3--2--3r-4--4-?,}
  {  Vi-ga-dozz im------------már * Si--on le--á---nya, uj-jong re-pes-ve, Je-ru-zsá-lem le-á--nya, al-le-lu-ja.}

% Magnificat (A) (Adv26)
\defsheet[tone={8szo-c}]{SpiritusSanctus}
  {<-1-3-----3--3---4t-4---:-4--4-----5--4---3--4t-5-,-5--556uU6-tT4-4--:-4----4---5--4---3--eE2-1----:-1--3--4---5---4--4-?,}
  {  A Szent-lé-lek Is-ten * le-száll re-ád  Má-ri-a,  ne félj   te--hát, mert kit mé-hed-be fo--gadsz, az Is-ten-nek Fi-a.}

% Magnificat (B) (Adv49)
\defsheet[tone={7do-d}]{NeTimeasMaria}
  {<-8--6-----89p-oO8-8--:-©-----------------6---7---8--tT4-4--,-¦-------------5u-6----:-5----4--5-------7--6--4--4-?,}
  {  Ne félj, Má--ri--a, * mert_kegyelmet_ta-lál-tál az Úr--nál. Íme,_méhedben fo-gansz, s_fi-at szülsz, al-le-lu-ja.}

% Magnificat (C) (Adv64)
\defsheet[tone={1la-a}]{BeataEsMaria}
  {<-1t--5u--5---:-5--5u-5--:-5--5---rR3r-1--,-1----3---2--1--1---2---3---4tT4-3-:-2---3--4--3---2----1--qQ0-1--eE2-1--1-?,}
  {  Bol-dog vagy, Má-ri-a, * mi-vel hit--tél, mert tel-je-se-dik ben-ned mind-az, mit az Úr mon-dott né-ked al-le--lu-ja.}

% --- Feria II ---

% Benedictus (Adv57)
\defsheet[tone={1la-a}]{Leva}
  {<-1tu--5----:-7--5--4---4tT4-:-2-3----4---3---2--1-1--,-1--34t--4-:-5--4--3----4--4---5--,-1-2--3---4--4--4--eE2-:-1----3---eE2-1---1-3--eE2-1----1--?,}
  {  Nézz fel, * Je-ru-zsá-lem,   e-meld föl sze-me-i-det, és lásd meg ki-rá-lyod ha-tal-mát, í-me jön Üd-vö-zí-tőd,  hogy fel-old-jon a bi-lin-csek-ből.}

% Magnificat (Adv51)
\defsheet[tone={1la-a}]{AngelusDomini}
  {<-1--3--3---rR3-1---qQ0-:-3--4----5---4--3--3r-3-,-3---4--3--2---1e-eD0-:-2e-4-----eE2-1---qQ0--1--eE2-1--1-?,}
  {  Az Úr-nak an--gya-la  * kö-szön-töt-te Má-ri-át, s_ő mé-hé-ben fo-gant  a  Szent-lé--lek-től, al-le--lu-ja.}

% --- Feria III ---

% Benedictus (Adv90)
\defsheet[tone={1la-a}]{EgredieturVirga}
  {<-0----1---1tu-5---:-[--------------tT4¨5uj5T4-:-5--4---3--3-3r--3-:-3--2--3---4--eE2-1--1-,-3-----2---1--2--3---4t--4---3--:-3--2--3---4--eE2-1--1-?,}
  {  Vesz-sző sar-jad * Jessze_gyökeré-ből,         be-töl-ti a föl-det az Úr-nak di-cső-sé-ge, s_meg-lát-ja ak-kor min-den test az Úr-nak üd-vös-sé-gét.}

% --- Feria IV ---

% Benedictus (Adv52)
\defsheet[tone={7la-d}]{IerusalemRespice}
  {<-4t-6--7i--8----:-8--9----ü---8--9---8-,-8--8---45z--6--5z-4--4-?,}
  {  Je-ru-zsá-lem, * te-kint Nap-ke-let-re, és lás-sál, al-le-lu-ja.}

% Magnificat (Adv53)
\defsheet[tone={2-f}]{DeSion}
  {<-1--1r-3---:-3---0-1---1---,-3--3--1-0q-1-:-3--0--1eE2-3---1-?,}
  {  Si-on-ból * jön a Tör-vény, az Úr i-gé-je  Je-ru-zsá--lem-ből.}

% --- Feria V ---

% Magnificat (Adv53)
\defsheet[tone={4-a}]{BenedictaTu}
  {<-1--2----4t---5--:-4t-6--5---4----5--5---,-2--5----4-Ï3-2--1---wW1-:-0---1---2r---2--2-?,}
  {  Ál-dott vagy te * az as-szo-nyok kö-zött, ál-dott a te mé-hed-nek   gyü-möl-cse, Jé-zus.}

% --- Feria VI ---

% Benedictus (Adv54)
\defsheet[tone={1la-a}]{EcceVenietDeus}
  {<-1tu-5---:-5u-5---5--4--3---4--1-,-3--2---1---3--3--4tT4-:-3--3---eE2-1--4---eE2-qQ0-:-1--eE2-1--1-?,}
  {  Í---me, * el-jön az Is-ten-Em-ber Dá-vid-nak há-zá-ból,   el-fog-lal-ni tró-nu--sát,  al-le--lu-ja.}

% Magnificat (Adv54)
\defsheet[tone={4-a}]{ExAegyptoVocavit}
  {<-1-2----4t--5---:-4---5---6--tT4-5--5-5--,-5---2----5----4--Ï3-wW1-wW1-:-0--1--–2e--4--3--2--2-?,}
  {  E-gyip-tom-ból * hív-tam az én  fi-a-mat, jön már, hogy üd-vö-zít-se    né-pe-met, al-le-lu-ja.}

% --- Sabbato ---

% Benedictus (Adv55)
\defsheet[tone={4-a}]{SionNoliTimere}
  {<-4t-5---:-4--5z---tT4-5--5---,-2-5--4-Ï3-2--1--wW1-:-0--1---2---–3---4--3--2--2?,}
  {  Si-on, * ne félj már to-vább, í-me a te Is-te-ned   el-jön hoz-zád, al-le-lu-ja.}

% --- Hebdomada 2 Adventus ---

% --- Dominica ---

% ··· Laudes ···

% Antiphona 1 (Adv59)
\defsheet[tone={1la-a}]{UrbsFortitudinis}
  {<X-1t-5-:-[---------7--5---4--5-----:-3----5---4---5--rR3-2e-4--eE2-1--1-,-5uj5-5-5-5i--8iK5-5-,-5---7----5--5---4---5--rR3-4-:-2----4--2----4--wW1-0--:-1e-2--1--1-?,}
  {   Si-on  erősséges vá-ros ne-künk, * mert biz-tos fa-la  az Üd-vö--zí-tő, és   ő a bás-tyá--ja, hát nyis-sá-tok meg ka-pu--it, mert ve-lünk az Is--ten, al-le-lu-ja.}

% Antiphona 2 (Adv44)
\defsheet[tone={7do-a}]{OmnesSitientes}
  {<ô-tg3¨5z-5---:-5--5zZ5-4--4----:-4---4---4--rR3-wW1-2---1--:-5--6---7iI7-zZ5-4--5uU6-5--:-4---2---4--5---4---eE2-1--:-2e-2e-1--1-?,}
  {   Ó,     mind, ti szom-ja-zók, * jer-tek az é---lő  víz-hez, ke-res-sé---tek az U----rat, míg meg-ta-lál-hat-já--tok, al-le-lu-ja.}

% Antiphona 3 (Adv60)
\defsheet[tone={3-c}]{EcceDominusNosterCumVirtute}
  {<-4-5---7--7--:-7--7--6---5---7--tT4-4--,-5----7----6--5-4---5---4-:-4---3--4--5--eE2-1-:-2--3r-2--2-?,}
  {  Í-me, az Úr * ha-ta-lom-mal ér-ke--zik, hogy szol-gá-i-nak sze-mét meg-vi-lá-go-sít-sa, al-le-lu-ja.}

% Benedictus (B) (Adv68)
\defsheet[tone={4-a}]{EcceMitto}
  {<-1-2--4--4t--5---:-4---5--6--tT4-5---5--5--,-2-5--4-Ï3--wW1-wW1-:-0--1--2---4---2-2-?,}
  {  Í-me el-kül-döm * hír-nö-kö-met hoz-zá-tok, a-ki u-tat ké--szít  az én szí-nem e-lőtt.}

% ··· II Vesperas ···

% Antiphona 1 (Adv58)
\defsheet[tone={1la-a}]{EcceInNubibusCaeli}
  {<-1--0--1--3---2--1-qQ0-:-3--4--5---4----tT4-3-:-3r---2--eE2-1---qQ0-:-1--eE2-1--1-?,}
  {  Ím az ég fel-hő-i-ben * ér-ke-zik majd az  Úr  nagy ha-ta--lom-mal,  al-le--lu-ja.}

% Antiphona 2 (Adv59)
\defsheet[tone={1la-a}]{EcceApparebit}
  {<-1t-5-:-[----------7--tT4t-:-3---4--5---eE2-1--,-5--5--4-----5--5u--5--:-5--5----55i-8iK5-5-,-5----7--5----4----5--rR3r-:-3-----3---4--eE2-qQ0-:-1--eE2-1--1-?,}
  {  Í--me, megjelenik az Úr,  * nem ál-tat té--ged, ha ha-laszt-ja jöt-tét, te csak vár-jad  őt, mert bi-zony-nyal el-jön,   s_nem kés-le-ke--dik,  al-le--lu-ja.}

% Magnificat (AB) (Adv53)
\defsheet[tone={8do-c}]{Veniet}
  {<-7---7---7-7iI7-7----:-¦-------6----5u-4--,-3--5---7---7--7----7i--7-:-7---7--6--5---55--4---5zZ5-4--3r-4-?,}
  {  Jön már u-tá---nam, * aki_erő-sebb ná-lam, ki-nek nem va-gyok mél-tó  még sa-ru-szí-ját sem meg--ol-da-ni.}

% --- Feria II ---

% Benedictus (Adv82)
\defsheet[tone={8szo-c}]{DicitDominus}
  {<-4---3----1--3ed134tT4-:-0--1--2---3r-4--5----4tT43r-rR3-,-5----7--4--5---4----5--4-:-4-3---4----5--4---4rR3E1-:-2--3r-4--4-?,}
  {  Így szól az Úr:       * Pe-ni-ten-ci-át tart-sa-----tok,  mert el-kö-zel-gett im-már a men-nyek or-szá-ga,      al-le-lu-ja.}

% --- Feria III ---

% Magnificat (Adv15)
\defsheet[tone={5-c}]{VoxClamantisInDeserto}
  {<-¢-----------4---3-:-5-7----7ui-7----:--8--8----8--7i--5-:-5--5--5---4t-3--,-3-4---4---5--3--3r--4-:-7--5--4----5--4--3---3-?,}
  {  A_Kiáltónak sza-va  a pusz-tá--ban: * „Ké-szít-sé-tek el  az Úr-nak út-ját, e-gye-nes-sé te-gyé-tek Is-te-nünk ös-vé-nye-it!”}

% --- Feria IV ---

% Benedictus (Adv61)
\defsheet[tone={8szo-c}]{SuperSoliumDavid}
  {<-4--4---1eE2-3r--:-4tT4-3--4tT4¨5uj5-,-4--4--5-7--4---5--4-:-4-4---eE2-1---3----4t-4---4rR3d1-:-2--3r-4--4-?,}
  {  Dá-vid trón-ján * ül   az Úr,         és az ő or-szá-gá-ban u-ral-ko--dik mind-ö--rök-ké,      al-le-lu-ja.}

% Magnificat (Adv68)
\defsheet[tone={4-a}]{SionRenovaberis}
  {<-4t-5---:-4--5----6----tT4-5-5----,-2---5--4---Ï3-2-1--wW1-:-0-1--2--4---2---2-?,}
  {  Si-on, * bi-zony majd meg-ú-julsz, meg-lá-tod az i-ga-zat,  a-ki el-jön hoz-zád.}

% --- Feria V ---

% Benedictus (Adv84)
\defsheet[tone={5-c}]{PonentDominoGloriam}
  {<-¦-----------uU6-5-:-5-5--4---tT4-3---,-4--3--5-77i-7-:-8-ö--8--¦-----------7i--5--,-5--8---7-6u-tT4-5-:-4--4t--4t--3--3-?,}
  {  Megadják_az Úr--nak a di-csé-re--tet * és az ő hí--rét a tá-vo-li_szigetek zen-gik, mi-vel í-me el--jő  és már nem ké-sik.}

% Magnificat (Adv69)
\defsheet[tone={4-a}]{QuiPostMeVenit}
  {<-1-2--4-4--5z--5----:-4-5----zh4--5--5--,-2---5--4----Ï3-wW1-wW1-:-0---1--2--3--4--–3-2---2-?,}
  {  A-ki u-tá-nam jön, * e-lőbb volt ná-lam, nem va-gyok én mél-tó    meg-ol-da-ni sa-ru-szí-ját.}

% --- Feria VI ---

% Benedictus (Adv69)
\defsheet[tone={1re-a}]{DicitePusillanimes}
  {<-3----4t--5---:-7--5-----4--2e-4-:-4-3--2---1ed1-0q--1--,-4-4---2--4tT4-3--:-3-4--2-eE2--1--1-?,}
  {  Mond-já-tok: * ki-csiny-hi-tű-ek, e-rő-söd-je---tek meg, í-me, az Is---ten, a mi U-runk el-jő.}

% --- Sabbato ---

% Benedictus (Adv70)
\defsheet[tone={1la-a}]{Levabit}
  {<-1----0--1---2e-4-----eE2--1--qQ0-:-3-5---4--3---4-3---,-3--3---4---tT4--3-:-3r---2-3----2-----1--0q-1-?,}
  {  Nagy je-let tá-maszt majd az Úr  * a nem-ze-tek e-lőtt, és ösz-sze-gyűj-ti  mind a szét-szórt Iz-rá-elt.}

% --- Hebdomada 3 Adventus ---

% ··· I Vesperas ···

% Antiphona 1 (Adv74)
\defsheet[tone={8szo-c}]{IerusalemGaude}
  {<-1--2--3---4----eE2-1-----:-3----3-4t--4--,-5u-7-:-5u-5---4----5---4-:-3---4--eE2-1--:-3--4t-4--4-?,}
  {  Je-ru-zsá-lem, ör--vendj * nagy ö-röm-mel, í--me, el-jön majd hoz-zád sza-ba-dí--tód, al-le-lu-ja.}

% Magnificat (Adv72)
\defsheet[tone={1la-a}]{AnteMe}
  {<X-1---3----1--0---3-4t--5----:-6-----5----4---4-3t-4--,-4--4---4--2---4---5---4---3-:---2----3r--eE2--1---0q--1-?,}
  {   Nem volt Is-ten e-lőt-tem, * s_nem lesz más u-tá-nam, ne-kem ha-jol meg min-den térd, s_en-gem vall min-den nyelv.}

% ··· Laudes ···

% Antiphona 1 (Adv73)
\defsheet[tone={1re-a}]{VenietDominus}
  {<-5--4---4--5--4----:-3---4--5-----eE2-1---1--:-1---1--2--3--4t--4-:-4-5--4---3---4t--3--3--,-3--4---5---4---5---4--2--3---4---2---3---qQ0-:-1--eE2-1--1-?,}
  {  El-jön az Is-ten, * nem ha-laszt-ja  jöt-tét, meg-vi-lá-go-sít-ja  a sö-tét-ség rej-te-két, ki-nyi-lat-koz-tat-ja ma-gát min-den nép-nek,  al-le--lu-ja.}

% Antiphona 2 (Adv75)
\defsheet[tone={5-c}]{MontesEtOmnesColles}
  {<-¢-----------4---3---:-3---3-5---7i--7--:-7-7---7--7i-8-:-7--8-ö---9--iI7-7--:-8---7i-5-:-5--5--5--4t-3---,-3----3r-4t---3r-4---:-7--7---4--tT4-3-?,}
  {  A_hegyek_és hal-mok * meg-a-láz-kod-nak, a gör-be-sé-gek ki-e-gye-ne-sed-nek, s_a rö-gök út-tá si-mul-nak, jöjj el hát, U--runk, ne kés-le-ked-jél!}

% Antiphona 3 (Adv74)
\defsheet[tone={1la-a}]{DaboInSionSalutem}
  {<X-1---3-1---0--3--4---3--4t-5----:-6----5--4---4---4---3--4---5--eE2-qQ0-:-1--eE2-1--1-?,}
  {   Meg-a-dom az Üd-vöt Si-on-nak, * s_Je-ru-zsá-lem-nek di-cső-sé-ge--met,  al-le--lu-ja.}

% ··· II Vesperas ···

% Antiphona 3 (Adv74)
\defsheet[tone={1la-a}]{IusteEtPieVivamus}
  {<-1-0--1---2e----4--eE2-1--1---,-0---1--3---1-:-1-0---1---4--3----:-0q-3--3---0e--1-?,}
  {  I-ga-zul s_jám-bo-ran él-jünk, vár-va vár-juk a bol-dog re-ményt, az Úr-nak jöt-tét.}

% Magnificat (C) (Adv68)
\defsheet[tone={8do-c}]{VenietFortior}
  {<-7---7---¨-7-7iI7-7--:-¦-------6----5u-4--,-3--5---7---7--7----7i--7-:-7---7--6--5---55--4---5zZ5-4--3r-4-?,}
  {  Jön már * u-tá---nam, aki_erő-sebb ná-lam, ki-nek nem va-gyok mél-tó  még sa-ru-szí-ját sem meg--ol-da-ni.}

% --- Feria II ---

% Benedictus (Adv66)
\defsheet[tone={8szo-c}]{DeCaeloVeniet}
  {<-4-7-----7---8--9--8----:-9-7---8--7--4--4--4---,-7--7--8-7--5u-7-:-7--7iI7-tT4-3--4--4-?,}
  {  A menny-ből el-jö-vend * U-ral-ko-dó Is-te-nünk, és az ő ke-zé-ben or-szág és  ha-ta-lom.}

% Magnificat (Adv82)
\defsheet[tone={8szo-c}]{BeatamMeDicent}
  {<-5u--5---4---5---4--4---:-5--4---5---rR3-4---4t-3--3--,-4----4-4--3--5---7----7--8--7iI7-:-5---5---7--7--4--4-?,}
  {  Bol-dog-nak mon-da-nak * en-gem min-den nem-ze-dé-kek, mert a-lá-za-tos szol-gá-ló-ját    meg-lát-ta az Is-ten.}

% --- Feria III ---

% Benedictus (Adv83)
\defsheet[tone={4-a}]{ElevareElevare}
  {<-1----2----4t---5----:-4-5---6-----5--4--5---5--,-2----5--4---Ï3-wW1-wW1-:-0--1--2---4--2-2-?,}
  {  Kelj föl, kelj föl, * e-mel-kedj, Je-ru-zsá-lem, vesd le már i--gá--dat,  Si-on rab le-á-nya.}

% Magnificat (Adv52)
\defsheet[tone={1la-a}]{AntequamConvenirent}
  {<X-1--1-3----1---0--34t-5----:-6---5--£------------3t-4-,-£--------------2---4---5--4--2-:-2e-4-----eE2-1---qQ0--1--eE2-1--1-?,}
  {   Mi-e-lőtt egy-be-kel-tek, * úgy ta-láltatott_Má-ri-a,  hogy_gyermeket hor-doz mé-hé-ben a  Szent-lé--lek-től, al-le--lu-ja.}

% --- Feria IV ---

% Benedictus (Adv85)
\defsheet[tone={2-f}]{Consolamini}
  {<-3--3----3--eE2-1----:-1--0----1--3rR3-3----1--0----1----3rR3-3--:-eE2-qQ0--1--eE2-1--1-?,}
  {  Vi-gasz-ta-lód-jál, * vi-gasz-ta-lód--jál, vá-lasz-tott né---pem, így szól az Úr--is-ten.}

% Magnificat (Adv68)
\defsheet[tone={7do-d}]{TuEsQuiVenturusEs}
  {<-8--8----6--8o-8--7iI7U6-4---:-5--5-----7-7--:-5---6---5--4--4---,-£-------------3-----4t-eE2-1--:-3--4t-4--4-?,}
  {  Te vagy az el-jö-ven----dő, * te vagy, U-ram, kit vár-va vá-runk, hogy_megszaba-dítsd né-pe--det, al-le-lu-ja.}

% --- Feria IV ---

% Benedictus (Adv82)
\defsheet[tone={2-f}]{ConsurgeConsurge}
  {<-1r-3---:-eD0q-1-----:-3----1---3-3---1--0q-1--,-3-0----1--3---2---1---1-?,}
  {  In-dulj, in---dulj, * vedd fel e-rős-sé-ge-det, U-runk ha-tal-mas kar-ja!}

% --- Feria IV ---

% Magnificat (Adv84)
\defsheet[tone={1la-a}]{HocEstTestimonium}
  {<X-1--qQ0-3--4--3---4t-5----:-6--5----4--4---3t--4---,-4--4-2--4---tT4-3---:-2-3r---eE2--1--1-?,}
  {   Ez a   ta-nú-ság-té-tel, * me-lyet Já-nos mon-dott: ki u-tá-nam jön majd, e-lőbb volt ná-lam.}

% ----------------------------------------------------------------------------
% b. Post diem 16 decembris
% ----------------------------------------------------------------------------

% --- Hebdomada 4 Adventus ---

% ··· I Vesperas ···

% Antiphona 1 (Adv119)
\defsheet[tone={1la-a}]{EcceVenietDesideratus}
  {<X-1-1--3-1----0--34t-5------:-6---5---4--4---4--3t-4--,-£----------2----4---tT4-3-:-3--3--4---eE2-qQ0-:-1--eE2-1--1-?,}
  {   Í-me a vágy-va vá--gyott, * min-den-né-pek re-mé-nye, betelik_fé-nyes-ség-gé--vel az Úr-nak há--za,   al-le--lu-ja.}

% Antiphona 2 (NZS26)
\defsheet[tone={4-a}]{VeniDomine}
  {<-1----2---4-4t-5-----:-4--5---6--5---4--5x--5x-,-5--2----5---4--Ï3---2--1--w\\W1\\-:-0--1---–2e-4--3--2y-2y-?,}
  {  Jöjj el, ó U--runk, * és már ne kés-le-ked-jél, ol-dozd föl go-nosz te-te-it        né-ped-nek Iz-ra-el-nek.}

% Antiphona 3 (NZS232)
\defsheet[tone={4-a}]{EcceIamVenit}
  {<-1-2--4---4t--5---:-5--4-5---6---tT4-5--5-,-2----5---4--Ï3--2--1---w\\W1\\-:-0-----1--2--4-2----2-?,}
  {  Í-me már itt van * az i-dők tel-jes-sé-ge, mely-ben Is-ten el-kül-dé        szent Fi-át a föld-re.}

% ··· Laudes ···

% Antiphona 1 (Adv118)
\defsheet[tone={1la-a}]{CaniteTuba}
  {<-3----0--1----1tu-5----:-5--5--tT4¨5uj5T4-:-4----5--4---33--3--23r-eE2-1---1-,-3-2---1---1---1---2--3--4tT4-3--:-3--qQ0-1--4---2--eE2-1--1-?,}
  {  Szól-ja-tok, kür-tök, * Si-on-ban,         mert kö-zel van az Úr--nak nap-ja. Í-me, jön meg-men-te-ni min--ket, al-le--lu-ja, al-le--lu-ja.}

% Antiphona 2 (Adv120)
\defsheet[tone={1la-a}]{DominusVeniet}
  {<-qQ0-3---4t-5---:-4----7---5--4---7-55-1---3----1--,-5uj55-4----7--7--7uU6Z5-5--:-5--4---3--4t--4-:-5---4--3---4--1--,-4-rF1-3--1----3-3---2e-qQ0-:-0--1--eE2-3r-4--:-0--1--3r-2e--4--eE2-1--1-?,}
  {  Jön már az Úr, * most fus-sa-tok e-lé ezt mond-ván: ó,    nagy Fe-je-de-----lem, ki-nek or-szá-ga  nem ér-het vé-get! E-rős Is-ten, U-ral-ko-dó,   bé-ke Ki--rá-lya, al-le-lu-ja, al-le--lu-ja.}

% Antiphona 3 (Adv121)
\defsheet[tone={2-f}]{Omnipotens}
  {<-3---3---3--1--¨-1-0----1e-3--3--0---,-2-3--4--3--1---3---3---2e--4-eE2-qQ0--:-1--eE2-1--1-?,}
  {  Min-den-ha-tó * I-géd, ó  Is-te-nünk, a te ki-rá-lyi szé-ked-ből a-lá--száll, al-le--lu-ja.}

% Benedictus (B) (Adv68)
\defsheet[tone={7la-d}]{QuomodoFietIstud}
  {<-4--7---7---8--7---5u--:-¦-----------4---5uj5-:-7--7---8o--8oO8-:-7----7--5--4--5---7--:-7---4--3r-4--,--8o---8---9õ-õÔö-oO8-9--8---9---8oO8-:-8-9-----8--7---7i-7---8--7-----5u-7-,-6--7-8--7---7---uU6-5--uU6-tT4-4-:-5---5--5---zZ5-4--4-?,}
  {  Mi-kép-pen le-het ez, * Istennek_an-gya-la,    mi-vel fér-fit,   hogy tő-le fo-gan-jak, nem is-me-rek? „Hall-jad Má-ri--a,  Is-ten szü-ze,    a Szent-lé-lek Is-ten le-száll re-ád, és a Ma-gas-ság-be--li e---re--je  meg-ár-nyé-koz té-ged.”}

% Benedictus (C) (Adv85)
\defsheet[tone={4-a}]{ExQuoFactaEst}
  {<-1--2-----4---4t--5---:-4--5----6--tT4-5---5--,-2t--ÏrR3-wW1-wW1-:-0--1---2---4-2----2-?,}
  {  Mi-helyt hal-lot-tam * kö-szön-té-sed sza-vát, fel-uj---jon-gott  mé-hem-ben a gyer-mek.}

% ··· II Vesperas ···

% Antiphona 2 (Adv120)
\defsheet[tone={1re-a}]{EruntPravaInDirecta}
  {<-3-0---1--1t-5u-:-[---------tT4--5----:-3r-5--4---2---3---4--3--2--1---1--,-3----4---3--3-0--:-2e-4--3---3--eE2-qQ0-:-1--eE2-1--1-?,}
  {  A gör-be-sé-gek  egyenessé lesz-nek, * a  rö-gös föl-dek út-tá si-mul-nak, jöjj el, Ó, U-runk és ne kés-le-ked-jél,  al-le--lu-ja.}

% --- Antiphonae ad Laudes et Vesperas a die 17 ad diem 23 decembris ---

% ··· Feria II ···

% Antiphona 1 (NZS231)
\defsheet[tone={4-a}]{EcceVenietDominus}
  {<-1-2--4--4t-5---:-4---5---6--5--4----5-5-,-2---5----4-Ï3-2--1----w\\W1\\-:-0-1-2--4--2---2-?,}
  {  Í-me jő az Úr, * min-den ki-rá-lyok U-ra, bol-dog, a-ki ké-szen áll      ő-e-lé-be fut-ni.}

% Antiphona 2 (Adv69)
\defsheet[tone={7la-d}]{CantateDominoCanticum}
  {<-4--7--7---8----9--8----:-8--8-9--7---5--7--7--,-7--7---5--7--iI7--7-:-7-7iI7-5---5---4--4-?,}
  {  Az Úr-nak zeng-je-tek, * új é-ne-ket az Úr-nak, di-csé-re-te szól-jon a föld min-den vé-gén.}

% ··· Feria III ···

% Antiphona 1 (NZS233)
\defsheet[tone={4-a}]{EgredieturDominus}
  {<-1--2---4t-5--:-4--5-6-----tT4-5---5--,-2---5----4---Ï3-wW1-w\\W1\\-:-0---1---2--4--2--2-?,}
  {  Ki-lép az Úr * az ő szent he--lyé-ről, jön mint sza-ba-dí--tó,       meg-men-te-ni né-pét.}

% --- Antiphonae ipsis diebus decembris a 17 ad 24 assignatae ---

% ··· Die 17 decembris ···

% Benedictus (Kar16)
\defsheet[tone={8szo-c}]{Scitote}
  {<-5---3--3r--4--:-4----5--4---3---4--4---5--3y--3y-,-3--5----7i--6---7---5--:-4---5---4---3--eE2-1--:-3--4t-4x-4x-?,}
  {  Tud-já-tok meg, hogy kö-zel van Is-ten or-szá-ga;  bi-zony mon-dom nék-tek: már nem kés-le-ke--dik, al-le-lu-ja.}

% Magnificat
\defsheet[tone={2-f}]{OSapientia}
  {<-1e3+-ed1-0q---1----:-Ÿ----------------------3---1e---1--0í-,-Ÿ-----------0q-1-:-Ÿ-------------3--1--1---3-3--1---3--3--4tT4¨5uj5-,-5-----5--4---3--3--2---1e-0í-,-0e---3-:-3----3----1---0q-1qQ0Ôð-:-0q-3-3---0---0e-1í-?,}
  {  Ó    Böl-cses-ség, * ki_a_Fölségesnek_szájá-ból szár-ma-zol, ki_a_minden-sé-get egyik_végétől má-si-kig e-rő-sen át-fo-god,        s_sze-lí-den el is ren-de-zed: Jöjj el, s_mu-tasd meg né-künk     az o-kos-ság út-ját!}

% ··· Die 18 decembris ···

% Benedictus (Adv98)
\defsheet[tone={1la-a}]{VigilateAnimo}
  {<-1---1----0---1---3---2--1---qQ0--,-3--4---tT4-3-:-3r-2---eE2-1--1-?,}
  {  Vir-rasz-sza-tok lel-ke-tek-ben, * kö-zel van már az Úr, az  Is-ten.}

% Magnificat
\defsheet[tone={2-f}]{OAdonai}
  {<-1e3+-ed1-0q-1----:-Ÿ-----------3---1--1--0í-,-0--Ÿ------------------3---1----1---3---1--3---4tT4¨5uj55\%-,-5----4--¢-----------3---2-----1e-0í-,-0e---3y--1--3-----1---0q--1qQ0Ôð-:-0--1----3----0e--1í-?,}
  {  Ó    Á---do-náj, * Izrael_házá-nak ve-zé-re,  ki Mózesnek_a_vöröslő tűz-láng-ban meg-je-len-tél,           s_ne-ki Sion_hegyén tör-vényt ad-tál: Jöjj el, és válts meg min-ket,     ki-nyúj-tott kar-ral!}

% ··· Die 19 decembris ···

% Magnificat
\defsheet[tone={2-f}]{ORadixIesse}
  {<-1e3+-3----1---ed1-0q-1---:-Ÿ------------------3--1e--1--0í--,-0---1---1--:-Ÿ----------3---1--3---4tT4¨5uj55\%-,-5--5--5-5--4---3--eE2-1e--0í-,-0e---33+-1--3---3--3----1---0q--1qQ0Ôð`-:-0q--3--3---0--0e--1í-?,}
  {  Ó    Jesz-sze gyö-ke-re, * ki_jelként_állsz_a né-pek kö-zött, kit lát-ván, a_királyok meg-né-mul-nak,           ki-re a né-pek vá-gya-koz-nak, Jöjj el, és sza-ba-díts meg min-ket,      már ne kés-le-ked-jél!}

% ··· Die 20 decembris ···

% Magnificat
\defsheet[tone={2-f}]{OClavisDavid}
  {<-1e3+-ed1-ed1-0q--1í---:-Ÿ-----------3---1--1--0í-,-1--0q--1-----"-1--3------1-1--3--3---4tT4¨5uj55\%-,-5--5------4--3------3--2---1e---0í-,-0e---33+-1--3----3--1-ed1-0q--1-"-Ÿ---------3--1---1--0í-,-1---3--3---1---0q-1qQ0Ôð`-:-0-0--1---3--0---0e-1í-?,}
  {  Ó    Dá--vid kul-csa, * Izrael_házá-nak jo-ga-ra   ki meg-nyitsz, és nincs, a-ki be-zár-ja,            be-zársz, és nincs, ki meg nyit-ja:  Jöjj el, és vidd ki a bör-tön-ből a_bilincs-be ver-te-ket, kik sö-tét-ben ül-nek,      a ha-lál ár-nyé-ká-ban.}

% ··· Die 21 decembris ···

% Magnificat
\defsheet[tone={2-f}]{OOriens}
  {<-1e3+-ed1ed1-0q-1í---:-1-1---3----1--1--0í-,-1-0---1---3---4tT4¨5uj55\%-,-5----5---5----4--3--3----2---1e--0í-,-1---3--3---1---0q-1qQ0Ôð`-:-0-0--1---3--0---0e-1í-?,}
  {  Ó    Nap----ke-let, * ö-rök fény su-ga-ra,  i-gaz-ság nap-ja,            Jöjj el, s_vi-lá-go-síts meg min-ket, kik sö-tét-ben ü-lünk,      a ha-lál ár-nyé-ká-ban.}

% ··· Die 22 decembris ···

% Magnificat
\defsheet[tone={2-f}]{ORexGentium}
  {<-1e3+-3--1---ed1-0q-1----:-Ÿ---------------1---3--1ed1Q0-0í-,-1--3----3---4tT4¨5uj55\%-:-5---5--4---3---3--2---1e-0í-,-0e---33+-1--3------3---1--0q-1qQ0Ôð-:-0---1-3---3---0--0e--1í-?,}
  {  Ó    né-pek Ki--rá-lya, * s_mindnyájuk_vá-gya-ko-zá-----sa,  te Szeg-let-kő,            két né-pet egy-be-fog-la-ló,  Jöjj el, és mentsd meg az em-bert,    kit a por-ból al-kot-tál!}

% ··· Die 23 decembris ···

% Magnificat
\defsheet[tone={2-f}]{OEmmanuel}
  {<-1e3+-ed1-ed1-0q-1í--:-1---1e-1-"---1--1---1----3--1ed1Q0í-,-3--3---1---3--3----4tT4¨5u-5x-:-5-----5--4---3--3--2--1e-0í-,-0e---3-!-1--3--3--1----0q--1qQ0Ôð`-:-0q-3-3y---0--0e-1í-?,}
  {  Ó    Em--má--nu-el! * Tör-vé-nyünk és Tör-vény-ho-zónk!     Né-pek-nek re-mény-sé------ge,  s_min-de-nek Üd-vö-zí-tő-je!  Jöjj el, és üd-vö-zíts min-ket,      mi U-runk Is-te-nünk!}

% --- Die 24 decembris ---

% ··· Laudes ···

% Antiphona 1 (Adv93)
\defsheet[tone={3-c}]{TuBetlehem}
  {<-4--5u--7--7----:-uU6-5u-7----4-:-5u--5----rR3-2---3r-4---,-5----¦-----------5-7--4--:-4--5---5--5----4--5--4--3----4--rR3-2-?,}
  {  Te Bet-le-hem, * Jú--da föld-je, nem vagy a   leg-ki-sebb, mert belőled_jön a Ve-zér, ki kor-má-nyoz-za né-pe-met, Iz-ra--elt.}

% Antiphona 3 (Adv131)
\defsheet[tone={8do-c}]{Crastina}
  {<-¦---------------¨-¦-----------5---7--6--5---4t-:-3---5----¦--------------5---6--5--4-?,}
  {  A_holnapi_napon * eljön_hozzá-tok az Üd-vös-ség, így szól a_seregeknek_U-ra, Is-te-ne.}

% Benedictus (Kar24)
\defsheet[tone={4-a}]{CompletiSunt}
  {<-4--4t--5---:-4--5--6-tT4-5---5--5-,-2t---ÏrR3-wW1--wW1-:-0--1--2---4----2--2y-?,}
  {  Be-tel-tek * Má-ri-á-nak nap-ja-i,  hogy meg--szül-je    el-ső szü-lött Fi-át.}

% ··· Tertia ···

% Antiphona (Adv129)
\defsheet[tone={8szo-c}]{IudaeaEtIerusalem}
  {<-4--4--rR3d1-3--4t-5--5---4----:-4--5u---4---4tT4R3-3-,-uU6-5--5i--6-7--tT4-4--:-4--rR3-1-:-1--2--3---4t---5--4rR3d1-:-2--3r-4--4-?,}
  {  Jú-de-a     és Je-ru-zsá-lem, * ne félj már töb----bé, jöj-je-tek e-lő hol-nap, és í---me, ve-le-tek lesz az Úr,      al-le-lu-ja.}

% ··· Sexta ···

% Antiphona (Adv129)
\defsheet[tone={8do-c}]{HodieScietis}
  {<-uU6-5--7---7---zZ5-4----:-4----5--4--34t-5-,-7--uU6-5---6---7---5--4-:-4--4-3--4t--4--4-?,}
  {  Még ma meg-tud-já--tok, * hogy el-jő az  Úr, és re--gel meg-lát-já-tok az ő di-cső-sé-gét.}

% ··· Nona ···

% Antiphona (Adv129)
\defsheet[tone={1la-a}]{CrastinaDie}
  {<X-1--1--3---1--0---34t-5---:-4-6----5---4--4----3t-4-,-4--4-4---2--4---tT4-3--:-3-2--3---4--3--2--1--1-?,}
  {   El-tö-röl-te-tik hol-nap * a föld-nek go-nosz-sá-ga, és u-ral-ko-dik raj-tunk a vi-lág Üd-vö-zí-tő-je.}

% ----------------------------------------------------------------------------
% II. Tempus Nativitatis
% ----------------------------------------------------------------------------

% ----------------------------------------------------------------------------
% a. Usque ad USQUE AD SOLLEMNITATEM EPIPHANIAE
% ----------------------------------------------------------------------------

% --- Pro Cunctis Diebus ---

% ··· Invitatorium ··· (NZS137)

\defsheet{ChristusNatusEstNobis}
  {<-33r---3-:-3--3---3---4--3----4t-5-----:-5---3--4t---4-4---3-.}
  {  Krisz-tus ma meg-szü-le-tett né-künk: * jöj-je-tek, i-mád-juk.}

% --- Die 25 decembris - In Nativitate Domini ---

% ··· Officium lectionis ···

% Antiphona 1 (NZS138)
\defsheet[tone={2-g}]{DominusDixit}
  {<ô-2--2y-¨-2---2--2r-2--:-2--4--44-2r---1í-,-3r-5--5----5--rR3-2r-2y-.}
  {   Az Úr * mon-dá né-kem: én Fi-am vagy te,  én ma szül-te-lek té-ged.}

% Antiphona 2 (Gondolkodó 4.)
\defsheet[tone={8szo-c}]{TamquamSponsus}
  {<-4----5--4--rR3-:-3--5--uU6--5u-zZ5-:-4--4-5----5--4--4-?,}
  {  Mint Vő-le-gény, ki-lé-pett az Úr    az ő nász-há-zá-ból.}

% Antiphona 3 (Gondolkodó 4.)
\defsheet[tone={1re-a}]{DiffusaEst}
  {<-3---4t--55--¨-2--3---4---eE2-1--1--,-3---4t-44---3--eE2-qQ0-1----eE2-1---1-?,}
  {  Ked-ves-ség * öm-lött el aj--ka-don: meg-ál-dott az Is--ten mind-ö---rök-ké.}

% ··· Laudes ···

% Antiphona 1 (Kar37)
\defsheet[tone={2-g}]{QuemVidistis}
  {<ô-4---4---4--2----2----4--4----:-1----2--4---3--2r---4--,-2--1--2---4--2---4--2---:-2--2--1--2-4---4---1---4--2---,-1w--4---4--4--2---4---2--:-2----1---2---4--3--2---4---4--,-4---5--3-4---2--1----2--:-1--2--4--4---1--1r-2y-2y-?,}
  {   Kit lát-ta-tok, pász-to-rok? * Mond-já-tok el gyor-san! Ad-já-tok hí-rül ne-künk, ki az ki a föl-dön meg-je-lent? Lát-tuk az Új-szü-löt-tet, s_an-gya-lok kó-ru-sát lát-tuk, kik az U-rat di-csér-ték, al-le-lu-ja, al-le-lu-ja.}

% Antiphona 2 (Kar40)
\defsheet[tone={7la-a}]{Angelus}
  {<ô-1--1t-5----:-[----------------4--56uU6Z5-5\\u\\-:-4-5---7----7-6--5---7---6--5---6---5--,-5--3---5---6---5--4----5-4---rR3-2e-:-1--3--4--wW1-0---1w-2--1í-1í-?,}
  {   Az an-gyal * így_szólt_a_pász-to-rok-----hoz:     í-me, nagy ö-rö-met hir-de-tek nek-tek, mi-vel meg-szü-le-tett e nap né--künk az Üd-vö-zí--tő, al-le-lu-ja.}

% Antiphona 3 (Kar45)
\defsheet[tone={8do-c}]{ParvulusFilius}
  {<-7--6-----5--7--4x-:-3---5--7----7-6--tT4-5--4x--,-7----7--7--8o---8v--:--9-8---7--8---uU6-5x--:-7--8--uU6-5-:-6--5z-4x-4x-?,}
  {  Ki-csiny-ke Fi-ú  * szü-le-tett e na-pon ne-künk, s_ne-ve ez lesz majd: „E-rős-sé-ges Is--ten”, al-le-lu--ja, al-le-lu-ja.}

% Benedictus
\defsheet[tone={8szo-c}]{Gloria}
  {<-3--4t--4---:-4-4--4---5---4---3--4t--5-,-5--6-7---6---5--4---4-:-4-4--4-4--5--4--tT4-3--eE2-1--:-3--4t-4x-4x-?,}
  {  Di-cső-ség * a ma-gas-ság-ban Is-ten-nek és a föl-dön bé-kes-ség a jó-a-ka-ra-tű em--be-rek-nek, al-le-lu-ja.}

% ··· II. Vesperas ···

% Antiphona 1 (Kar80)
\defsheet[tone={1re-a}]{TecumPrincipium}
  {<-3--0--1--1t-5u-5---:-[---------------tT45uj5T4-:-4-3-----4---55-rR3r-1--,-3--2---1---1----2--3---4tT4-3-:-3-3---2--3--4----eE2-1-1-?,}
  {  Ti-éd az el-ső-ség * hatalmadnak_nap-ján,        a szent-ség fé-nyé--ben, ma-gam-ból szül-te-lek Té---ged a haj-na-li csil-lag e-lőtt.}

% Antiphona 2 (Kar84)
\defsheet[tone={8szo-c}]{ApudDominum}
  {<-4--7\\-uJ4-:-6u-5--4---tT4-3y-,-3t-7--7--uJ4-6-7--5--5-6---5--4x-4x-?,}
  {  Az Úr--nál * az ir-gal-mas-ság, és bő-sé-ges ő-ná-la a sza-ba-dí-tás.}

% Antiphona 3 (Kar65)
\defsheet[tone={1la-a}]{InPrincipio}
  {<-1---1t--5----:-[-----------5u-5--:-5--5---4----3--4-1-,-3-2---1--1---1---2--3----4tT4-3---:-3-2--3---4--3--2--1í-1í-?,}
  {  Kez-det-ben, * minden_idők e--lőtt Is-ten volt az I-ge, Ő az, ki meg-szü-le-tett né---künk, a vi-lág Üd-vö-zí-tő-je.}

% Magnificat
\defsheet[tone={6-a}]{Hodie}
  {<X-3-4t-5---:-55T45z-5---4---3t-5x--,-3-3r-4-:-4--4--4--3--4t-zZ5-rR3-3y--,-3-4z-6-:-¥------------------4--zZ5T4-3t--5-:--3--4---6---6c--4--6---tT4-3t--5x-,-5-7--™7iI7uU6Z5-:-5--5---4--5---6--5-4--3t--5x--3r---4x-,-3--eE2-1í--1-2--3---4t--4x--3--eE2-q\\Q0\\-:-1--3--3y-3y-?,}
  {   E na-pon * Krisz--tus szü-le-tett, e na-pon az Üd-vö-zí-tő meg-je--lent, e na-pon a_földön_énekelnek az an----gya-lok, ör-ven-dez-nek az ark-an--gya-lok, e na-pon          új-jon-ga-nak az i-ga-zak is, mond-ván: Di-cső-ség a ma-gas-ság-ban Is-ten-nek,      al-le-lu-ja!}

% --- Dominica infra octavam Nativitatis - S. Familiae Iesu, Maria et Ioseph ---

% ··· I. Vesperas ···

% Antiphona 2 (Adv132)
\defsheet[tone={7do-d}]{IosephFiliDavid}
  {<-8--6-----8--9---8--8--:-©-------------6---7---8--5--4--4-,-¦-----------------5u-6---:-6-5-----4--5---uU6-4--4-?,}
  {  Jó-zsef, Dá-vid fi-a, * ne_félj_magad-hoz ven-ni Má-ri-át, mert_aki_tőle_szü-le-tett, a Szent-lé-lek-től va-ló.}

% ··· II. Vesperas ···

% Antiphona 3 (Kar122)
\defsheet[tone={7la-a}]{PuerIesus}
  {<ô-1-4----4---5z-5---:-rR3-2--4t-4-,--4---5---3-----4----wW1-2-:-2--qQ0-1---2--3--2--1---0q-1-?,}
  {   A gyer-mek Jé-zus * nö--ve-ke-dett kor-ban s_böl-cses-ség-ben az Is--ten és em-be-rek e--lőtt.}

% --- Die 26 decembris ---

% ··· Vesperas ···

% Magnificat (Kar121)
\defsheet[tone={8szo-c}]{DumMediumSilentium}
  {<-4--4---34t--4-----:-4--4--3----4t-eE2W1-12e-:-1----3r-5---4--5--3--5---7--6--7---tT4R3-3---,-3--4---5----4---3-4t-5---:-4---3---4--5--eE2-1-:-4--5--7---5---4---3----4t-5--4x-4x-?,}
  {  Mi-kor mély csend * bo-rí-tott be min---dent, s_az éj-sza-ka út-já-nak fe-lé-nél tar---tott, ak-kor jött el, ó U--runk, min-den-ha-tó i---géd ki-rá-lyi szé-ked-ből, al-le-lu-ja.}

% --- Die 29 decembris ---

% ··· Officium lectionis ···

% Antiphona 2 (Gondolkodó 4.)
\defsheet[tone={3-c}]{Orietur}
  {<-4u-7---:-7--7--4---3--5u-7-:-7-7--7---5---4---5--4--4-,-4---4-3---rR3-2-?,}
  {  Ki-kél * az Úr nap-ja-i--ban a bé-kes-ség-nek bő-sé-ge, s_Ő u-ral-ko--dik.}

% --- Die 30 decembris ---

% ··· Officium lectionis ···

% Antiphona 1 (Gondolkodó 4.)
\defsheet[tone={8szo-c}]{VeritasDeTerra}
  {<-4-7--uJ4-:-6---7--5---5----4-tT4--3--,-3--5--7-7---7-:-4-6--7--5----5-4-----4-?,}
  {  A hű-ség * fel-nö-vek-szik a föld-ből, és az i-gaz-ság a-lá-te-kint a menny-ből.}

% --- Die 31 decembris ---

% ··· Officium lectionis ···

% Antiphona 1 (Gondolkodó 4.)
\defsheet[tone={4-a}]{LaetenturCaeli}
  {<-1--2---4--4---4--5z-5---:-4----5---6---5---4-5---,-5--2t-4---Ï3-wW1-2wW1-:-0q---2--4---2--2-?,}
  {  Ví-gad-ja-nak az e--gek * s_ör-ven-dez-zen a föld, az Úr szí-ne e---lőtt   mert im-már el-jött.}

% Antiphona 3 (Gondolkodó 4.)
\defsheet[tone={6-a}]{Notum}
  {<-3---3--3---2--4--55-¨-3--3--1--0-,-3--3-3--2---4--55---5--tT4-rR3-3-?,}
  {  Meg-mu-tat-ta az Úr * al-le-lu-ja, az ő üd-vös-sé-gét, al-le--lu--ja.}

% ··· Laudes ···

% Benedictus
\defsheet[tone={7la-d}]{FactaEstCumAngelo}
  {<-4--4---4i--8---:-8--8--7-----6-8----8---8-9--8--7---9--9--iI7-7-:-©----------oO8--7---8o-iI7--7--,-7--tT4-4-:-4-4--5---tT4-3---4--4t--4--,-3--5-7---7---7--8o--8-:-£-----------5--7--tT4-3----4t-5--4--4-?,}
  {  Kö-rül-vet-te  * az an-gyalt a meny-nye-i se-re-gek so-ka-sá--ga, dicséretet zeng-vén és mond-ván: Di-cső-ség a ma-gas-ság-ban Is-ten-nek, és a föl-dön bé-kes-ség a_jóakaratú em-be-rek-nek, al-le-lu-ja.}

% --- Die 1 ianuarii -  In sollemnitate Sanctae Dei Genetricis Mariae ---

% ··· I. Vesperas ···

% Antiphona 1 (Kar128)
\defsheet[tone={6-a}]{OAdmirabileCommercium}
  {<X-3ed1Q0-:-¢----------3r--4--4z----rR3-3---:-¢---------4---3---4--4t-3y-3y-4zZ5T4-4---6-----5---4--3-:-4--3----3----4---3---3--3-3zZ5T4tT4-4--,-4----4t--3--3----eE2-q\\Q0\\-:-0--1----3-3---3--4t--4-:-tT4R3-4tT4-3y-3y-?,}
  {   Ó        csodálatra mél-tó szent cse-re: * az_emberi nem-nek Al-ko-tó-ja tes----tet s_lel-ket vé-ve  ke-gyes volt szü-let-ni a Szűz------től, s_ki-lép-ve mint em--ber,      ne-künk a-ján-dé-koz-ta  is----ten--sé-gét.}

% ··· Laudes ···

% Antiphona 3 (NZS33)
\defsheet[tone={2-g}]{GenuitPuerperaRegem}
  {<ô-4----3--4-----2------3-2--1--3---2y-----:-2--2--wW1--2---4-2eE2W1-1í-,-2----ñw-2-2----rR3-2--3r-5zZ5¨6uU6c-:-6--5----4--3--4----5-3--rR3-2y-,-2---4----3---4---3--2---eE2-1w-2wW1×ñ-:-1--2-rR3-1--3r-2y--2y-:-ËqQ0-1--2y-2y-?,}
  {   Nagy Ki-rályt szült, í-me az Asz-szony, * ne-ve fény-lik ö-rök----ké,  mint jó a-nya, ví--ga-do-zott,        de szűz ma-ra-dott i-ga-zá--ban. Nem volt ily cso-da még so--ha-sem,     az-u-tán se ha-son-ló,  al---le-lu-ja.}

% Benedictus
\defsheet[tone={1re-a}]{MirabileMysterium}
  {<-1---3--1--0---3----4--3t-5--:-5--4--2---3---4-eE2-1í-1í--,-1---3-1---0-3--4t---5---:-5----X6-5---4--4----4--3t-4--:-4-4--2-----4--5---4--3---:-2-----3---4--3-2--1---q\\Q0\\-,-0---1----3---3---3--4t--4---:-4-----3----4t--eE2-1---0q--1í-?,}
  {  Cso-dá-la-tos misz-té-ri-um * je-le-nik meg e nap né-künk: meg-ú-jul a te-remt-mény, mert em-ber-ré lett az Is-ten; a-mi volt, az meg-ma-radt, s_föl-vet-te a-mi nem volt,     nem szen-ved-vén ve-gyü-lést, s_nem szen-ved-vén meg-osz-lást.}

% ----------------------------------------------------------------------------
% III. Tempus Quadragaeismae
% ----------------------------------------------------------------------------

% ----------------------------------------------------------------------------
% a. Usque ad sabbatum hebdomadae quintae
% ----------------------------------------------------------------------------

% --- Pro cunctis diebus ---

% ··· Invitatorium ···

\defsheet{ChristumDominumProNobis}
  {<-0-----1ed1-2-23rR32¨2e1wW1-1--:-1--1--1eE2-3r-4t--4---4---3--4tT4R3-2eE2W1-1----,-0---012e-wW1--1eE2¨3r-4tT4R3r-eE2-.}
  {  Krisz-tus  U-run-----------kat, az ér-tünk kí-sér-tet-tet és gyö----tör----tet, * jöj-je---tek, i-------mád-----juk!}

% vel:
% TODO

% ··· Officium lectionis ···

% Hymnus (In officio dominicali)
\defsheet{ExMoreDoctiMystico}
  {<-3--4t----5----5--rR3-4----4t--5-:-4-2---3--4---eE2-1---1w-2---;-2----4---4t--tT4-wW1-3---wW1-1w-:-0--2----3r--2---qQ0-1w---wW1-1-?,-3-4t-5-5-rR3-4-4t-5-:-4-2-3-4-eE2-1-1w-2-;-2-4-4t-tT4-wW1-3-wW1-1w-:-0-2-3r-2-qQ0-1w-wW1-1-?,-3-4t-5-5-rR3-4-4t-5-:-4-2-3-4-eE2-1-1w-2-;-2-4-4t-tT4-wW1-3-wW1-1w-:-0-2-3r-2-qQ0-1w-wW1-1-?,-3-4t-5-5-rR3-4-4t-5-:-4-2-3-4-eE2-1-1w-2-;-2-4-4t-tT4-wW1-3-wW1-1w-:-0-2-3r-2-qQ0-1w-wW1-1-?,-3-4t-5-5-rR3-4-4t-5-:-4-2-3-4-eE2-1-1w-2-;-2-4-4t-tT4-wW1-3-wW1-1w-:-0-2-3r-2-qQ0-1w-wW1-1-?,-1wW1-0q-.}
  {  Meg-szen-telt ré-gi  szép szo-kás a böj-ti meg-tar-tóz-ta-tás;  tart-suk meg új--ra  gon-do--san, mi negy-ven nap-ra  szab-va  van. Mó-zes és a szent lát-no-kok pél-dá-ja már e-lénk ra-gyog, s_Krisz-tus, ki Úr min-den fö-lött, ma-gá-ra vett ily böjt-i-dőt. Kí-mé-lőn szól-jon hát sza-vunk, le-gyen sze-gé-nyebb asz-ta-lunk, ne győz-zön paj-zán tré-fa-ság, se ál-mos, tét-len lus-ta-ság. Ke-rül-jünk min-den vét-ke-set, szí-vün-ket rossz ne ront-sa meg, el-len-ség hogy-ha meg-kí-sért, hi-tünk a csap-dát zúz-za szét. Bol-dog Há-rom-ság, add meg ezt, osz-tat-lan Egy-ség, tedd meg ezt, ne-künk az üd-vös böjt-i-dő le-gyen sok jót gyü-möl-csö-ző. Á-men.}

% Hymnus (In officio feriali)
\defsheet{NuncTempusAcceptabile}
  {<-3--4t---5-5---rR3-4--4t-5-:-4--2----3---4---eE2-1----1w-2---;-2-4---4t---tT4--wW1-3---wW1-1w-:-0-2----3r-2---qQ0-1w--wW1-1-?,-3-4t-5-5-rR3-4-4t-5-:-4-2-3-4-eE2-1-1w-2-;-2-4-4t-tT4-wW1-3-wW1-1w-:-0-2-3r-2-qQ0-1w-wW1-1-?,-3-4t-5-5-rR3-4-4t-5-:-4-2-3-4-eE2-1-1w-2-;-2-4-4t-tT4-wW1-3-wW1-1w-:-0-2-3r-2-qQ0-1w-wW1-1-?,-3-4t-5-5-rR3-4-4t-5-:-4-2-3-4-eE2-1-1w-2-;-2-4-4t-tT4-wW1-3-wW1-1w-:-0-2-3r-2-qQ0-1w-wW1-1-?,-1wW1-0q-.}
  {  Ez most a ked-ve--ző i--dő, ke-gyel-mek nap-ja  ránk ra-gyog: a lan-kadt föld-re  hoz-za  már  a böjt a  bol-dog gyó-gyu-lást. Krisz-tus nagy fé-nye tün-dö-köl, üd-vünk-nek nyit-ja nap-ja-it, míg bűn-nel seb-zett szí-ve-ket fe-gyel-me gyó-gyít, jó-ra int. U-runk, tart-has-sunk üd-vö-sen test-ben-lé-lek-ben böj-tö-ket, s_u-tunk-on biz-tos fény le-gyen ö-rök Hús-vé-tod ö-rö-me. Ke-gyes Há-rom-ság, té-ge-det i-mád-ja-nak már min-de-nek, s_ki-ket me-gú-jít nagy ke-gyed, zeng-jük di-cső, új é-ne-ked. Á-men.}

% ··· Laudes ···

% In officio feriali

% Hymnus
\defsheet{IamChristeSolIustitiae}
  {<X-3--3---4---5---4---4-----3--2---:-4--2-----3----1---0--1----3--3--,-3----3---4--5----4---4--3-2---:-4--2---3--1---0---3----1--1-?,-02eE2W1Q0-1-.}
  {   Üd-vös-ség nap-ja, Krisz-tu-sunk, ha-sítsd szét vak-sá-gunk kö-dét, hogy fel-ra-gyog-jon az e-rény, mi-dőn re-ánk nap-palt ho-zol. Á---------men.}

% Responsorium breve
\defsheet{IpseLiberabitMe}
  {<X-3-3---3--4---3---4t-tT4-,-3-4--5--4----ed1--3r-4--3-?,-5--5-4---3--3--4t---5-?,-5--5---5-.}
  {   Ő sza-ba-dít meg en-gem * a va-dá-szok csap-dá-já-tól. És a dur-va be-széd-től. Di-cső-ség...}

% In officio dominicali
\defsheet{PrecemurOmnesCernui}
  {<X-3---3---4--5----4----4--3-2--:-4----2----3---1--0--1-3--3--,-3----3---4---5---4--4----3---2-:-4-2---3---1---0--3-----1--1-?,-02eE2W1Q0-1-.}
  {   Tér-del-ve mond-junk hő i-mát, zeng-jünk bűn-bá-nó é-ne-ket, tart-sák fel for-ró köny-nye-ink a vét-ket tor-ló szent ke-zet! Á---------men.}

% Responsorium breve
\defsheet{ChristeFiliDeiVivi}
  {<X-¢----------------4--3---4t-5--,-5--4---3----rR3-3--?,-[------------------4--5--4--:-5---rR3-4t--5-?,-5--5---5-.}
  {   Krisztus,_az_élő Is-ten Fi-a, * Ir-gal-mazz ne--künk. Akit_a_mi_gonoszsá-ga-in-kért tör-tek ösz-sze. Di-cső-ség...}

% ··· Tertia ···

% Hymnus
% TODO

% Antiphona
% TODO

% ··· Sexta ···

% Hymnus
\defsheet{QuaChristus}
  {<-0-1----2e-2-qQ0-1----2--1---:-1--2-----3---rR3-2--12e-eE2-2---,-2-4---4t---tT4--eE2--12e-eE2-2--:-0-1---2e--2-qQ0-1----2---1-.}
  {  U-runk ez ó-rán szom-ja-zott, ke-reszt-fán tes-te lan-ka--dott: a-kik most zsol-tárt mon-da--nak, a-zok-nak é-gi  szom-jat ad.}

% Antiphona
\defsheet[tone={tirr-b}]{VivoEgoDicitDominus}
  {<-uj5-4z--6---:-6----7--uU6-4t-,-7---7-7--6---5--4--5---5-rR3-2e-wW1-,-1--2---45z-5----5----rR3-2e--4---rR3-2-2---?,}
  {  É---lek én, * mond-ja az  Úr:  nem a-ka-rom ha-lá-lát a bű-nös-nek,  ha-nem in--kább hogy meg-tér-jen és  él-jen.}

% ··· Nona ···

% Hymnus
% TODO

% Antiphona
% TODO

% ··· Vesperas ···

% In officio dominicali

% Hymnus
\defsheet{AudiBenigneConditor}
  {<-qQ0-1---eE2-3-----2e---1--0-1---:-eE2-3r--rR3-1--3--1w--1--0--,-0-1----2----eE2-1--2---3-4--:-0q-3r---2---eE2---1---wW1-0--1-.}
  {  Sza-vun-kat halld meg, jó U-runk: a   hoz-zád es-dő szó-za-tot, a-mely-lyel sír-va kér i-mánk e  negy-ven szent nap al--ko-nyán.}

% In officio feriali

% Hymnus
\defsheet{IesuQuadragenariae}
  {<-3--2--4----tT4---3----4t--tT4R3-2-:-1--1t--5---5---tT4-4z-tT4-4t--;-5--5----4--tT4-3---4----eE2-1--:-3----2--4---tT4-3---4t-tT4R3-2-.}
  {  Úr Jé-zus, szent negy-ven-na-pos    böj-töd vál-lal-juk új-ra most: üd-vünk-re pél-dát adsz ne-künk, hogy fé-kez-zük ter-mé-sze-tünk.}

% Responsorium breve
\defsheet{EgoDixi}
  {<X-3---4----3----4t-5---,-4--5---rR3--4---3-?,-5---5-----5---4--4--,-4----4---4---5---rR3-4t-5-,-5--5---5-.}
  {   Így szól-tam: U-ram, * Kö-nyö-rülj raj-tam! Gyó-gyíts meg en-gem, mert vét-kez-tem el-le-ned. Di-cső-ség...}

% ············································································
% Hebdomada II Quadragesimae
% ············································································

% --- Feria VI ---

% ··· Vesperas ···

% Magnificat
\defsheet[tone={3-c}]{QuaerentesEumTenere}
  {<-4--5---6--7-5---7---6----:-4--4---5---5-zZ5-4--,-6----6---5--uU6-4-:-5--rR3-rR3-2-?,}
  {  El-fog-ni i-gye-kez-tek, * de fél-tek a nép-től, mert pró-fé-tá-nak tar-tot-ták őt.}

% --- Sabbato ---

% ··· Laudes ···

% Benedictus
\defsheet[tone={1la-a}]{VadamAdPatremMeum}
  {<X-1--1t-5----:-7-5----4-:-rf2-23r-eE2-1--1-,-3-123E2s0-:-34tT45z-5----4t-4---3---3---2e-4--eE2-1--1-?,}
  {   El-me-gyek * a-tyám-hoz és  mon-dom ne-ki: A-tyám,     fo------gadj en-gem nap-szá-mo-sa-id  kö-zé.}

% ············································································
% Hebdomada III Quadragesimae
% ············································································

% --- Dominica ---

% ··· I. Vesperas ···

% Antiphona 1
\defsheet[tone={8szo-c}]{ConvertiminiAdDominum}
  {<-4---4u--5---rR3-4t-4----:-8--7--5--7---7--,-©----------------8Oo8-uU67iI7-:-5---5--4----7--7-4----4--?,}
  {  Tér-je--tek az  Úr-hoz, * Is-te-ne-tek-hez: jóságos_és_irgal-mas  ő,        meg-bo-csát-ja a rosz-szat.}

% Antiphona 2
\defsheet[tone={8szo-c}]{Sacrificabo}
  {<-3--5---7-zh4-:-6--7--5--4---4--4--tT4-3-,-3t-7--zh4-6u-5---4--4--?,}
  {  Di-cső-í-tő  * ál-do-za-tot mu-ta-tok be, és az Úr  ne-vét hí-vom.}

% Antiphona 3
\defsheet[tone={3-c}]{NemoTollitAMe}
  {<-4---5z-¥------------------:-7--uU6-5--7--uU6-,-7--8---8--6---7-5---5-4-:-4--4--4t--eE2-1----2e---4---2-?,}
  {  Sen-ki sem_veszi_el_tőlem * az é---le-te-met,  ha-nem ma-gam a-dom o-da  és ma-gam ve--szem visz-sza azt.}

% ··· Laudes ···

% Antiphona 1
\defsheet[tone={1re-a}]{RegnavitDominus}
  {<-3--1---1---1t-5---:-5--4--5---rR3-34t-eE2-1---:-1--1eE2-3---4-5---rR3-2e-4---eE2-1--1-?,}
  {  Or-szá-gol az Úr, * ha-ta-lom-ba  öl--tö--zött, és most fel-ö-vez-te  ma-gát e---rő-vel.}

% Antiphona 2
\defsheet[tone={1la-a}]{VimVirtutisSuae}
  {<-0--1---1-1--1tu-5---:-[------------tT45uj5T4-,-5----4--3-3--3---3---3--3---3r-3-:-2e--4---eE2-1-?,}
  {  Sa-ját e-re-jé--ből * megkezdett_a tűz,        hogy fi-a-id meg-sza-ba-dul-ja-nak sér-tet-le--nül.}

% Antiphona 3
\defsheet[tone={2-g}]{RegesTerrae}
  {<-3----3--2--1---1-:-2--3---4---5--4--:-2--3----4--3---2--1--1-?,}
  {  Föld ki-rá-lya-i * és min-den né-pek, di-csér-jé-tek az Is-tent!}

% Benedictus (Anno A)
\defsheet[tone={8do-c}]{Aqua}
  {<-7-uU6--:-tT4-5z-tT4--5-4--,-4-4--2--4---eE2-3---:-1---2e---4t-5---4---4-?,}
  {  A víz, * a---me-lyet a-dok, a-ki ab-ból i---szik, nem szom-ja-zik töb-bé.}

% ············································································
% Sacrum triduum paschale
% ············································································

% --- Feria V in Cena Domini ---

% ··· Vesperas ···

% Responsorium breve
\defsheet{ChristusFactusEstBrevis}
  {<X-33r---3-:-¢---------------4tg3E2¨11rf2¨33r-3-:--3t¨34t77¨iI7i™uU6u¨55-44-:-4-5--7uj5g3¨zh4tT4t¨3rR3-3-.}
  {   Krisz-tus engedelmes_volt ér---------------tünk mind          @       @    a ha-lá------------------lig.}

% --- Feria VI in Passione Domini ---

% ··· Vesperas ···

% Responsorium breve
\defsheet{ChristusFactusEstLonga}
  {<X-33r---3-:-¢---------------4tg3E2¨11rf2¨33r-3-:--3t¨34t77¨iI7i™uU6u¨55-44-:-4-5--7uj5g3¨zh4tT4t¨3rR3-3-:--3-3--5z----445z¨uU6h44t¨zZ5T4-3--3--3--eD0¨3r¨tT4R3E24t3tT4t¨rR3-.}
  {   Krisz-tus engedelmes_volt ér---------------tünk mind          @       @    a ha-lá------------------lig, a ke-reszt-nek                ha-lá-lá-ig.}

% --- Dominica Paschae in Resurrectione Domini ---

% ··· Completorium ···

% Responsorium breve
\defsheet{HaecDies}
  {<-3--3--4-3--:-3--3----3--3--3----2----4--5---,-5--tT4-rR3-3---:-1----3---4----rR3-3-.}
  {  Ez az a nap, me-lyet az Úr szer-zett ne-künk, ör-ven-dez-zünk, s_ví-gad-junk raj-ta.}

% ----------------------------------------------------------------------------
% IV. Tempore Paschale
% ----------------------------------------------------------------------------

% ············································································
% b. Usque ad Ascensionem Domini
% ············································································

% --- Pro cunctis diebus ---

% ··· Hora Media ···

% Hymnus (Sexta)
\defsheet{VeniteServiSuplices}
  {<-5-5---5--5----5----5---4--4--,-5---5--5---4---3--4---5--5--,-5----5--5---4---3--4--3--2--,-3----4---4---3--4----5---4--4-?,-4tT4-3\\r\\-.}
  {  E-sen-gő szol-gák, jöj-je-tek, tár-já-tok fel ma szí-ve-tek, nyis-sá-tok dal-ra aj-ka-tok, hogy bol-dog Is-tent áld-ja-tok! Á----men.}

% --- In Ascencione Domini ---

% ··· Vesperas ···

% Hymnus
\defsheet{IesuNostraRedemptio}
  {<ô-2--2----1---2---3---4---2--2-:-2-3---4---3---1--2----4--334tT4R3y-,-2--3---4--5---3--4--3--2-:-2-3---4---3---2--3----2--1-----?,-1wW1-Ë0\\q\\-.}
  {   Jé-zus, meg-vál-tás kút-fe-je, a-kit sze-ret-ve szom-ja-zunk,       vi-lág te-rem-tő Is-te-ne, i-dők tel-jén ki föld-re szállsz, Á----men.}

% Antiphona
\defsheet[tone={8szo-c}]{AlleluiaTPSexta}
  {<-7--7--7--7---¨-5--4--5u-7-:-7--tT4-3r-4-?,}
  {  Al-le-lu-ja, * al-le-lu-ja, al-le--lu-ja!}

% ----------------------------------------------------------------------------
% V. Tempus per Annum
% ----------------------------------------------------------------------------

% ············································································
% In Dominicis
% ············································································

% --- Ad I. Vesperas (Cyclus unius anni) ---

% Magnificat
\defsheet[tone={1la-a}]{SustinuimusPacem}
  {<-0q--0---1-eE2-qQ0-:-2e-rR3-2eE2W1-0q-1--,--tT4-5---3--4---rR3-:-2--4t-wW1-3--wW1-0q--1--,-5--4--5---uU6--uj5-tT4R3-4--5--5--qQ0-,-2e-4--3--wW1-:-2-3---4t-qQ0-0---:-3--3--eE2-1wW1-0q-1-?,}
  {  Vár-tuk a bé-két, * és nem jött,  U--runk; a   jót ke-res-tük,  és í--me, za-var van itt: el-is-mer-jük, U---runk, bű-ne-in-ket;  ne ha-ra-gudj  ö-rök-ké mi--ránk, Iz-ra-el  Is---te-ne.}

% --- Anno A ---

% Dominica XIV

% Benedictus
\defsheet[tone={1re-a}]{ConfiteorTibiPater}
  {<X-0--1---1t-5u-5----:-tT45uj5-4---:-3--4--5----rR3r-1-,-3----2--Ÿ------------------------------2---3---4t-4--3--,-¢-----------------2--3---4---eE2-1---1-?,}
  {   Di-cső-í--te-lek, * A-------tyám, ég és föld U----ra, mert az okosak_és_a_bölcsek_elől_elrej-tet-ted e--ze-ket, és_a_kicsinyeknek ki-nyi-lat-koz-tat-tad.}

% Dominica XV

% Benedictus
\defsheet[tone={8szo-c}]{CumTurbaPlurima}
  {<-4-7--5---rR3--4t-4---:-7---7--8--uj5-7--7---,-©-------------------7--8--8oO8-uU67iI7-,-5---5--5--7---7---4---4---:-4--4----4-4---eE2-1-:-3----3--4---5--4---4-?,}
  {  A-mi-kor nagy tö-meg * vet-te kö-rül Jé-zust, és_a_városokból_hoz-zá si-et---tek,      pél-da-be-szé-det mon-dott: Ki-ment a mag-ve--tő, hogy el-ves-se mag-ját.}

% Dominica XXI

% Benedictus
\defsheet[tone={8szo-c}]{QuemDicuntHomines}
  {<-rF1-eE2-3r---4---:-4--4t-tT4-3---5--5zZ5-4---3r-4-;--4uU6-7i-iI7-5---7-uU6h4-5---rR3-5---zZ5-4--,-4--4t--tT4tT4R3-5uu-5zZ5-4---4-;--rR3d13r-r3ttT4tg3r35u-7i----8--:-8--uU6-4tT4R3-5uj5-zZ5-4--4-,--4--4t-5---4t--tT4R3-3--:-3t---7i78o-uU6--5--5--;-5i-8--uU6-5-4--4t---rR3-3-:-3--3---4t-5èu-zZ5-5-57iiI7-tT4R3-5zZ5-4-?,}
  { „Ki--nek tart-ják * az em-be--rek az Em---ber-fi-át?” mond-ta Jé--zus a ta----nít-vá--nya-i---nak. Fe-lel-vén      Pé--ter, mon-dá: „Te      vagy          Krisz-tus, az é---lő     Is---ten Fi-a.” „Én is mon-dom ne----ked, hogy te    vagy Pé-ter, és er-re  a kő-szik-lá--ra  fo-gom é--pí--ten-i egy----há----za---mat.”}

% Magnificat
\defsheet[tone={8szo-c}]{Quodcumque}
  {<-3r-4--¨-5---4-----3r-2e--1--;-1---0----34t-5--5uj5-5z----4---4-,-7--uj5-7i--8---8o-uU6--4-34t-5--;-7u8µ4-4----5--rf2-3--4-2-----3---1-;--5----tT4-5--5--5--5uj5-zZ5-4---4-?,}
  { „A-mit * meg-kötsz a  föl-dön, meg lesz köt-ve a    menny-ben is; és a---mit fel-ol-dasz a föl-dön, fel   lesz ol-doz-va a menny-ben is,” mond-ta  az Úr Si-mon  Pé--ter-nek.}

% --- Anno C ---

% Dominica XV

% Benedictus
\defsheet[tone={8szo-c}]{HomoQuidam}
  {<-4---4u-5---:-[------------------rR3-4t-4---:-7---7---8--7--5--7--7---:-¦------------------8--7---8oO8-uU67i-7--:-5--4---5---7---7---4----4-?,}
  {  Egy em-ber * Jeruzsálemből_Jeri-kó--ba ment. Rab-lók ke-zé-be ke-rült. Ezek_kifosztották, vé-res-re   ver---ték, és fél-hol-tan ott-hagy-ták.}

% Dominica XIX

% Benedictus
\defsheet[tone={8szo-c}]{NoliteTimere}
  {<-4--4u----5--:-3--4t----5--4---,-8----8---8o--8-----7-7----uU67iI7-:-5----5---5---5--7--7--4---4-?,}
  {  Ne félj, te * ki-csiny-ke nyáj, mert úgy tet-szett A-tyád-nak       hogy nek-tek ad-ja or-szá-gát.}

% ············································································
% In Solemnitatibus Domini
% ············································································

% --- Sanctissimi Corporis et Sanguinis Christi ---

% I. Vesperas

% Antiphona 1
\defsheet[tone={2-g}]{MiseratorDominus}
  {<ô-4--4--4---2---rR3-2---:-4-4--4----2-4----5z-5x--,-5--5--5--6--5---rR3--2--:-2---2--1---ñ--1w-2-?,}
  {   Az ir-gal-mas Is--ten * e-le-delt a-dott né-künk, em-lé-ke-ze-tet szer-zett cso-da-tet-te-i--nek.}

% Antiphona 2
\defsheet[tone={5-g}]{QuiPacem}
  {<-2--0--2---4t-4----:-4--4---5---7--6--tT4-4-,-4-5--4--3--4t-rR3E2-2-:-2-1w-eE2-1--0í--0í-?,}
  {  Ki bé-két te-remt * az Egy-ház ha-tá-ra--in, a ga-bo-na ja-vá----val e-lé-git ki min-ket.}

% Magnificat
\defsheet[tone={1la-a}]{OQuam}
  {<-3--0--1----1tu-5---:-[----------------tT45uj5T4-:-4----5---4---3--3r-3-:-3-2--3-4---3---2--1e--3-:-3--4---5---3----2--1--qQ0-:-1--eE2-1í-1í-?,}
  {  Ó, mi-lyen é---des * hozzánk_a_te_lel-ked:        hogy ked-ves-sé-ge-det a fi-a-kon meg-mu-tas-sad ke-nye-ret adsz az ég-ből,  al-le--lu-ja.}

% Laudes

% Antiphona 1
\defsheet[tone={7do-a}]{AngelorumEsca}
  {<ô-5--5---3---5-6--7--zZ5-5---:-5--5--3---4-5--wW1-1--,-£------------2r----3-:-2--1---2---4----3----1í-1í-?,}
  {   An-gya-lok e-le-de-lé--vel * él-te-ted a Te né--ped, s_kenyeret_a menny-ből bő-ség-gel nyúj-tasz né-kik.}

% Hora Medias

% Antiphona (tertia)
\defsheet[tone={4-a}]{Desiderio}
  {<-1----2--4t---5----:-4----5-6---5--4--5--5--5x-,-2--5----4---Ï3--2--1--wW1-:-0--1-2r---2----2--2y-.}
  {  Vágy-va vágy-tam, * hogy e hús-vé-ti la-ko-mát  el-költ-hes-sem ve-le-tek,  mi-e-lőtt szen-ve-dek.}

% --- Sanctissimi Cordis Iesu ---

% Laudes

% Antiphona 2
\defsheet[tone={1re-a}]{VeniteAdMe}
  {<-3---0--1---2tu-5----:-5--5----5---tT45uj5T4-:-4-4---5--4--3r--3--:-2-----3r--eE2-1í---1í-,-1----3---wW1-2--3----4t-4--3--:-2----3--4--eE2-1í----1í-.}
  {  Jöj-je-tek ho--zám, * ti mind-nyá-jan,        a-kik fá-ra-doz-tok, s_ter-hel-ve  vagy-tok, s_én meg-eny-hít-lek ti-te-ket, mond-ja az Úr  Krisz-tus.}

% Hora Media

% Antiphona (tertia)
\defsheet[tone={7do-a}]{PopuleMeus}
  {<ô-5tg3-56uj5-5----:-5---5--3---4--5--wW1-1í-,-£-------------------2r-3y-:-2--1----2rR3-1í-1í-?,}
  {   Én   né----pem, * mit te-tem én el-le--ned? Vagy_miben_ártottam né-ked? Fe-lelj meg  né-kem!}
\makeatletter

\makeatletter

% ············································································
% Hymnus
% ············································································

\newcommand{\sheet@TeLucisAnteTerminum}{%
  \begin{gregosheet}
    \melody{<-7--7---7-7---7--7--6--5---:-7--7---7-----7--4----5--4--3--,-5----5--6----7--6--5----4--5zZ5-:-5-5---4----3--4----5--4--4-.}
    \lyrics{  Im-már a nap le-ál-do-zott, Te-rem-tőnk, ké-rünk té-ge-det, légy ke-gyes és ma-radj ve-lünk,  ő-riz-zed, óv-jad  né-pe-det!}
  \end{gregosheet}
}

% vel:
\newcommand{\sheet@ChristeQuiSplendor}{%
  \begin{gregosheet}
    \melody{<-7-----7----7---7---7--7---6--5---:-7--7--7--7----4-5----4--3--,-5---5----6---7----6---5----4-5zZ5-:-5--5----4----3--4---5-----4--4-.}
    \lyrics{  Krisz-tus, tün-dök-lő nap-pa-lunk, az éj ár-nyát e-losz-la-tod, kit fény-ből fény-nek vall a hit    és fény-nyel ál-dod szent-je-id.}
  \end{gregosheet}
}

% ············································································
% Responzorium breve
% ············································································

% Responsorium breve in Tempore Paschali
% TODO

% Responzorium breve
\newcommand{\sheet@InManusTuasDomine}{%
  \begin{gregosheet}
    \melody{<X-3r-3---,-4--4---4--3-4t--4---ed1-3r-3-,,-5---5---5z--5---5---4---4t-4-:--5--4--4---rR3-4-4t-,,-5--5---5---5--5-5z--5-:--5-5--4-4-:--5--5-4-----3--4t--tT4-.}
    \lyrics{   U-ram, * ke-zed-be a-ján-lom lel-ke-met. Meg-vál-tot-tál min-ket U-runk, hű-sé-ges Is-te-nünk. Di-cső-ség az A-tyá-nak, a Fi-ú-nak, és a Szent-lé-lek-nek.}
  \end{gregosheet}
}

% Nunc dimittis
\newcommand{\sheet@SalvaNos}{%
  \begin{gregosheet}
    \melody{<X-1-----3----1-0----3---4----3--4t---5x---:-6-5----4---4--4--3t---4y-,-4----4-4-----2---4---5---4----3y---:-3--2--3---4---eE2-1---0q--1í-,,-3----4--[--------------------4-3----,-[-----------------4--3---4t--4-,,-5----5---5---4---3--3--,-5-5--5--4---3--4t--4-,,-5-5---4---3---3---,-5---5---5--5---4---3--4t-4--,,-3r---[------------------------5---5--4--3---,-[------------------------4--3--4t-4-,,-3--4---[-----------------4--3-3---,-5--5-----4--3---4--4t--4-,,-3--4---[---------------------------4---3---3---,-[----------------4---3---4t-4rR3E2W1í-,,-1-----3-.}
    \lyrics{   Tarts meg, U-ram, vir-rasz-tá-sunk-ban, * ő-rizz meg al-vá-sunk-ban, hogy a Krisz-tus-sal vir-rasz-szunk, és bé-kes-ség-ben nyu-god-junk! Most bo-csátod_el_szolgádat, U-ram, * a_te_igéd_szerint bé-kes-ség-ben. Mert lát-ták sze-me-im * a te üd-vös-sé-ged-et.  A-kit ren-del-tél * min-den né-pek szí-ne e--lőtt, Hogy fény_legyen_a_népek_vilá-gos-sá-gá-ra, * és_dicsőségére_népednek, Iz-rá-el-nek. Di-cső-ség_az_Atyának_és Fi-ú-nak * és Szent-lé-lek Is-ten-nek, Mi-kép-pen_kezdetben_vala,_most_és min-den-kor * és_mindörökkön_ö-rök-ké! Á--men.         Tarts meg...}
  \end{gregosheet}
}

% ············································································
% Dominica Post I Vesp.
% ············································································
% Antiphona 1
\newcommand{\sheet@MiserereMihiDomine}{%
  \begin{gregosheet}[tone=8. szó tónus]
    \melody{<-4--4---4----5---4----tT4-3y---,-5--7----6u---5x-:-5-4---3--4t-4y-,,-7-7-6-7-5-4-.}
    \lyrics{  Kö-nyö-rülj raj-tam, U---ram, * és hall-gasd meg  i-mád-sá-go-mat!}
  \end{gregosheet}
}

% Antiphona 2
\newcommand{\sheet@InNoctibus}{%
  \begin{gregosheet}[tone=4. tónus]
    \melody{<X-0--1--3r--4-:-5---6--5---4--4t-4-,,-7-7-6-7-6-4-.}
    \lyrics{   Éj-je-len-te  áld-já-tok az U--rat!}
  \end{gregosheet}
}

% ············································································
% Dominica Post II Vesp.
% ············································································
% Antiphona
\newcommand{\sheet@QuiHabitat}{%
  \begin{gregosheet}[tone=8. szó tónus]
    \melody{<-4-4--4---4---4t-tT4-3---:-3--5---5u-7---4--3r-4--,-5-7---7---6u-5-:-5--4---3--4t--4--4-,,-7-7-6-7-6-4-.}
    \lyrics{  A ma-gas-ság-be-li--nek * ol-tal-má-ban la-ko-zik, a min-den-ha-tó  ár-nyé-ká-ban ma-rad.}
  \end{gregosheet}
}

% ············································································
% Feria II
% ············································································
% Antiphona
\newcommand{\sheet@IntendeVoci}{%
  \begin{gregosheet}[tone=8. szó tónus]
    \melody{<-4--6----uU6-5z--:-5---5--uj5-zZ5-4-,,-7-7-6-7-5-4-.}
    \lyrics{  Kö-nyör-gé--sem * sza-vá-ra  fi--gyelj.}
  \end{gregosheet}
}

% ············································································
% Feria III
% ············································································
% Antiphona
\newcommand{\sheet@NeIntres}{%
  \begin{gregosheet}[tone=8. szó tónus]
    \melody{<-4t-5------:-5u--7i-7----6---5----4-4-,,-7-7-6-7-5-4-.}
    \lyrics{  Ne szállj * per-be szol-gád-dal, U-ram.}
  \end{gregosheet}
}

% ············································································
% Feria IV
% ············································································
% Antiphona 1
\newcommand{\sheet@EstoMihi}{%
  \begin{gregosheet}[tone=6. tónus]
    \melody{<-3----4--5----rR3-3----:-3--1---3--4--4--3--3-,,-6-6-5-4-5-4-.}
    \lyrics{  Légy né-kem, U---ram, * ol-tal-ma-zó Is-te-nem!}
  \end{gregosheet}
}

% Antiphona 2
\newcommand{\sheet@DeProfundis}{%
  \begin{gregosheet}[tone=8. szó tónus]
    \melody{<-7-7----7--7---5---:-5--4--5---6---5----4--3r-4-,,-7-7-6-7-5-4-.}
    \lyrics{  A mély-sé-gek-ből * ki-ál-tok hoz-zád, ó, U--ram}
  \end{gregosheet}
}

% ············································································
% Feria V
% ············································································
% Antiphona
\newcommand{\sheet@CaroMea}{%
  \begin{gregosheet}[tone=4. tónus]
    \melody{<-tT4-3---:-4--3----wW1--3--rR3--2---2-,,-5-5-4-5-4-2-.}
    \lyrics{  Tes-tem * el-nyug-szik re-mény-ség-ben.}
  \end{gregosheet}
}

% ············································································
% Feria VI
% ············································································
% Antiphona
\newcommand{\sheet@DomineDeusSalutisMeae}{%
  \begin{gregosheet}[tone=4. tónus]
    \melody{<-3-1----3---3--4--3--4--4t-4--,-3--2---2e--4-:-3--2---3---4--3--2-,,-5-5-4-5-4-2-.}
    \lyrics{  U-ram, sza-ba-dí-tó Is-te-nem, éj-jel-nap-pal te-hoz-zád ki-ál-tok!}
  \end{gregosheet}
}

% ············································································
% Antiphonae finales ad Beatam Mariam Virginem
% ············································································
% Tempus Quadragaesimae
% TODO

% Simpl.
\newcommand{\sheet@AveReginaCaelorum}{%
  \begin{gregosheet}
    \melody{<X-3-2--1--0--1--3---4--3----:-5-7---6--4--5--4--3--5--4--;-3---2--1--0----1---3---4---3-:-5--4---6---5--4---1---4--3-,-3---4---5---5--4---5--6-5-:-7--6---5--4---3---1--4-3-,-6--5---4-6---5--4--34t-5-;-7-6---67i-5--5-----4---3--44-3-.}
    \lyrics{   A-ve Re-gi-na cae-lo-rum, * a-ve, Do-mi-na an-ge-lo-rum, sal-ve ra-dix, sal-ve, por-ta, ex-qua mun-do lux est or-ta. Gau-de, Vir-go glo-ri-o-sa, su-per om-nes spe-ci-o-sa; va-le, o val-de de-co-ra, et pro no-bis Chris-tum ex-o-ra.}
  \end{gregosheet}
}

% Tempus Pachale
% TODO

% Simpl.
\newcommand{\sheet@ReginaCaeli}{%
  \begin{gregosheet}
    \melody{<X-3--4--3--4---5--:-6---5--4-:-6--5--4--3-,-3---7-7----8--7--6--5--4---5--6-:-6--5--4--3-,-7--7---8---7-:-7--3---4--3--:-4--5--6--7-,-7-3--4---6--5---4--3-:-2--4--4--3-.}
    \lyrics{   Re-gi-na cae-li * lae-ta-re, al-le-lu-ja: Qui-a quem me-ru-is-ti por-ta-re, al-le-lu-ja: Re-sur-rex-it, si-cut di-xit, al-le-lu-ja: O-ra pro no-bis De-um, al-le-lu-ja.}
  \end{gregosheet}
}

% Tempus Per Annum
% TODO

% Simpl.
\newcommand{\sheet@SalveRegina}{%
  \begin{gregosheet}
    \melody{<-0---2---4--5--4---,-5--7---6--5--4--5---4--4-,-7--4---5---33-1-:-2--3----4---2----wW1-0-,,-4--5--6---7--4-:-5--6--7---6--5--4-5-4-,,-7--4--5---3--1--2--:-2--4---5---7--5----4-:-5--4---3---2--1--2---1---0-,,-4-5---6--7-:-4--5--7--6--5---4--,-7--4---5--3--1--2--3--4---3---5-4--4-:-3--2---1---wW1-0-,,-4--5z-7--:-6--4--5---4---4----5---7----6---5--4-,-0--4---5----7---6--5--4--2--3--wW1-0-,,-2å4-22--0--,,-4ç6uU6-tT4-4-,,-77¨4tg3¨ed1-2e--4-:-0---3--2--qQ0-0-.}
    \lyrics{  Sal-ve, Re-gi-na, * ma-ter mi-se-ri-cor-di-ae, vi-ta, dul-ce-do, et spes nos-tra, sal-ve.  Ad te cla-ma-mus ex-su-les fi-li-i E-vae. Ad te sus-pi-ra-mus, ge-men-tes et flen-tes in hac lac-ri-ma-rum val-le.  E-ia, er-go, ad-vo-ca-ta nos-tra, il-los tu-os mi-se-ri-cor-des o-cu-los ad nos con-ver-te.  Et Ie-sum, be-ne-dic-tum fruc-tum vent-ris tu-i,  no-bis post hoc ex-si-li-um os-ten-de.  O   cle-mens, o      pi--a,   o           dul-cis Vir-go Ma-ri--a.}
  \end{gregosheet}
}

\makeatother

% ============================================================================
% DOMINICA
% ============================================================================

% ············································································
% I. Vesperas
% ············································································

% Hymnus
\defsheet{DeusCreatorOmnium}
  {<ô-2--2----2--1--2---4--3--2í-:-2--3--4---5---4---5---5--6x-,-6--6---5---6-4---5----rR3-2í-:-2--1---2--4----5---2--1---2í-?,-1e4R3E2W1-2í-.}
  {   Is-ten, vi-lá-got al-ko-tó,  vi-lá-got sar-kán for-ga-tó,  ki nap-nak é-kes fény-ru--hát, éj-nek ke-gyel-mes ál-mot adsz. Á---------men.}

% Antiphona 1
\defsheet[tone={8szo-c}]{DomineClamavi}
  {<-¦-----------5--4----45u-7c---:-5--4---5---4x-?,}
  {  Tehozzád_ki-ál-tok, U---ram, * si-ess hoz-zám!}

% Antiphona 2
\defsheet[tone={8szo-c}]{PortioMeaDomine}
  {<-¦-------------5--4----45u--7c--:-7--7-5---4---5----4x-?,}
  {  Osztályrészem te légy U----ram * az é-lők-nek föld-jén!}

% Antiphona 3
\defsheet[tone={1la-a}]{HumuliavitSemetipsum}
  {<-1---3-1---0--3--4t-5---:-XzZ5-4--3t-4--,-2-4----tT4-3-:-¢------------2--3--4----eE2-1---1-?,}
  {  Meg-a-láz-ta ön-ma-gát * az   Úr Jé-zus, e-zért Is--ten felmagasztal-ta őt mind-ö---rök-re.}

% Responsorium breve
\defsheet{QuamMagnificataSuntRBr}
  {<X-3----4----3---4t-5--,-5-5--5--4--3---rR3-3--?,-5---5----5z--4-:-4--3---4t-?,-[-------------5z--5-:-5-5--4-4-:-5--5-4-----3--4t--tT4-.}
  {   Mily nagy-sze-rű-ek * A te mű-ve-id, U---ram.  Böl-cses-ség-gel al-kot-tad.  Dicsőség_az_A-tyá-nak a Fi-ú-nak és a Szent-lé-lek-nek.}

% ············································································
% Invitatorium
% ············································································

\defsheet{VeniteExsultemusDomino}
  {<X-3---4t-5----:-5--5---5---3----4t-4\\-3\\-.}
  {   Jöj-je-tek, * ör-ven-dez-zünk az Úr--nak.}

% ············································································
% Officium lectionis
% ············································································

% Hymnus
\defsheet{PrimoDierumOmnium}
  {<ô-2--2--2---1---2---4---3--2í--:-2--3----4--5------4--5---5--6x--,-6----6----5---6--4----5--rR3-2í-:-2--1----2--4--5----2-1--2í-?,-1e4R3E2W1-2í-.}
  {   El-ső nap min-den nap kö-zött, me-lyen ég s_föld al-kot-ta-tott, s_me-lyen fel-tá-madt Al-ko--tónk le-győz-te ér-tünk a ha-lált, Á---------men.}

% Antiphona 1
\defsheet[tone={1re-a}]{InLege}
  {<-3--1äeE2s0-¨-3---5--4---tT45zZ5-:-5---4---3---3-4--eE2-1-:-3r-2---3--1---1-?,}
  {  Az Úr      * tör-vé-nyé-ben       van tel-jes a-ka-rat-tal éj-jel és nap-pal.}

% Antiphona 2
\defsheet[tone={1la-a}]{Praedicans}
  {<-1t--5z--5--¨-5--6--4---5--5---4---2-:-4--2-3-----4--eE2-1--0w--rf2-eE2-1-?,}
  {  Hir-det-ni * az Úr-nak pa-ran-csa-it  az ő szent he-gyé-re ren-del-te--tett.}

% Antiphona 3
\defsheet[tone={8szo-c}]{TuEsGloria}
  {<-3r-4----¨-5--4---eE2-3--:-1--2----3--4---5--rR3--4t-5--:-7--7-6--7----5---7--4--3-,-3--4---5----7---7i--7--6-:-7-5--rR3---4t-5----4-?,}
  {  Te vagy * di-cső-sé--gem, te vagy ol-tal-ma-zóm, U--ram, te e-mel-ted fel fe-je-met és meg-hall-gat-tál en-gem a te szent he-gyed-ről.}

% ············································································
% Laudes
% ············································································

% Hymnus
\defsheet{AeterneRerumConditor}
  {<-5--5---5--5---5--5--4--4-,-5--5--5----4-3---4--5--5-,--5-5---5-4--3---4---3--2----,-3-------4--4--3--4---5-4--4-?,-4tT4-3\\r\\-.}
  {  Vi-lág te-rem-tő Is-te-ne, ki Úr vagy é-jen és na-pon, i-dőt i-dő-vel vál-to-gatsz, s_nincs mű-ve-id-ben u-na-lom, Á----men.}

% Antiphona 1
\defsheet[tone={7la-d}]{AdTeDeLuce}
  {<-4---6---7--8--8---¨-9---8--7---7----7--7uU6-,-7----8oO8-uU67iI7-5--6uU6-4--4-?,}
  {  Hoz-zád éb-re-dek * vir-ra-dat-kor, Is-ten,   hogy lás--sam     ha-tal--ma-dat.}

% Antiphona 2
\defsheet[tone={8szo-c}]{HymnumDicamus}
  {<-4----4----tT4-3-----¨-5--5èuU6-5tT4-5z-5--4----rR3d1-3--3r-4--4-?,}
  {  Mond-junk him-nuszt * az Úr----nak, Is-te-nünk-nek,  al-le-lu-ja.}

% Antiphona 3
\defsheet[tone={8szo-c}]{FiliiSion}
  {<-7--7--7--uU6-5x-:-5--5----4--5---6--5--4----4----4--5z-4--4-?,}
  {  Si-on fi-a---i  * uj-jong-ja-nak Ki-rá-lyuk-ban, al-le-lu-ja.}

% Responsorium breve
\defsheet{ChristiFiliDeiVivi}
  {<X-¢----------------4--3---4t-tT4-,-5--4---3----rR3-3--?,-5z-4-:-4--4-5---rR3--4t--5-?,--[-------------5z--5-:-5-5--4-4-:-5--5-4-----3--4t--tT4-.}
  {   Krisztus,_az_élő Is-ten Fi-a,  * Ir-gal-mazz ne--künk. Ak-i   az A-tya jobb-ján ülsz. Dicsőség_az_A-tyá-nak a Fi-ú-nak és a Szent-lé-lek-nek.}

% ············································································
% Hora Media
% ············································································

% Hymnus
\defsheet{RectorPotens}
  {<-5-5-----5---5--5---5--4--4\\-,-5-5---5----5---4--5---5---5x-,-5---5---5--4--3---4--3---2\\-,-3--4---4---4--3--5--4---4x-?,-4tT4-3\\r\\-.}
  {  U-runk, fel-sé-ges Is-te-nünk, a lét rend-jét ki meg-sza-bod: reg-gel vi-lá-gos-sá-got adsz, és dél-ben tű-ző ra-gyo-gást, Á----men.}

% Antiphona 1
\defsheet[tone={8szo-b}]{BonumEst}
  {<XB-1----2e-3r-3---:-3--4---3---23r-wW1-0í-,-0----2--3--4---3---34zZ56uU6-:-4--3--2e--3y-?,}
  {    Jobb az Úr-ban * bi-zal-mat tar-ta--ni,  mint fe-je-del-mek-ben         bi-za-kod-ni.}

% Antiphona 2
\defsheet[tone={4-a}]{FortitudoMea}
  {<-0--1e-3-2---1w-2---:-2--4--5---tT4R3-rR3-1w-2-,-2e-4-eE2--wW1-0---1---1--3rR3-2--2-?,}
  {  Az én e-rős-sé-gem * és di-csé-re----tem az Úr  és Ő lett né--kem sza-ba-dí---tá-som.}

% Antiphona 3
\defsheet[tone={8szo-c}]{Confitebor}
  {<-3--3---3r-4---¨-5--7----7i-7--:-7----6---5----6u--5---4--4-?,}
  {  Há-lát a--dok * ne-ked, U--ram, mert meg-hall-gat-tál en-gem.}

% ············································································
% II. Vesperas
% ············································································

% Hymnus
\defsheet{LucisCreator}
  {<ô-2----2--2--1--2---4--3-2í--:-2--3---4---5---4--5-----5--6x-,-6----6--5----6--4--5---rR3-2í-:-2--1---2-4---5---2---1--2í-?,-1e4R3E2W1-2í-.}
  {   Fény al-ko-tó-ja, jó U-runk, ki min-den nap-ra fényt ho-zol, s_az új fény el-ső haj-na--lán  ad-tál a lét-nek kez-de-tet, Á---------men.}

% Antiphona 1
\defsheet[tone={1la-a}]{Virgam}
  {<-3-0--1-1tu-5---:-5---5--tT45uj5T4-,-5--4--3-----3--2--3r-eE2-1-1---,-1-2e--4tT4-3-:-¢----------2--3--4--3--2--1---1-?,}
  {  A te e-rőd-nek * pál-cá-ját         ki-bo-csájt-ja az Úr Si--on-ból: u-ral-kod--jál a_te_ellen-sé-ge-id kö-ze-pet-te.}

% Antiphona 2
\defsheet[tone={4-a}]{ExAegypto}
  {<-6-tT4--5z--5---¨-6--7--6---5---4---5z---6-5-?,}
  {  E-gyip-tom-ból * ki-ve-zet-tél min-ket, U-runk.}

% Antiphona 3
\defsheet{SalusEtGloria}
  {<----7--7--8oO8-7i-:-[---------------------------4--5----rR3-:-----4--4t-3--3-,------[--------------------------4--3--4t-5-:-----7--7--iI7-5-:-5--4tT4-3--3-?,-----7--7--8oO8-7i-:-[----------------------------------4----5--rR3-:---4--4t-3--3-,------[-----------------------------------4----3--4t-5-:--------7--7--iI7-5-:-5--4tT4-3--3-?,-----7--7--8oO8-7i-:-[----------------------5---5---5--5--4--5--rR3-:------4--4t-3--3-,------[------------------------------4---3--4t--5-:------7--7--iI7-5-:-5--4tT4-3--3-?,-----7--7--8oO8-7i-:-[----------------------------4--5---4---rR3-:----4--4t-3--3-,------[------------------4--3-----4t--5-:-------7--7--iI7-5-:-5--4tT4-3--3-?,-----7--7--8oO8-7i-:-[----------------------4--5-rR3-:-----4--4t-3--3-,------5--5-----5--4---3--4t--5-:-------7--7--iI7-5-:-5--4tT4-3--3-?,-----7--7--8oO8-7i-:-[--------------------------------4---5---rR3-:-----4--4t-3--3-,------[----------------4---3---4t-5-:-------7--7--iI7-5-:-5--4tT4-3--3-.}
  {<V.> Al-le-lu---ja.  Üdv,_dicsőség_és_hatalom_Is-te-nünk-nek, <F.> al-le-lu-ja, <V.> mert_igazak_és_jóságosak_í-té-le-te-i. <F.> Al-le-lu--ja, al-le---lu-ja. <V.> Al-le-lu---ja.  Dicsérjétek_Istenünket,_ti,_összes szol-gá-i, <F.> al-le-lu-ja, <V.> és_akik_félve_tisztelitek_őt,_kicsi-nyek és na-gyok. <F.> Al-le-lu--ja, al-le---lu-ja. <V.> Al-le-lu---ja.  Átvette_uralmát_Urunk, min-den-ha-tó Is-te-nünk, <F.> al-le-lu-ja, <V.> örvendezzünk,_ujjongjunk_és_di-cső-ít-sük őt! <F.> Al-le-lu--ja, al-le---lu-ja. <V.> Al-le-lu---ja.  Elérkezett_a_Bárány_menyegző-jé-nek nap-ja, <F.> al-le-lu-ja, <V.> már_fel_is_készült me-nyasz-szo-nya. <F.> Al-le-lu--ja, al-le---lu-ja. <V.> Al-le-lu---ja.  Dicsőség_az_Atyának_és Fi-ú-nak, <F.> al-le-lu-ja, <V.> és Szent-lé-lek Is-ten-nek. <F.> Al-le-lu--ja, al-le---lu-ja. <V.> Al-le-lu---ja.  Miképpen_kezdetben_vala,_most_és min-den-kor, <F.> al-le-lu-ja, <V.> és_mindörökkön_ö-rök-ké, á--men. <F.> Al-le-lu--ja, al-le---lu-ja.}

% Responsorium breve
\defsheet{BenedictusEs}
  {<X-3--4----3-----4t-5-----,-5--4--3----rR3-3--?,-5--5---4---5--4---4-:-5--5--5----4---3-4t--5-?,-[-------------5z--5-:-5-5--4-4-:-5--5-4-----3--4t--tT4-.}
  {   Ál-dott vagy, U--runk, * Az ég-bolt fö--lött. Di-csé-ret-re mél-tó  és ma-gasz-tos ö-rök-ké.  Dicsőség_az_A-tyá-nak a Fi-ú-nak és a Szent-lé-lek-nek.}

% ============================================================================
% FERIA II
% ============================================================================

% ············································································
% Invitatorium
% ············································································

\defsheet{IubilemusDeo}
  {<X-¢---------4--3--4t---5---,-5--5---3-----4t-4--3-.}
  {   Szabadító Is-te-nünk-nek * vi-gad-junk, az Úr-nak!}

% ············································································
% Officium lectionis
% ············································································

% Hymnus
\defsheet{SomnoRefectisArtubus}
  {<-3--3---4---5----4----4---3-2--:-4---2----3--1-0------1---3--3---,-3-3------4---5---4----4-3--2---:-4-2----3--1---0----3----1--1--?,-0w32W1Q0-1í-.}
  {  Al-vás-ban fris-sült már a test ott-hagy-va á-gyunk, fel-ke-lünk. A-tyánk, hoz-zád szól é-ne-künk; e-seng-ve kér-jük: légy ve-lünk. Á--------men.}

% ············································································
% Laudes
% ············································································

% Hymnus
\defsheet{SplendorPaternaeGloriae}
  {<-3--3---4---5----4----4---3-2--:-4---2----3--1-0------1---3--3---,-3-3------4---5---4----4-3--2---:-4-2----3--1---0----3----1--1--?,-0w32W1Q0-1í-.}
  {  A-tyá-nak fé-nye, fény-su-gár, nap-ból su-gár-zó nap-vi-lág, fény fé-nye, nap-pal nap-ja vagy: for-rás, hol min-den fény fa-kad. Á--------men.}

% Responsorium breve
\defsheet{BenedictusDominusASaeculo}
  {<X-4--3----4t-tT4-,-4-5---rR3-4-3---3-?,-5z-4-:-4-4---4---4--5---rR3-4t-5-?,-[-------------5z--5-:-5-5--4-4-:-5--5-4-----3--4t--tT4-.}
  {   Ál-dott az Úr, * Ö-rök-kön-ö-rök-ké.  A--ki  e-gye-dül mű-vel cso-dá-kat. Dicsőség_az_A-tyá-nak a Fi-ú-nak és a Szent-lé-lek-nek.}

% Benedictus
\defsheet[tone={6-a}]{BenedictusDeus}
  {<-34tg3r-3rR3d1-¨-3--4--5--4--3--3-?,}
  {  Ál-----dott   * Iz-ra-el Is-te-ne.}

% ············································································
% Hora Media
% ············································································

% Antiphona 1
\defsheet[tone={4-a}]{PraeceptumDomini}
  {<-3--1--2---3r-2---¨-eE2--1-:-3---4--5--4--4t--4--3-eE2-1w-2-?,}
  {  Az Úr tör-vé-nye * fény-lő, meg-vi-lá-go-sít-ja a sze-me-ket.}

% Antiphona 2
\defsheet[tone={8szo-c}]{DomineDeusMeusInTe}
  {<-4-3----5--7i-7--7----:-5--7---7---4y-4y-?,}
  {  U-ram, én Is-te-nem, * te-ben-ned bí-zom.}

% Antiphona 3
\defsheet[tone={8szo-c}]{DeusIudexIustus}
  {<-4--4---5z-5--4t-rR3-:-4-5---7--6----5----7--4-,-7--7---6--7---5----4t-:-2e--4t--4-?,}
  {  Is-ten i--gaz bí-ró * e-rős és hosz-szan tű-rő: va-jon ha-rag-szik-e    min-den nap.}

% ============================================================================
% FERIA III
% ============================================================================

% ············································································
% Invitatorium
% ············································································

\defsheet{RegemMagnum}
  {<-1-3----2e-1------,-3-3---4---5--rR3E2-2-.}
  {A nagy Ki-rályt, * i-mád-juk az U-----rat!}

% ············································································
% Officium Lectionis
% ············································································

% Hymnus
\defsheet{ConsorsPaterniLuminis}
  {<ô-2--2--2---1----2-4----3---2í-:-2--3----4---5--4----5---5--6x-,-6--6--5---6---4-----5-rR3-2í--:-2----1--2--4-----5---2---1--2í-?,-1e4R3E2W1-2í-.}
  {   Vi-lá-gos-ság, A-tyád-dal egy, te fény-nek fé-nye, nap-vi-lág, az éj-ből fel-száll é-ne--künk: jöjj el kö-zénk, így kér-le-lünk. Á---------men.}

% ············································································
% Laudes
% ············································································

% Hymnus
\defsheet{PergrataMundoNuntiat}
  {<--5---5---5---5---5---5z---5---5-:--5---5---5---4---3---4t---4---4-;-3---5---5---4---3---4---4---5-:---4---4---4---3---1---3---2---1--?,-1wW1-0\\q\\-.}
  {   Várt haj-nal ad már hírt ne-künk, a nap dár-dá-ja éjt sza-laszt, min-dent a-rany-ba önt a fény, s_a föld szí-nes ru-há-kat ölt.        Á----men.}

% Responsorium breve
\defsheet{DeusMeus}
  {<X-3--4--3--:-3--4t-tT4--,-4-5--rR3-4--3--3-?,-5--5---5z-4-:-5--rR3-4t--tT4-?,-[-------------5z--5-:-5-5--4-4-:-5--5-4-----3--4t--tT4-.}
  {   Is-te-nem, se-gí-tőm, * A-ki-ben re-mé-lek. Ol-tal-ma-zóm és meg-men-tőm.   Dicsőség_az_A-tyá-nak a Fi-ú-nak és a Szent-lé-lek-nek.}

% Benedictus
\defsheet[tone={7la-d}]{Erexit}
  {<-4---4-6---7--8---7---6--7i-8---:-9--ö--9---8---9-8--8--,-8--8-9----8--uU6-7--:-5--5---6---5--4--4-?,}
  {  Föl-e-mel-te né-künk az Is-ten * az üd-vös-ség e-re-jét, az ő szol-gá-já--nak, Dá-vid-nak há-zá-ban.}

% ············································································
% Hora Media
% ············································································

% Antiphona 1
\defsheet[tone={2-b}]{Beati}
  {<X-6---5z-4----:-£-----------------3----4z--4y-?,}
  {   Bol-do-gok, * kik_törvényed_sze-rint jár-nak.}

% Antiphona 2
\defsheet[tone={2-f}]{CantaboDomino}
  {<-1-3--2--3---2--1--1---:-3--1--qQ0-1-eE2--1--1-?,}
  {  É-ne-ke-lek az Úr-nak * ki jó-kat a-dott ne-kem.}

% Antiphona 3
\defsheet[tone={8szo-c}]{LaeteturIacob}
  {<-5-4---3---4t-4----:-5--7--uU6--tT4R3-5z-5--4-?,}
  {  Ö-rül-jön Já-kob, * és uj-jong-jon   Iz-ra-el.}

% ============================================================================
% FERIA IV
% ============================================================================

% ············································································
% Invitatorium
% ············································································

\defsheet{DominumQuiFecitNos}
  {<-3--3r-3---:-3-3--4--3--4t--5----,-5---3--4t---4-4---3-.}
  {  Az Is-tent, a mi al-ko-tón-kat, * jöj-je-tek, i-mád-juk!}

% ············································································
% Officium Lectionis
% ············································································

% Hymnus
\defsheet{RerumCreatorOptime}
  {<ô-4--4---4--4---4--4--3-2í--:-4----4-----4--4---1--2--1--0í---,-2-2---3---4--3---2---1---2eE2í-:-2---2---1---0---1--2--1--1í-?,--1wW1-0q-.}
  {   Vi-lág-te-rem-tő jó U-runk, nézz ránk, ve-zér-lő tá-ma-szunk! A-kik-nek al-vás eny-het ad,     bűn-ben tes-ped-ni so-se hagyd. Á----men.}

  % Antiphona 1
\defsheet[tone={6-a}]{DiligamTeDomine}
  {<-3---4---5---3r-3---eE2-1---:-3--4-4---4--3-?,}
  {  Sze-ret-lek té-ged U---ram * én e-rős-sé-gem.}

% Antiphona 2
\defsheet[tone={1re-a}]{FactusEst}
  {<-3--4tT4-5x--:-5--4---2e--4--eE2--1--0q-1í-?,}
  {  Se-gí---tőm * te let-tél ne-kem, Is-te-nem.}

% Antiphona 3
\defsheet[tone={4-a}]{RetribuitMihiDominus}
  {<-4---3--2---3--qQ0-1e-2--:-2-4---4t--rR3-2---2-?,}
  {  Meg-fi-zet ne-kem az Úr * i-gaz-vol-tom sze-rint.}

% ············································································
% Laudes
% ············································································

% Hymnus
\defsheet{NoxEtTenebrae}
  {<-3--3--4--5----4--4---3--2--:-4--2--3----1--0---1---3--3---,-3--3---4--5---4---4---3-2---:-4-----2---3---1----0--3---1--1-?,-02eE2W1Q0-1-.}
  {  Éj és ho-mály és fel-le-gek, ör-vé-nyek és köd-tor-la-szok, fe-hér az ég, már kél a fény: Krisz-tus van itt, tá-voz-za-tok! Á---------men.}

% Responsorium breve
\defsheet{InclinaCorMeumDeus}
  {<X-¢------------------4--3--4t--5------:-4--3--3-?,-5-5--5-5z-4-:-4--3----4t-5--?,--[-------------5z--5-:-5-5--4-4-:-5--5-4-----3--4t--tT4-.}
  {   Szívemet_tanításod fe-lé for-dítsd, * Is-te-nem. A te u-ta-don él-tess en-gem!   Dicsőség_az_A-tyá-nak a Fi-ú-nak és a Szent-lé-lek-nek.}

% ············································································
% Hora Media
% ············································································

% Antiphona 1
\defsheet[tone={1la-a}]{InTotoCordeMeo}
  {<X-1---0---3r--3---5---:-6--5---4---4--4----3t-4--,-4--4-----5--4---3-:-2--3---4---eE2-1-?,}
  {   Tel-jes szí-vem-ből * ke-res-lek té-ged, U--ram, ne hagyj el-tér-ni  pa-ran-csa-id--tól.}

% Antiphona 2
\defsheet[tone={7do-d}]{InclinaDomine}
  {<-8---9-----7i-8----:-8-6----7--6--4-,-6--6----6u---5---6---5--4-?,}
  {  For-dítsd ho-zám, * U-ram, fü-le-det és hall-gasd meg sza-va-mat.}

% Antiphona 3
\defsheet[tone={5-c}]{EgoAutem}
  {<-tT4-tT4-3---:-3-5---77i-7---7--6--7i--5-:-uU6-5---5-4-----tT4-3-?,}
  {  Én  pe--dig * i-gaz-ság-ban je-le-nek meg szí-ned e-lőtt, U---ram.}

% ============================================================================
% FERIA V
% ============================================================================

% ············································································
% Invitatorium
% ············································································

\defsheet{AdoremusDominum}
  {<-3-3---4---3--4t-5-----,-5-3--4--5----4---3-.}
  {  I-mád-juk az Is-tent, * ő al-ko-tott min-ket!}

% ············································································
% Officium Lectionis
% ············································································

% Hymnus
\defsheet{NoxAtraRerum}
  {<-5-5---5---5---5---5----4--4-,-5--5---5--4-----3--4---5--5-,-5-5---5--4-----3---4----3--2-,-3--4---4-3----4--5---4--4-?,-4tT4-3\\r\\-.}
  {  A föl-dön min-den szín-re még sö-tét fá-tyolt bo-rít az éj. I-gaz Bí-ránk, szí-vünk fe-léd ál-dón e-seng és kér-ve kér. Á----men.}

% ············································································
% Laudes
% ············································································

% Hymnus
\defsheet{SolEcceSurgit}
  {<ô-2---2---2-1----2-4---3--2í-:-2--3----4--5----4--5---5--6x-,-6--6---5--6---4--5--rR3-2í-:-2--1--2--4---5---2----1--2í-?,-1e4R3E2W1-2í-.}
  {   Már kél a nap, a csu-pa tűz, és ránk pi-rít, bá-nat-ra int; ki tud-na bűn-be es-ni  már, ha ez fé-nyé-vel ránk te-kint! Á---------men.}

% Responsorium breve
\defsheet{ClamaviInTotoCorderMeo}
  {<X-3---3---3---4---3---4--4t-tT4--:-5----4----3----rR3-3--,-5--5---5---5z-4-:-5--rR3-4t--tT4--?,--[-------------5z--5-:-5-5--4-4-:-5--5-4-----3--4t--tT4-.}
  {   Tel-jes szí-vem-ből ki-ál-tok, * Hall-gass meg, U---ram. Pa-ran-csa-i--dat én meg-tar-tom. Dicsőség_az_A-tyá-nak a Fi-ú-nak és a Szent-lé-lek-nek.}

% Benedictus
\defsheet[tone={7la-a}]{InSanctitate}
  {<ô-1----1---3----4--5z-5---:-77----zZ5z-5x-,-5--5---5---6--5---rR3-4-:-2--2---3--2--1í--1í-?,}
  {   Szol-gál-junk az Úr-nak * szent-ség--ben, és meg-sza-ba-dít min-ket el-len-sé-ge-ink-től.}

% ============================================================================
% FERIA VI
% ============================================================================

% ············································································
% Invitatorium
% ············································································

\defsheet{DominumDeumNostrum}
  {<-3--3r-3-:--3-3--4--3--4t--5----,-5---3--4t---4-4---3-.}
  {  Az U--rat, a mi Is-te-nün-ket, * jöj-je-tek, i-mád-juk!}

% ············································································
% Officium Lectionis
% ············································································

% Hymnus
\defsheet{TuTrinitasUnitas}
  {<ô-4-----4--4---4---4---4--3--2í--:-4--4---4---4--1-2----1--0í--,-2-----2---3--4---3--2-1--2eE2í-:-2--2---1--0--1----2----1---1í-?,-1wW1-0\\q\\-.}
  {   Szent-há-rom-ság egy Is-te-nünk, ha-tal-mas Úr a föld fe-lett, halld meg di-csé-rő é-ne-künk,   éb-red-ve me-lyet zeng szí-vünk. Á----men.}

% ············································································
% Laudes
% ············································································

% Hymnus
\defsheet{AeternaCaeliGloria}
  {<-3-3---4---5---4--4---3--2--:-4--2---3--1---0---1---3--3---,-3-3---4--5---4---4---3--2-:-4-2----3--1----0---3---1--1-?,-02eE2W1Q0-1-.}
  {  Ö-rök tün-dök-lő Nap-vi-lág, ha-lan-dó nép-nek egy re-mény, a Fön-sé-ges-nek egy Fi-a,  a tisz-ta Szűz-nek mag-za-ta:  Á---------men.}

% Responsorium breve
\defsheet{AuditamFacMihiMane}
  {<X-¢--------------------------4--3----4t--5--:--4--3--rR3-3-?,-5--5-5----4---5--4---4--:-4--3----4t-?,--[-------------5z--5-:-5-5--4-4-:-5--5-4-----3--4t--tT4-.}
  {   Add,_hogy_már_korán_reggel ta-pasz-tal-jam * Jó-sá-go--dat! Az u-tat, mer-re men-jek, mu-tasd meg.   Dicsőség_az_A-tyá-nak a Fi-ú-nak és a Szent-lé-lek-nek.}

% Benedictus
\defsheet[tone={5-c}]{VisitavitEtFecit}
  {<-3---5--7--7---uU6-¨-tT4-7---8---5---4--tT4-3--4--4t-3--3-?,}
  {  Meg-lá-to-gat-ta  * és  meg-vál-tot-ta az  Úr az ő  né-pét.}

% ============================================================================
% SABBATUM
% ============================================================================

% ············································································
% Invitatorium
% ············································································

\defsheet{PopulusDomini}
  {<-1--2--4t-5-:-5--4--5--6--5--4---5--5--5--,-6u--6--5----zZ5T4-5z--6---5-.}
  {  Az Úr né-pe, és le-ge-lő-jé-nek ju-ha-i, * jöj-je-tek, i-----mád-juk őt.}


% ============================================================================
% DOMINICA
% ============================================================================

% ············································································
% Vesperas I.
% ············································································

% Antiphona 3
\defsheet[tone={1la-a}]{DeditIlliDominus}
  {<X-1-1---1--0---34t-5---:-6-5----4--4--3t-4-,-4--2-4---tT4-3-:-3--3-----2--3--4-eE2-1--1-?,}
  {   Ö-rök di-cső-sé--get * a-dott ne-ki az Úr, és ö-rök ne--vet ha-gyott ne-ki ö-rök-sé-gül.}

% ············································································
% Officium lectionis
% ············································································

% Hymnus
\defsheet{MediaeNoctis}
  {<-1----0--1---3----4--3--2---1-:-4---4---3-4---5--3---2---1-,-1--1----2---0--ð----0----1--1--:-1-0----1----3--4-3----2--1-?,-1wW1-0q-.}
  {  Csak fé-lig múlt az éj-sza-ka, szó-lít a pró-fé-ták sza-va: di-csér-jük Is-tent szün-te-len, A-tyát s_Fi-út e-gyüt-te-sen, Á----men.}

% vel:
% TODO

% Antiphona 1
\defsheet[tone={3-c}]{Confessionem}
  {<-4---4ç6u-6--:-5--7---6--5---5--6--5--4--,-4-----5z-6----5--4---tT4-:-2e---rR3-2-?,}
  {  Fön-ség--be * és mél-tó-ság-ba öl-tö-zöl, fényt öl-tesz ma-gad-ra,   mint ru-hát.}

% Antiphona 2
\defsheet[tone={8szo-c}]{DomineDeusMeus}
  {<-rR3E2-1----3r-5--4----:-4--5----4---3----5u-6--4-?,}
  {  U-----ram, Is-te-nem, * ma-gasz-tos vagy i--ga-zán.}

% Antiphona 3
\defsheet[tone={tirr-e}]{QuamMagnificataSuntE}
  {<-2----2---1--2--1---:-1-0--1w-3--1---3-2-?,}
  {  Mily fön-sé-ge-sek * a te mű-ve-id, U-ram.}

% ············································································
% Laudes
% ············································································

% Hymnus
\defsheet{EcceIamNoctis}
  {<X-5--4----5--4---1í-:-4--4--3---4--5x-5x-,-5---7---8-7---5x-:-5--4---3---4--5x---5x-,-5---4---5----4---1í-:-4-4--3---4-tT4-3y-,-3---3---2-1í---1í-?,-1wW1-0\\q\\-.}
  {   Sá-pad, el-osz-lik  kö-de már az éj-nek, pir-kad a haj-nal, ró-zsa-szí-nű fény kél; for-dul-junk buz-gón  a vi-lág U-rá--hoz  han-gos i-mánk-kal,  Á----men.}

% Antiphona 1 (Nb41)
\defsheet[tone={8szo-c}]{DominusMihiAdiutor}
  {<-2--3---4--4--:-5--4--3--4t--eE2-1í-,-3-4----5---4--45uU67iI7-:-5---5---4---3--4---5--4x-4x-?,}
  {  Ve-lem az Úr * az én se-gít-sé--gem, a-zért nem fé-lek,        mit árt-hat né-kem az em-ber?}

% Antiphona 2
\defsheet[tone={8szo-c}]{TresExUno}
  {<-4---5---¨-4---5u---uU6Z5-:-iI7-iol7i-7---7iI77U6Z5-3-5--5---5--uU6--tT4-5-4-,-5--5uU6-tT4-3---3--:-4--5----6--5--4-?,}
  {  Hár-man * egy száj-jal     ki-ál-----tot-tak       a ke-men-ce láng-ja--i-ban és é----ne--kel-ték: ál-dott az Is-ten.}

% Antiphona 3
\defsheet[tone={8szo-c}]{LaudateDominum}
  {<-4--4----4--4---3t-7-6---¨-4----6--uj5-rR3-:-5--5--7--5--zZ5-4-?,}
  {  Di-csér-jé-tek az U-rat * nagy-sá-gá--nak   so-ka-sá-ga sze-rint.}

% Responsorium breve
\defsheet{ConfitebimurTibiDeus}
  {<X-3--3----3--3----4--3--:-4--4t-5-----,-4--4--5--rR3-4---3---3-?,-5--5--5z---4-:-5---5--4---3--4t-5-?,-5--5---5-.}
  {   Ma-gasz-ta-lunk té-ged, Is-te-nünk, * És ne-ve-det szó-lít-juk. El-be-szél-jük cso-da-tet-te-i--det. Di-cső-ség...}

% ············································································
% Vesperas II.
% ············································································

% Antiphona 1 (SzKD43)
\defsheet[tone={8szo-c}]{IuravitDominus}
  {<-4---6--8--8----7i-7---:-7--6---4---zZ5-4-,-44--eE2--1--3----4t-4x--4x-?,}
  {  Meg-es-kü-dött az Úr, * és meg nem bán-ja: pap vagy te mind-ö--rök-ké.}

% ============================================================================
% FERIA II
% ============================================================================

% ············································································
% Officium lectionis
% ············································································

% Hymnus
\defsheet{IpsumNuncNobis}
  {<-1----0---1--3-4---3-2--1-:-4--4-3---4--5--3--2---1--,-1---1--2--0---ð-0--1--1--:-1--0-----1--3---4---3-2--1-?,-1wW1-0q-.}
  {  Most itt az ó-ra, a-mi-kor az e-van-gé-li-um sze-rint meg-ér-ke-zik a Vő-le-gény és menny-or-szá-got a-la-pít. Á----men.}

% Antiphona 1
\defsheet[tone={8szo-c}]{InTuaIustitia}
  {<-7-7--7--7---5--4---45u-7---:-7---7--7----7---5--4----5-4-?,}
  {  A te ig-gaz-sá-god ál--tal * sza-ba-díts meg en-gem, U-ram!}

% Antiphona 2
\defsheet[tone={8szo-c}]{FaciemTuam}
  {<-5-3--5--7--6----4-5----:-7--7----7----7---6----7---5--4-?,}
  {  A te ar-co-dat, U-ram, * ra-gyog-tasd fel szol-gád fö-lött!}

% Antiphona 3
\defsheet[tone={2-f}]{DiligiteDominum}
  {<-1---1---1--0---1--eE2-1----:-3---3---3-----4--5-:-3--3---3--2-3--1---0---0q--3----3r-eE2-1-?,}
  {  Sze-res-sé-tek az U---rat, * min-den szent-je-i,  mi-vel az i-ga-za--kat meg-tart-ja az  Úr.}

% ············································································
% Laudes
% ············································································

% Hymnus
\defsheet{LucisLargitorSplendide}
  {<ô-4--4----4---4---4---4--3--2í-:-4--4-----4-4-----1--2---1--0í-,-2--2--3---4---3----2---1--2eE2í-:-2---2---1--0--1--2---1--1í-?,-1wW1-0q-.}
  {   Fé-nyes-ség-osz-tó, bő-ke-zű,  ki fényt á-raszt-va gaz-da-gon  az éj-nek gyá-szát osz-la-tod,    s_a nap vi-lá-ga tűz le ránk. Á----men.}

% Antiphona 1
\defsheet[tone={8szo-c}]{Sitivit}
  {<-4u---6--4---¨-6-7--5---4-tT4-3--:-5z-5-4-?,}
  {  Szom-ja-zik * u-tá-nad a lel-kem, én U-ram!}

% Antiphona 2
\defsheet[tone={3-c}]{Ostende}
  {<-4--5z---6---¨-7--tT4---5z-tT4-:-4çz-5---4---4--3----2--2-?,}
  {  Mu-tasd meg * ne-künk, U--ram,  ir--gal-mad fé-nyes-sé-gét!}

% Antiphona 3
\defsheet[tone={2-f}]{Opera}
  {<-3--eE2-1---¨-33r-3---3-:-3tT4t-4--eE2-0-äe----1-?,}
  {  Ke-zé--nek * mű--ve--it  hir---de-ti  a menny-bolt.}

% Responsorium breve
\defsheet{ExultateIustiInDomino}
  {<X-¢----------------4----3--4t-tT4--,-[---------------4---3---rR3-3-?,-5-5z--4-:-5--5----5-4--3---4t-5-?,-5--5---5-.}
  {   Örvendjetek,_iga-zak, az Úr-ban, * A_szentekhez_di-csé-ret il--lik. E-lőt-te  új dalt é-ne-kel-je-tek. Di-cső-ség...}

% Benedictus
\defsheet[tone={4-a}]{BenedictusDominusQuia}
  {<-2--1----wW1-0---:-2----ed1-2--3--4t-5--:-tT4-3---4t--4----2---2-?,}
  {  Ál-dott az  Úr, * mert meg-lá-to-ga-tott és  meg-vál-tott min-ket.}

% ············································································
% Hora Media
% ············································································

% Hymnus
\defsheet{DicamusLaudesDomino}
  {<-5--5----5---5---5---5--4-4-,--5---5--5----4---3---4---5--5--,-5--5-5---4-3---4---3--2-,-3-----4---4---3-4---5---4--4-?,-4tT4-3\\r\\-.}
  {  Di-csér-jük dal-lal az U-rat, buz-gó kész-ség-gel lel-ke-sen: az ó-ra, í-me, dél-re jár s_ben-nün-ket i-mád-ság-ra hív. Á----men.}
  
% ============================================================================
% FERIA III
% ============================================================================

% ············································································
% Officium lectionis
% ············································································

% Hymnus
\defsheet{NocteSurgentes}
  {<X-5---4----5---4--1í-:-4----4--3---4---5x--5x---,-5---7--8-7----5x-:-5---4--3--4---5x--5x-,-5----4---5---4-----1í-:-4-4--3---4---tT4--3y-,-3-3---2-1í---1í-?,-1wW1-0\\q\\-.}
  {   Kel-jünk fel éj-jel  s_kö-zö-sen vir-ras-szunk, egy-re a zsol-tár  sza-va-it fon-tol-va,  szár-nya-lón száll-jon  a di-csé-ret hang-ja   é-des U-runk-hoz.  Á----men.}

% Antiphona 1
\defsheet[tone={6-a}]{Revela}
  {<-3----4---5--4--3---:-4-3--3---?,}
  {  Tárd fel az Úr-nak * u-ta-dat.}

% Antiphona 2
\defsheet[tone={2-f}]{IuniorFui}
  {<-eE2-1--2e--2---¨-1--1---1--wW1-1--:-0q-3---3r--3---3-4--3---wW1-eE2--1-?,}
  {  If--jú vol-tam * és meg-vé-nül-tem, de nem lát-tam i-ga-zat el--hagy-va.}

% Antiphona 3
\defsheet[tone={6-a}]{SperaInDomino}
  {<-3---3---2--4--5---:-5--5-4---5---4--3-?,}
  {  Bíz-zál az Úr-ban * és a jót cse-le-kedd.}

% ············································································
% Laudes
% ············································································

% Hymnus
\defsheet{AeternaeLucisConditor}
  {<ô-2----2--2--1--2---4--3--2í--:-2-3---4--5--4-----5---5--6x-,-6--6---5---6----4-----5--rR3-2í--:-2---1--2---4---5---2---1--2í-?,-1e4R3E2W1-2í-.}
  {   Fény al-ko-tó-ja, Is-te-nünk, ö-rök ve-rő-fény, nap-vi-lág, va-lód-ból csak fényt is-me--rünk, nem ér fel hoz-zád vak ho-mály. Á---------men.}

% Antiphona 1
\defsheet[tone={6-a}]{Salutare}
  {<-3--4t-3---:-3--4---eE2-1---3--4--4--3-?,}
  {  Én or-cám * üd-vös-sé--ge, én Is-te-nem.}

% Antiphona 2
\defsheet[tone={3-c}]{CunctisDiebus}
  {<-4-4--5u---7---¨-5---7---zZ5-4-:-4-----3---4---5----rR3E2-2-?,}
  {  É-le-tünk-nek * min-den nap-ján ments meg min-ket, U-----runk!}

% Antiphona 3
\defsheet[tone={2-f}]{TeDecetHymnus}
  {<-3--3---3--2----3--1----:-1-1--1---1---0--1e-1-?,}
  {  Té-ged il-let, Is-ten, * a di-csé-ret Si-on-ban!}

% Responsorium breve
\defsheet{VocemMeamAudiDomine}
  {<X-¢----------------4---3---4t-5----,-4----4--4----3--4--5---4-ed1-3r-3-?,-[------------4--5--4--4-:-5---4---3---4t-5--?,-5--5---5-.}
  {   Tekints_könyörgé-sem-re, U--ram, * Mert na-gyon re-mé-lek i-gé--id-ben. Téged_szólít ki-ál-tá-som már haj-nal e--lőtt. Di-cső-ség...}

% Benedictus
\defsheet[tone={1re-a}]{DeManuOmnium}
  {<-3----5-4---5--3r-4----:-1-1---3---3---2---2e--3-:-4---4--4----2---eE2-1----0q-1-?,}
  {  Mind-a-zok ke-zé-ből, * a-kik gyű-löl-nek min-ket sza-ba-díts meg min-ket, U--runk.}

% ============================================================================
% FERIA IV
% ============================================================================

% ············································································
% Officium lectionis
% ············································································

% Hymnus
\defsheet{OSatorRerum}
  {<X-5---4--5---4---1í-:-4--4---3-4--5x-5x-,-5----7--8-----7-----5x-:-5--4---3--4---5x-5x-,-5---4--5--4---1í--:-4-4--3---4--tT4-3y-,--3----3-2---1í---1í-?,-1wW1-0\\q\\-.}
  {   Min-de-nek Aty-ja,  vi-lág ú-jí-tó-ja,  nagy Ki-rály, Krisz-tus, fé-lel-me-tes Bí-ró,  áld-va di-csé-rünk, e-se-dez-ve ké--rünk, nézz a szí-vünk-re!   Á----men.}

% Antiphona 1
\defsheet[tone={6-a}]{UtNonDelinquam}
  {<-3--4---5---3---4----eE2-1----:-3---4----4-3-?,}
  {  Ne vét-kez-zem nyel-vem-mel, * add meg, U-ram!}

% Antiphona 2
\defsheet[tone={7do-a}]{NonSisMihi}
  {<ô-5--3----5z-5---:-3--4t--wW1-1í-,-4--4--4--2r-3y-:-3-2---1---2--4---3---1í--1í-?,}
  {   Ne légy né-kem * fé-lel-mem-re,  én re-mé-sé-gem  a szo-ron-ga-tás-nak nap-ján.}

% Antiphona 3
\defsheet[tone={4-a}]{Expectabo}
  {<-3--1----2--3---¨-4--3---2---1e-2--,-eE2W1-23r-eE2-?,}
  {  Re-mény-ke-dem * ne-ved-ben U--ram, mert  jó  az.}

% ············································································
% Laudes
% ············································································

% Hymnus
\defsheet{FulgentisAuctor}
  {<-5--5--5--5--1---3--2--1í--:-5--5--6---4----5-----7----6--5x-,-5-----6---7--8----7----5---7--zZ5T4x-:-5---5--6---4--2---4---5--5x-?,-5zZ5-4\\t\\-.}
  {  Ég Al-ko-tó-ja, Is-te-nünk, ki éj-jel hold-fényt adsz ne-künk s_nap-pal na-pot, mely kör-be fut,     s_e-lő-írt út-ján cél-ba jut.  Á----men.}

% Antiphona 1
\defsheet[tone={8szo-c}]{DeusInSanctoViaTua}
  {<-4u-5----¨-4-----5-rR3-4t-4-,-6u-8-7----5----7--6--:-4----6u-5--4--4--4-?,}
  {  Is-ten, * szent a te  U--tad ki o-lyan nagy is-ten, mint a  mi Is-te-nünk?}

% Antiphona 2
\defsheet[tone={8szo-c}]{Exsultavit}
  {<-7--7---6u-5---:-5--5--5---4---3--4t-4-?,}
  {  Ör-ven-de-zik * az én szí-vem az Úr-ban.}

% Antiphona 3
\defsheet[tone={3-c}]{DominusRegnavit}
  {<-4--5u-uU6-tT4-5---¨-4--3---4t--4---rR3-2-?,}
  {  Az Úr or--szá-gol * ör-ven-dez-zen a --föld!}

% Responsorium breve
\defsheet{BenedicamDominum}
  {<X-3--4---3--4t-tT4-,-5---rR3-4-3--3-?,-5--5---4--4-:-5----5--5--5---4---3--4t-tT4-?,-5--5---5-.}
  {   Ál-dom az U--rat * Min-den i-dő-ben. Di-csé-re-tét szün-te-le-nül zen-gi aj-kam.   Di-cső-ség...}

% Benedictus
\defsheet[tone={8szo-c}]{Liberati}
  {<-4---4---5--5z--zZ5-:-5----7---uU6--5--zZ5-4t---7--6--5----5-:-5z----5---4-?,}
  {  Meg-sza-ba-dul-ván * szol-gál-junk az Úr--nak, Is-te-nünk-nek szent-ség-ben.}

% ============================================================================
% FERIA V
% ============================================================================

% ············································································
% Officium lectionis
% ············································································

% Hymnus
\defsheet{AlesDieiNuntius}
  {<-3-3---4---5----4----4---3--2-:-4--2---3--1----0----1--3-3---,-3-3---4---5--4----4--3--2-:-4-----2----3---1---0---3--1--1-?,-0\\2eE2W1Q0-1\\-.}
  {  A nap-pal szár-nyas hír-nö-ke  je-len-ti már, hogy jő a fény; a lel-kek éb-resz-tő-je is, Krisz-tus, ben-nün-ket él-ni hív: Á-----------men.}

% Antiphona 1
\defsheet[tone={8do-c}]{Extende}
  {<-5---3r----4---:-4t-4-----rR3-4---4t-3--,-5--¦-----------------6--7--5--:-3--4--5z---5---4-?,}
  {  Ter-jeszd ki, * U--runk, ha--tal-ma-dat, és szabadíts_meg_lel-ke-in-ket, el ne pusz-tul-junk.}

% Antiphona 2
\defsheet[tone={2-c}]{Confundantur}
  {<-4---5---7---7--7i--8----:-7-8---9--iI7-8--8---8--,-8--8--5--8---7---6---tT4-5--:-4--3-4----5--7--4--4--?,}
  {  Szé-gye-nül-je-nek meg, * a-kik en-gem ül-döz-nek, és én ne szé-gye-nül-jek meg, én U-ram, én Is-te-nem.}

% Antiphona 3
\defsheet[tone={2-f}]{MementoDomine}
  {<-1wW1Q0-1---eE2-3rR3-2eE2W1-:-3--3--4----5---4---3---4--3r-4-:-3-1----qQ0-3-eE2-1-?,}
  {  Em-----lék-ezz U----ram,   * és mu-tasd meg nek-ünk ma-ga-dat a szük-ség i-de--jén!}

% ············································································
% Laudes
% ············································································

% Hymnus
\defsheet{IamLucisOrto}
  {<-7---7---7-7----7---7----6--5-:-7-7----7--7---4---5--4-3--,-5---5---6--7---6--5----4--5zZ5-:-5--5----4--3--4---5-4--4-?,--4tT4-3\\r\\-.}
  {  Már kél a fény-nek csil-la-ga: e-seng-ve kér-jük az U-rat, jár-jon ma min-de-nütt ve-lünk,  ne ront-sa ár-tás é-le-tünk. Á----men.}

% Antiphona 1
\defsheet[tone={2-d}]{TuamDomine}
  {<-7---7i-----8----:-7--8---9--iI7--8-8---,-5i-uU6--tT4-:-5----5---7---5--4----4t--5-?,}
  {  Ser-kentsd fel, * ha-tal-ma-dat, U-runk, és jöjj el,   hogy meg-sza-ba-díts min-ket!}

% Antiphona 2
\defsheet[tone={8szo-c}]{ConversusEst}
  {<-4---4---4----5-4--5--4--3----:-3-----5--7----7---4---3r-4-?,}
  {  Meg-for-dult a te ha-ra-god, * s_meg-vi-gasz-tal-tál en-gem.}

% Antiphona 3
\defsheet[tone={1re-a}]{ExsultateDeo}
  {<-5--5---5---5--4---5--4---rR3--:-2-3--4--eE2-1----1-?,}
  {  Ör-ven-dez-ze-tek Is-ten-nek, * a mi se-gí--tőnk-nek!}

% Responsorium breve
\defsheet{InMatutinisDomine}
  {<X-3---4--3----4t-5----,-5--5---4--3---rR3-3-?,-5z-4-:-5--5---4--3---4t--tT4-?,-5--5---5-.}
  {   Haj-na-lig, U--ram, * Ró-lad el-mél-ke--dem. Mi-vel vé-del-me-zőm let-tél.   Di-cső-ség...}

% Benedictus
\defsheet[tone={7do-d}]{DaScientiam}
  {<-6--7-----8--7--6----7-6----,-5--rR3-4--5--7iI7-:-7--5--4--4-?,}
  {  Ta-nítsd né-pe-det, U-ram; * bo-csá-na-tá-ra     bű-ne-ik-nek.}

% ============================================================================
% FERIA VI
% ============================================================================

% ············································································
% Officium lectionis
% ············································································

% Hymnus
\defsheet{GalliCantu}
  {<X-3---4---4---1---3-4--4---3--:-6--5--4--5---3r--5-4--,-5--6--7--4---6---5--4--3-:-6----5---4--3---5-zZ5-4-,-3--4--4---1---3---4---4--3---:-6-----5---4--5-----3r---5--4-?,--4tT4-3r-.}
  {   Fel-har-san már a ka-kas-szó, éj sö-té-tet old e jel, és az éj-fél vak ré-me-it  mint fel-le-get ű-zi  el. Jó-sá-gos Úr, hoz-zád té-rünk, s_buz-gón ké-rünk, légy ve-lünk. Á----men.}

% Antiphona 1
\defsheet[tone={8do-c}]{NeInIraTua}
  {<-7---7-uU6-5---:-7--7---7----8--8---7-7-?,}
  {  Har-a-god-ban * ne bün-tess en-gem U-ram.}

% Antiphona 2
\defsheet[tone={2-d}]{ConfundanturEtRevereantrur}
  {<-4---5---7---7--7i--8---:-7i-9---iI7-8--,-8-5---8---7---6--tT4-5--:-3r---5--7--5---5-?,}
  {  Szé-gye-nül-je-nek meg * és fél-je--nek, a-kik lel-kem ke-re--sik, hogy el-ra-gad-ják.}

% Antiphona 3
\defsheet[tone={8szo-c}]{NeDerelinquasMe}
  {<-4--4-----5--7--7---¨-7--7-7----7--7--7--7--:-7--7--zZ5--zZ5-4--4-?,}
  {  Ne hagyj el en-gem * én U-ram, én Is-te-nem, ne tá-vozz-el  tő-lem.}

% ············································································
% Laudes
% ············································································

% Hymnus
\defsheet{DeusQuiCaeliLumen}
  {<X-3--3----4--5-----4---4--3---2-:--5--2---3---1----0---1---3--3-,-3-3------4-5---4----4----3--2-:-4--2-----3--1---0---3----1--1-?,-0\\2eE2W1Q0q\\-.}
  {   Is-ten, ki menny-nek fé-nye vagy és min-den fény-nek kút-fe-je, A-tyánk, e-get tart fenn ka-rod és menny-or-szá-got nyit ke-zed, Á----men.}

% Antiphona 1
\defsheet[tone={1la-a}]{AmpliusLavaMe}
  {<-0----1---3--2----1-qQ0--:-3---4---tT4-3--,-3----4----2---3--2--1--1-?,}
  {  Moss meg en-gem, U-ram, * tel-jes-ség-gel, tisz-títs meg bű-ne-im-től!}

% Antiphona 2
\defsheet[tone={tirr-e}]{DomineAudivi}
  {<-2-1w---3---1---1---:-0-1--3---3--2---1--2---2-?,}
  {  U-ram, hal-lot-tam * a te sza-va-dat és fél-tem.}

% Antiphona 3
\defsheet[tone={tirr-e}]{Lauda}
  {<-2--1------¨-1--0--1---2e---1--3-2-?,}
  {  Di-csérd, * Je-ru-zsá-lem, az U-rat!}

% Responsorium breve
\defsheet{ClamaboAdDominumAltissimum}
  {<X-¢-------------4---3---4--4t-5----,-£--------5--rR3-4---3--3-?,-5--5----5--5z-4--:-5--5---5---4--3---4t-5-?,-5--5---5-.}
  {   A_fölséges_Is-ten-hez ki-ál-tok, * Aki_jóra vi-szi sor-so-mat. El-küld az ég-ből, és meg sza-ba-dít en-gem. Di-cső-ség...}

% Benedictus
\defsheet[tone={1re-b}]{PerViscera}
  {<ô-1--2---2z-6---:-6----6--5--4-5z-6-,-6---6--6--5--4----4t--4x-:-4-3--4---5---rR3-2--1w-2y-?,}
  {   Ir-gal-má-nak * mély-sé-ge-i ál-tal meg-lá-to-ga-tott min-ket  a ma-gas-ság-ból tá-ma-dó.}

% ············································································
% Vesperas
% ············································································

% Hymnus
% TODO

% Antiphona 1
\defsheet[tone={1re-a}]{InclinavitDominus}
  {<-3---3----4--tT4-5-:-2-3---4--eE2-1---1---?,}
  {  Meg-hall-ja az  Úr  i-mád-sá-gom sza-vát.}

% Antiphona 2
\defsheet[tone={2-f}]{AuxiliumMeum}
  {<-3--eE2W1-1-:-0-1--3--3---0e-1----?,}
  {  Az Úr-tól    a mi se-gít-sé-günk.}

% vel ad libitum:
% TODO

% Antiphona 3
\defsheet[tone={8szo-c}]{InConsilioIustorum}
  {<-£---------5--4--5---4-:-3t-7---7--6--5--4u-zZ5-:-4--4----5--5--4--4--4-?,}
  {  Az_igazak ta-ná-csá-ban és gyü-le-ke-ze-té-ben   na-gyok az Úr mű-ve-i.}

% ============================================================================
% SABBATO
% ============================================================================

% ············································································
% Officium lectionis
% ············································································

% Hymnus
\defsheet{LuxAeterna}
  {<-3-4---4--1----3----4--4--3---:-6----5--4---5---3r--5--4--,-5-6---7----4---6--5---4---3-:-6-5-----4--3--5---zZ5-4--,-3---4---4---1--3---4--4---3-:-6--5----4---5---3r--5--4-.}
  {  Ö-rök su-gár, szép ve-rő-fény, fogy-ha-tat-lan nap-vi-lág, a gyá-szos éjt te győ-zöd le, a fényt új-ra meg-ho--zod, pok-lun-kat te rom-bo-lod le, el-ménk ben-ned föl-de-rül.}

% Antiphona 1 (Extra tempus per annum)
\defsheet[tone={8do-c}]{VisitaNosDomine}
  {<-7--7--5----4---45u-7-:-7--5---4--5---4-?,}
  {  Lá-to-gass meg min-ket üd-vös-sé-ged-del!}

% Antiphona 2 (Extra tempus per annum)
\defsheet[tone={1la-a}]{OblitiSuntDeum}
  {<-0--1--3---3---3---4t-4---:-4--2---3r--2---2e-1--1-?,}
  {  El-fe-lej-tet-ték Is-tent, ki meg-men-tet-te ő-ket!}

% Antiphona 3 (Extra tempus per annum)
\defsheet[tone={1la-a}]{SalvosNosFacDomine}
  {<-0-----1---2e--4----eE2-1--:-3------eE2-4t--eE2-qQ0-1-1---eE2-1---0q-1-?,}
  {  Ments meg min-ket, U---ram, gyűjts ösz-sze min-ket a nem-ze--tek kö-zül!}

% Antiphona 1 (Tempore per annum)
\defsheet[tone={3-c}]{ConfiteminiDomino}
  {<-4--5--5---5u-6---5--4--tT4-:--2e---4t-2---2-?,}
  {  Ad-ja-tok há-lát az Úr-nak, * mert jó hoz-zánk.}

% Antiphona 2 (Tempore per annum)
\defsheet[tone={3-c}]{Quoniam}
  {<-7-uU6-tT4-5--4--:-4--4---4---rR3-4t-rR3-2--2-?,}
  {  Ö-rök-ké--va-ló * ir-gal-mas-sá--ga az  Úr-nak.}

% Antiphona 3 (Tempore per annum)
\defsheet[tone={1re-a}]{TuEsDeus}
  {<-5--4----5--4--3----:-3--2---3r-eE2-1---0q--1-?,}
  {  Te vagy az Is-ten, * ki cso-dá-kat cse-lek-szel.}

% ············································································
% Laudes
% ············································································

% Hymnus
\defsheet{DieiLuceReddita}
  {<-7-7---7----7--7---7---6--5---:-7----7---7---7--4--5-4--3---,-5--5----6---7--6---5----4--5zZ5-:-5-----5---4--3---4-----5--4--4-?,-4tT3-2\\e\\-.}
  {  A nap-fény új-ból fel-ra-gyog, csen-dül-jön há-la-é-ne-künk, di-csér-jük Is-ten nagy ne-vét,   Krisz-tus ke-gyel-mét, ér-de-mét. Á----men.}

% Antiphona 1
\defsheet[tone={8szo-c}]{QuamMagnificataSunt}
  {<-7----6---7--5--4---:-6-7--6--5--4---5-4-?,}
  {  Mily fön-sé-ge-sek * a te mű-ve-id, U-ram.}

% Antiphona 2
\defsheet[tone={6-a}]{DateMagnitudinem}
  {<-3--4----4---eE2-1---:-3--4--4---3-?,}
  {  Ma-gasz-tal-já--tok * Is-te-nün-ket.}

% Antiphona 3
\defsheet[tone={1re-a}]{QuamAdmirabileEst}
  {<-0----1---1tz-5--tT4-:-3-4--5--rR3-3r-4-:--3--4-eE2--1----1-?,}
  {  Mily cso-dá--la-tos * a te ne-ved U--runk az e-gész föld-ön.}

% Responsorium breve
\defsheet{ExultabuntLabiaMea}
  {<X-3--3---4--3---4t--tT4--,-4-4--4---5--rR3-4-3--3-?,-[-----------5z--4--:-[-----------4--3--4t--5--?,-5--5---5-.}
  {   Aj-kam ör-ven-dez-zen, * A-mi-kor ne-ked é-ne-kel. Nyelvem_nap-hos-szat igazságodat em-le-ges-se!. Di-cső-ség...}

% Benedictus
\defsheet[tone={1re-a}]{InViamPacis}
  {<-3--3-5---rR3-4--4--5--:-3--2----4---eE2-1-1-?,}
  {  Az i-gaz-ság út-já-ra * ve-zess min-ket U-runk.}


% ============================================================================
% DOMINICA
% ============================================================================

% ············································································
% I. Vesperas
% ············································································

% Hymnus
% Ut in hebdomada I

% Antiphona 1 (ÉE789)
\defsheet[tone={7do-a}]{SitNomenDomini}
  {<ô-4--4----5zZ5-4----:-4--4--4--2--3----2-1---1-?,}
  {   Le-gyen ál---dott * az Úr ne-ve mind-ö-rök-ké.}

% ············································································
% Officium lectionis
% ············································································

% Hymnus
% Ut in hebdomada I

% Antiphona 1
\defsheet[tone={8szo-c}]{EcceDeusMeus}
  {<-4-rF1-3--4t-4--4---:-4--tT4-3---4t--4-,-5-7---7--4--eE2-3-:-1--2--3---4t-5---4-?,}
  {  Ő az  én Is-te-nem * és di--csé-rem őt, a-tyá-im Is-te--ne  és di-cső-í--tem őt.}

% Antiphona 2
\defsheet[tone={1re-a}]{RegnumTuumDomine}
  {<X-3-4--5--6---5---4-5----:-5-4---2--3--4--eE2-1--?,}
  {   A te or-szá-god U-ram, * ö-rök-ké va-ló or--szág.}

% Antiphona 3
\defsheet[tone={8do-c}]{InAeternum}
  {<-7---uU6-5---:-rR3-4----5-6---5---5-4t--5-?,}
  {  Min-den-kor * és  mind-ö-rök-kön ö-rök-ké.}

% ············································································
% Laudes
% ············································································

% Hymnus
% Ut in hebdomada I

% Antiphona 1
\defsheet[tone={4-a}]{RegnavitDominusDecorem}
  {<-2--1---2---4--4t--:-5-5---4---3--4--3--2-?,}
  {  Or-szá-gol az Úr, * é-kes-ség-be öl-tö-zött.}

% Antiphona 2
\defsheet[tone={7la-d}]{Benedicite}
  {<-4uU6-7--iI7-8--8o-8--:-89pP9-8--7---5i--7--6--5--6--5z-4--4--?,}
  {  Áld--já-tok az U-rat * az    Úr min-den mű-ve-i, al-le-lu-ja.}

% Antiphona 3
\defsheet[tone={1re-a}]{Laudate}
  {<-5--5----5--4---:-2--3-4---3-eE2-1-?,}
  {  Di-csér-jé-tek * az U-rat a men-nyből.}

% ············································································
% Hora Media
% ············································································

% Antiphona 1 (Nb41)
\defsheet[tone={8szo-c}]{InTribulatione}
  {<-4-4---3---2--3r--4---:-£----------5---4---5--4-3y-,-3t-7---7----7--5----7--7c-:-8--7--5---4---3r-4í-?,}
  {  A szo-ron-ga-tás-ban * segítségül hív-tam az U-rat, és meg-hall-ga-tott en-gem, és tá-gas-ság-ra vitt.}

% Antiphona 2 (Nb42)
\defsheet[tone={2-g}]{Circumdantes}
  {<ô-4--3---2--2---:-4--4---2---4---5z-5x-,-5--5--6--5--rR3-2-:-2--1----ñ---1w-2-?,}
  {   Kö-rül-vé-vén * kö-rül-vet-tek en-gem, de az Úr ne-vé--ben el-vesz-tém ő--ket.}

% Antiphona 3 (Nb40)
\defsheet[tone={2-g}]{ODomine}
  {<ô-4--rR3E2-2----:-2--2---2--1---2r-4x-,-4-5zZ5-4--:-3--4---5--4-3--2--1w-2-?,}
  {   Ó, U-----ram, * ad-jál üd-vös-sé-get, ó U----ram, ad-jál jó e-lő-me-ne-telt.}

% ············································································
% II. Vesperas
% ············································································

% Antiphona 1 (NZS14)
\defsheet[tone={1la-a}]{DixitDominus}
  {<-11--0--1--3---wW1-q\\Q0\\-:-3---4--tT4-4x-3r-2e\\-2-1í--1í-?,}
  {  Ülj az én job-bom-ra,     * mon-dá az  Úr az én   U-ram-nak.}

% Antiphona 2 (NZS14)
\defsheet[tone={4-a}]{Fidelia}
  {<-2-4t-5---:-44---2--3-4---eE2-1e-3y-,-3---3---3---01e-3y-:-3-2--3---4---3---2-1w--2-?,}
  {  E-rő-sek * mind az ő vég-zé--se-i;   szi-lár-dan áll-nak, i-ga-zak min-den i-dők-re.}

% ============================================================================
% FERIA II
% ============================================================================

% ············································································
% Officium lectionis
% ············································································

% Hymnus
% Ut in hebdomada I

% Antiphona 1 (PSI86)
\defsheet[tone={8do-c}]{DeusDeorum}
  {<-7--uU6-5--4---5--4--4---:-4t-zZ5-4---3r--4-?,}
  {  Az is--te-nek Is-te-ne, * az Úr  meg-szó-lalt.}

% Antiphona 2
\defsheet[tone={tirr-a}]{ImmolaDeo}
  {<X-5--4----4--5--4---4---:-3-4--5---6---4--6--6--5--?,}
  {   Ál-dozd az Is-ten-nek * a di-csé-ret ál-do-za-tát!}

% Antiphona 3
\defsheet[tone={7do-d}]{DixiIniquis}
  {<-7---7----8---6-7--8oO8-7----7----:-5--7--6----4t-4---6u-6----5--5---rR3-3t--4--4-?,}
  {  Azt mond-tam a go-no---szok-nak: * ne be-szél-je-tek go-nosz-sá-got Is--ten el-len.}

% ············································································
% Laudes
% ············································································

% Hymnus
% Ut in hebdomada I

% Antiphona 1
\defsheet[tone={8szo-c}]{BeatiQuiHabitant}
  {<-4---4u--5---rR3-4t--5--4--:-3--4--5---6---tT4-?,}
  {  Bol-dog-ok, a---kik az Úr * há-zá-ban lak-nak.}

% Antiphona 2
\defsheet[tone={2-g}]{Venite}
  {<-4---rR3-2----:-4---4----4---2--4--5--5z--5-:-[------------6---5-rR3-2---:-¡-------------------1--ñ-1w-2-?,}
  {  Jöj-je--tek, * men-jünk fel az Úr he-gyé-re, mert_Sionból jön a tör-vény, és_Jeruzsálemből_az Úr i-gé-je.}

% Antiphona 3 (NZS124)
\defsheet[tone={8do-c}]{CantateDomino}
  {<-¦--------------45u-6----,-7--7---7--6---5--4-5--4-?,}
  {  Énekeljetek_az Úr--nak, * és áld-já-tok az Ő ne-vét.}

% ============================================================================
% FERIA III
% ============================================================================

% ············································································
% Officium lectionis
% ············································································

% Hymnus
% Ut in hebdomada I

% Antiphona 1
\defsheet[tone={7la-d}]{ExsurgatDeus}
  {<-4---6---7--iI7oO8-8---:-8--9--8---uU6-4t--6--7---5--4--4-?,}
  {  Fel-kél az Is-----ten * és el-szé-led-nek el-len-sé-ge-i.}

% Antiphona 2
\defsheet[tone={1la-a}]{DeusNoster}
  {<X-1-0--3--4t-5----¨-5-6---5--4--5--eE2-1-?,}
  {   A mi Is-te-nünk * a sza-ba-dí-tó Is--ten.}

% Antiphona 3
\defsheet[tone={8szo-c}]{InEcclesiis}
  {<-3---3--3--4--4tT4-3---:-5---7--tT4-3--4t-4-?,}
  {  Gyü-le-ke-ze-tek--ben * áld-já-tok az Is-tent!}

% ············································································
% Laudes
% ············································································

% Hymnus
% Ut in hebdomada I

% Antiphona 1
\defsheet[tone={6-a}]{Benedixisti}
  {<-3--4---5---3r---¨-eE2-1--:-3-4--4---3--3-?,}
  {  Meg-ál-dot-tad, * U---ram, a te föl-de-det.}

% Antiphona 2
\defsheet[tone={7do-d}]{DeNocte}
  {<-ik68o-8----:-8oO8I7-7iI7-:-4u-7i--uj5--zZ5-4--4-?,}
  {  Éj----jel  * lel----kem    te-rád vár, Is--te-nem.}


% Antiphona 3
\defsheet[tone={tirr-e}]{IlluminaDomine}
  {<-2--2--1--2-----3--1---:-0-1--3--2--1----2-2-?,}
  {  Vi-lá-go-sítsd re-ánk * a te or-cá-dat, U-runk!}

% ============================================================================
% FERIA IV
% ============================================================================

% ············································································
% Officium lectionis
% ············································································

% Hymnus
% Ut in hebdomada I

% Antiphona 1
\defsheet[tone={8szo-c}]{DeusQuiGlorificatur}
  {<-4u-5-----¨-4---5--4---3-4t-4-:-7-8----8---7--5--7---7--,-©-------------8oO8-uU67iI7-:-[-----------------------7--7---4--4?,}
  {  Is--ten, * kit di-cső-í-te-nek a szen-tek ta-ná-csá-ban, nagy_és_rette-ne---tes       mindazok_fölött,_kik_őt kö-rül-ve-szik.}

% Antiphona 2
\defsheet[tone={8szo-c}]{DeFructu}
  {<-4u--zh4-¨-6u--tT4-tT4-3-:-7--6--7---5---4---4-?,}
  {  Tes-ted * gyü-möl-csé-ből ül-te-tek szé-ked-re.}

% Antiphona 3
\defsheet[tone={8szo-c}]{Potens}
  {<-5--4---3---¨-5u-zh4--6---7--tT4-tT4-3--,-£---------3---5---7--6--6u--5--4--4-?,}
  {  Ha-tal-mas * lé-szen mag-va a   föl-dön, az_igazak nem-zet-sé-ge meg-ál-da-tik.}

% ············································································
% Laudes
% ············································································

% Hymnus
% Ut in hebdomada I

% Antiphona 1
\defsheet[tone={2-f}]{InclinaDomine}
  {<-1---3-----2--3----2-1----:-eE2-3-2---1--0----1e---2---1--1-?,}
  {  For-dítsd ho-zám, U-ram, * fü-le-det és hall-gass meg en-gem.}

% Antiphona 2
\defsheet[tone={8szo-c}]{BeatiQuiAmbulant}
  {<-5---3r-4----:-5èu-7---6u-5----4---5----4---3----4t-4-?,}
  {  Bol-do-gok, * kik tör-vé-nyed sze-rint jár-nak, U--ram.}


% Antiphona 3
\defsheet[tone={tirr-a}]{QuiaMirabila}
  {<X-5----4---5--4---:-3---4--5--6----4z-5-?,}
  {   Mert cso-dá-kat * cse-le-ke-dett az Úr.}

% ============================================================================
% FERIA V
% ============================================================================

% ············································································
% Officium lectionis
% ············································································

% Hymnus
% Ut in hebdomada I

% Antiphona 1
\defsheet[tone={1la-a}]{OmnesInimiciMei}
  {<X-1--1---3--1--0--3----4t--5---:-6---5---4---4--4--4----3t-4x-,-4----4---2---4---tT4-3-:-3----2---3---4--3---2---1í-1í-?,}
  {   El-len-sé-ge-im mind-nyá-jan * hal-lot-ták ba-jo-mat, U--ram, s_ör-ven-dez-tek raj-ta, hogy ezt cse-le-ked-ted ve-lem.}

% Antiphona 2
\defsheet[tone={6-a}]{BenedictusDominus}
  {<-3r-5----eE2-1í-:-3----4-4---3y-?,}
  {  Ál-dott az  Úr * mind-ö-rök-ké.}

% Antiphona 3
\defsheet[tone={1la-a}]{ConvertereDomine}
  {<X-1w--3----4---2-----eE2-1ed1qQ0-:-qQ0-3r---5--4---3--4--3-:-2-3--4----eE2-1--1-?,}
  {   For-dulj hoz-zánk, U---runk,   * és  légy kö-nyö-rü-le-tes a te szol-gá--id-hoz.}

% ············································································
% Laudes
% ············································································

% Hymnus
% Ut in hebdomada I

% Antiphona 1
\defsheet[tone={6-a}]{FundamentaEius}
  {<-3--4--5-3r--eE2-1-:-3-4-----4--3---3y-?,}
  {  Si-on a-lap-ja--i * a szent he-gye-ken.}

% Antiphona 2
\defsheet[tone={2-f}]{EcceDominusNoster}
  {<-1-1---3-2--0q-1---:-1--1--3---2---0---0-,-2--0--1----3---eE2-1-?,}
  {  Í-me, a mi U-runk * ha-ta-lom-mal jön el, és ma-gasz-tos kéz-zel.}


% Antiphona 3
\defsheet[tone={tirr-e}]{IustusEtSanctus}
  {<-2-2---1--wW1---:-0-1--2--2--2-?,}
  {  I-gaz és szent * a mi Is-te-nünk.}

% ============================================================================
% FERIA VI
% ============================================================================

% ············································································
% Officium lectionis
% ············································································

% Hymnus
% Ut in hebdomada I

% Antiphona 1
\defsheet[tone={1la-a}]{ZelusDomusTuae}
  {<X-1-1--1--3----1--0--3---4t-5-:-6---5-4----4----3t-4--,-4-----4--4--2--4---5--4--3-:-2---3---4---eE2--1-1-?,}
  {   A há-za-dért va-ló buz-gó-ság meg-e-mész-tett en-gem, s_gya-lá-zó-id-nak át-ka-i   rám-hul-lot-tak, U-ram. }

% Antiphona 2
\defsheet[tone={8do-c}]{ConsolantemMe}
  {<-7--uU6-5----¨-¦-------------zZ5-4-:-4--2---3--4---4--,-4--4-4t--3--5---7--uU6-5-4--5--tT4-:-2--2----3r-2--eE2-1-:-1-2---3---4-3---4t--4--4-?,}
  {  Ke-res-tem, * aki_megvígasz-tal-na  és nem ta-lál-tam: és e-pét ad-tak ne-kem e-le-de-lül   és szom-jú-sá-gom-ban e-cet-tel i-tat-tak en-gem.}

% Antiphona 3
\defsheet[tone={8szo-c}]{QuaeriteDominum}
  {<--4--4---4--3---5--7-uU6-¨-6--7--5-5--6---5--4-?,}
  {   Ke-res-sé-tek az U-rat * és él a ti lel-ke-tek.}

% ············································································
% Laudes
% ············································································

% Hymnus
% Ut in hebdomada I

% Antiphona 1
\defsheet[tone={8szo-c}]{TibiSoli}
  {<-2--3r-4x---:-3t---¦-----------5---6uc-:-5--5---5----6---5----4-4-?,}
  {  El-le-ned, * csak ellened_vét-kez-tem,  kö-nyö-rülj raj-tam, U-ram.}

% Antiphona 2
\defsheet[tone={8szo-c}]{NonNosDerelinquas}
  {<-7--6-----tT4-5---4----:-3-4t---6--5--4-?,}
  {  Ne hagyj el  min-ket, * U-runk so-ha-se.}

% Antiphona 3
\defsheet[tone={tirr-a}]{ServiteDomino}
  {<X-3----5---6--5---4--4--3---:-3--3---4t--5-?,}
  {   Szol-gál-ja-tok az Úr-nak * ví-gas-ság-gal.}

% ============================================================================
% SABBATO
% ============================================================================

% ············································································
% Officium lectionis
% ············································································

% Hymnus
% Ut in hebdomada I

% Antiphona 1
\defsheet[tone={4-a}]{ConfiteanturDomino}
  {<-1--2--3---4--2e--4--3--2----:-0--1---2---3--eE2-2-?,}
  {  Ad-ja-nak há-lát az Úr-nak, * ir-gal-mas-sá-gá--ért.}

% Antiphona 2
\defsheet[tone={8do-c}]{DeNecessitatibusNostris}
  {<-7---7---7--5--4--7---7---:-¦-------------8---8----7-7-?,}
  {  Szo-ron-ga-tá-sa-ink-ból * szabadíts_meg min-ket, U-runk.}

% Antiphona 3
\defsheet[tone={1la-a}]{QuisIntelleget}
  {<-1--qQ0-34t-tT4--:-3--3---3---4--2--3-2----1-1-?,}
  {  Ki ér--ti  meg, * ir-gal-mas-sá-ga-i-dat, U-runk?}

% ············································································
% Laudes
% ············································································

% Hymnus
% Ut in hebdomada I

% Antiphona 1
\defsheet[tone={3-b}]{InMatutinisDomineMeditabor}
  {<XB-3---4--3--4z--6c---:-zZ5-4x---3--2---3---4---eE2W1-1í-?,}
  {    Már ko-ra reg-gel, * U---ram, ró-lad gon-dol-ko----zom.}

% Antiphona 2
\defsheet[tone={2-c}]{MecumSitDomine}
  {<-7--6----7--5----tT4-6----:-6-7--5---4----6u-5-?,}
  {  Le-gyen ve-lem, U---ram, * a te böl-cses-sé-ged.}

% Antiphona 3
\defsheet[tone={4-a}]{VeritasDomini}
  {<-2--1--2-4---4t-5--:-5-4---3--4---3--2-?,}
  {  Az Úr i-gaz-sá-ga * ö-rök-ké meg-ma-rad.}
% ============================================================================
% DOMINICA
% ============================================================================

% ············································································
% Vesperas I.
% ············································································

% Hymnus
% Ut in hebdomada II

% Antiphona 1
\defsheet[tone={1re-a}]{FiatPax}
  {<-5--4t---5--4---:-2-3--4-eE2-1-?,}
  {  Le-gyen bé-ke, * a te e-rőd-ben.}

% Antiphona 2 (SzKD66)
\defsheet[tone={8szo-c}]{Sustinuit}
  {<-4--4u---5--4---3-4t--4-:-4--5--4-5--rR3-3y-,-3--5--7---7---6u-5--4--3--4t-4-?,}
  {  Re-mény-ke-dik a lel-kem az Úr i-gé-jé--ben, az én lel-kem ő--rá vá-ra-ko-zik.}

% Antiphona 3
% Ut in hebdomada II

% ············································································
% Officium lectionis
% ············································································

% Hymnus
% Ut in hebdomada II

% Antiphona 1
\defsheet[tone={8szo-c}]{DominiEstTerra}
  {<-4--4u-tT4R3-4t-4-----¨-7--8---7--5u--6--,-7-8----7--8oO8-uU67iI7-:-4--5-----5-5---7--7---4---4-?,}
  {  Az Ú--ré    a  föld, * és ben-ne min-den, a föld-ke-rek--ség       és mind, a-kik ab-ban lak-nak.}

% Antiphona 2
\defsheet[tone={6-a}]{BenediciteGentes}
  {<-3---4--5----rR3-4--1----:-3--4--4---3-?,}
  {  Áld-já-tok, nem-ze-tek, * Is-te-nün-ket.}

% Antiphona 3
\defsheet[tone={2-d}]{Introibo}
  {<-5--7i-8---:-7-8--9--8---7---8-8--,-5i-uU6-5---4--tT4-:-3r----5----7--5---5-?,}
  {  Be-lé-pek * a te há-zad-ba, U-ram, és i---mád-ko-zom   szent temp-lo-mod-ban.}

% ············································································
% Laudes
% ············································································

% Antiphona 1 (D39)
\defsheet[tone={7la-a}]{OmnisSpiritus}
  {<ô-1---3---5z-5--:-4--2----3--2--1-1í-?,}
  {   Min-den é--lő * di-csér-je az U-rat.}

% ············································································
% Hora Media
% ············································································

% Antiphona 3 (vö. GH363)
\defsheet[tone={7do-a}]{Vovete}
  {<ô-5--5---3---5--6u-zZ5-5---:-5--3--4t-wW1-1-,-4--4-4--2r--3-:-2--1--2rR3-1---1-?,}
  {   Te-gye-tek fo-ga-dal-mat * és ad-já-tok meg ti U-ra-tok-nak Is-te-ne---tek-nek.}

% ············································································
% Vesperas
% ············································································

% Antiphona 2 (NZS14)
\defsheet[tone={8do-c}]{InMandatisEius}
  {<-7---6---7--5x---:-4--3--5--5--6---5----4---4--4x-?,}
  {  Bol-dog em-ber, * ki az Úr pa-ran-csát ked-ve-li.}

% ============================================================================
% FERIA II
% ============================================================================

% ············································································
% Laudes
% ············································································

% Antiphona 1 (D60)
\defsheet[tone={6-a}]{DomineRefugium}
  {<X-3-4----5--3---3---:-3-1--3--4--4--3?,}
  {   U-ram, te let-tél * a mi me-ne-dé-künk.}

% ············································································
% Hora Media
% ············································································

% Antiphona 2 (Nok145)
\defsheet[tone={8do-c}]{IusteIudicate}
  {<-7-7--7---5-4---45u-7----:-7--7--5--4---5--4-4x-?,}
  {  I-ga-zul í-tél-je--tek, * ti em-be-rek fi-a-i!}

% Antiphona 3 (NZS256)
\defsheet[tone={2-g}]{Clamavi}
  {<ô-4--4--rR3E2-2y---:-2--1---2----4--4----1r-2y-?,}
  {   Ki-ál-tot---tam, * és meg-hall-ga-tott en-gem.}

% ············································································
% Vesperas
% ············································································

% Antiphona 1 (D184)
\defsheet[tone={3-c}]{QuoniamInAeternum}
  {<-7-uU6-tT4-5--4x-:-4--4---4---rR3-4t-rR3-2y-2y-?,}
  {  Ö-rök-ké--va-ló * ir-gal-mas-sá--ga az  Úr-nak.}

% ============================================================================
% FERIA III
% ============================================================================

% ············································································
% Officium lectionis
% ············································································

% Antiphona 1 (NZS150)
\defsheet[tone={1re-a}]{ClamorMeus}
  {<-5--4--5--rR3-:-2---3---rR3-4t---eE2-1í-1í-?,}
  {  Ki-ál-tá-som * jus-son e---léd, ó   Is-ten!}

% Antiphona 3 (SzKD87)
\defsheet[tone={8szo-c}]{IustiConfitebuntur}
  {<-£---------5--4---4--3---:-5-7--7--uU6-tT4--5-4-,-£----------------eE2-1-:-3---3---3-4t---4x--4x-?,}
  {  Az_igazak há-lát ad-nak * a te ne-ved-nek, U-ram s_az_egyenes_szí-vű--ek  szí-ned e-lőtt lak-nak.}

% ············································································
% Vesperas
% ············································································

% Antiphona 1 (D184)
\defsheet[tone={8szo-c}]{HymnumCantate}
  {<-4-4--4---5--4---tT4-3y----:-3--5--7--7--7-6u--5--4x-?,}
  {  É-ne-kel-je-tek him-nuszt * Si-on da-la-i-ból ne-künk.}

% Antiphona 2 (D193)
\defsheet[tone={5-c}]{InConspectu}
  {<-5--5---4---3---5---7--7iI7U6-5----,-5----5--5--4---5--4----3-3y-?,}
  {  An-gya-lok-nak szí-ne e------lőtt * zsol-tá-ro-zok né-ked, U-ram.}

% ============================================================================
% FERIA IV
% ============================================================================

% ············································································
% Officium lectionis
% ············································································

% Antiphona 1 (Nok213)
\defsheet[tone={1re-a}]{Benedic}
  {<-5---4----¨-2e-4---3----2--1-1-?,}
  {  Áld-jad, * én lel-kem, az U-rat!}

% ============================================================================
% SABBATO
% ============================================================================

% ············································································
% Hora Media
% ············································································

% Antiphona 1
\defsheet[tone={1la-a}]{FiatManusTuaDomine}
  {<X-0-1---1tu-5----5--rR3--4t-5----:-3----2---3---4--3---2---1--1--,-3----2-1--1--2---3---4tT4-3-:-2e-4----eE2-1-?,}
  {   A-zon le--gyen ke-zed, U--ram, * hogy meg-sza-ba-dít-son en-gem; mert a te pa-ran-csa-i----dat vá-lasz-tot-tam.}

% Antiphona 2
\defsheet[tone={6-a}]{Eructavit}
  {<X-3--4--5---3---:-4--eE2-1--3---3---4---4-3-?,}
  {   Az én szí-vem * ün-ne--pi szó-zat-tól á-rad.}

% Antiphona 3
\defsheet[tone={4-a}]{VidiIerusalemNovam}
  {<-2---2---2--1--3--2--qQ0-1w-2----:-1----2-3---4t-4---4---5-4--3--4----4----rR3-2-?,}
  {  Lát-tam az új Je-ru-zsá-le-met, * mint a fér-jé-nek föl-é-ke-sí-tett meny-as--szonyt.}

\makeatother
% ============================================================================
% IUNIUS
% ============================================================================

% ············································································
% Die 24: In Nativitate Sanctis Ioannis Baptistae
% ············································································

% --- Laudes ---

% Antiphona 3
\defsheet[tone={7la-a}]{TuPuer}
  {<ô-1---1t---5----:-5-5--6---7---zZ5-5-:-5---3--4--5--wW1-1---,-£--------------2r---3---:-2--1--2---2--2--4--3-1--1-?,}
  {   Te, gyer-mek, * a Ma-gas-ság-be--li  pró-fé-tá-ja le--szel, az_Úr_orcája_e-lőtt mégy, el-ké-szí-te-ni az ő út-ját.}

% ············································································
% Die 27: S. Ladislai
% ············································································

% --- Officium Lectionis ---

% Antiphona 1
\defsheet[tone={1la-a}]{HicIniquisNonConsesnsit}
  {<-0--1---1tu-5---¨-tT4-3---4tT4-eE2W1-,-1--0--34t-5---44-3--44--1---,-44t-4-eE2-1---2----0--1---1-,-5--4---5---7i---78o-uU6Z5-4t--5-,-5-4---55--3-----4-3---4---5---,-5-4--eE2-1----2ed1-0q--1---1-?,}
  {  Go-no-szok-kal * nem cim-bo---rált,   út-ja-i---kon so-ha-sem járt. Úr  i-gá--ját váll-ra vet-te; Is-ten ked-vét: azt ke----res-te. A-mit mun-kált, a-mit gon-dolt, é-le-te  mind ta---ní--tás volt.}

% Antiphona 2
\defsheet[tone={2-f}]{RexAChristoConstitutusEst}
  {<-0--1--12ed1-1---¨-qQ0---ð---0q--1-,-0--1--23r-rR3E2W1-2--0--1---1-,-3-4--4----3--4--3---4t--5-,-5-4-----eE2-1----qQ0-ð---0q--1-?,}
  {  Őt ki-rály--lyá * Krisz-tus tet-te, tá-mo-gat-ta,     se-gí-tet-te, e-ré-nyek-be öl-töz-tet-te, a szent he--gyen vé--del-mez-te.}

%% Antiphona 3
%\defsheet[tone={4-a}]{LaudateTeRexGloriae}
%  {<-3---4z--3--rR3--1w-3-3---,-3r--5?,}
%  {  Áld-jon té-ged, ég U-ra, * Egy-há-zunk szent szó-za-ta, ün-ne-pel-jen zeng-ve ma é-des é-nek dal-la-ma.}

% ============================================================================
% IULIUS
% ============================================================================

% ············································································
% Die 3: S. Thomae
% ············································································

% --- Laudes ---

% Antiphona 1
\defsheet[tone={8szo-b}]{StetitIesus}
  {<-4---3--5----7--7---:-4-6--7---5--4----tT4-3---,-4---4---4---3--5----7--7--7--:--4--6u-tT4-3---4t-5--4x-4x-?,}
  {  Meg-ál-lott Jé-zus * a ta-nít-vá-nyok kö--zött: "Bé-kes-ség le-gyen ve-le-tek!" Al-le-lu--ja, al-le-lu-ja.}

% Antiphona 2
\defsheet[tone={1la-a}]{ThomasQuiDiciturDydimus}
  {<X-1--1----1--3-1---0---3---4t-5----:-5---5----5u-5--:-5-4--3---2e-4----eE2-1--,-1------1---3--1--0-3--4---3t-5---:-5---4---4t-eE2-qQ0--1--eE2-1í-1í-?,}
  {   Ta-más, ki I-ker-nek mon-da-tik, * nem volt ve-lük, a-mi-kor el-jött Jé--zus. S_mond-ták ne-ki a ta-nít-vá-nyok: lát-tuk az U---rat, al-le--lu-ja.}

% Antiphona 3
\defsheet[tone={4-a}]{MitteManumTuam}
  {<ô-1-------2--4-4--5z-5----:-4-5------6-5---4---5--5---,-5--2--5----4--3---2--1---wW1-:-0--1---2-–2e--4--3--2--2-?,}
  {   Nyújtsd ki a ke-ze-det, * é-rintsd a sze-gek he-lyét, és ne légy im-már hi-tet-len,  ha-nem hí-vő, al-le-lu-ja.}

% Benedictus
\defsheet[tone={8szo-c}]{QuiaVidistiMe}
  {<-4--4---3--5---7--7----:-4--6u---5--4---tT4-3--,-4---4--4----3-5---7---7---7--:-4--6u---5----6--5z-4x-4x-?,}
  {  Mi-vel lá-tál en-gem, * Ta-más, te-hát hit-tél, bol-do-gok, a-kik nem lát-nak, és hisz-nek, al-le-lu-ja.}

% --- Vesperas ---

% Magnificat
\defsheet[tone={8szo-c}]{Misi}
  {<-4--4----4u--5tT4-¨-3-4t-5--4-:-7-8---uj5-7--7---7-:-7--8--7--8oO8-uU67iI77-:-5--5-7--7--4x-4x-:-4-----eE2-3--:-1--2-3----4--4--rR3--4t-5--4x-4x-?,}
  {  Ki-nyúj-tot-tam  * a ke-ze-met a szö-gek he-lyé-be, és uj-ja-i----mat        az ő ol-da-lá-ba;  s_mon-dot-tam: én U-ram, Is-te-nem, al-le-lu-ja.}

% ············································································
% Die 11: S. Benedicti
% ············································································

% --- Officium Lectionis ---

% Hymnus
\defsheet{InterAeternas}
  {<X-5-4--5--4----1í--:-4--4---3---4--5x--5x-,-5---7--8----7---5x-:-5--4-3---4--5x-5x-,-5---4--5-------4---1í-:-4--4--3---4--tT4-3y-,-3--3--2--1í-1í-?,-1wW1-0\\q\\-.}
  {   É-gi fé-nyek közt, ki-ket áld-va hir-det  har-ci győz-tes-ként di-a-dal-mi é--nek, ott ra-gyogsz, bol-dog  Be-ne-dek su-gár-zó   ér-de-me-id-del.  Á----men.}

% --- Laudes ---

% Hymnus
\defsheet{LegiferPrudens}
  {<X-5-----4--5--4-----1í--:-4--4--3----4--5x-5x-,--5-----7--8----7--5x-:-5-4--3---4--5x-----5x--,-5---4--5--4---1í-:-4-4--3----4--tT4--3y-,-3--------3--2---1í--1í-?,-1wW1-0\\q\\-.}
  {   Bölcs ta-ní-tónk, Szent Be-ne-dek, vi-lá-gíts, Krisz-tu-sunk fé-nyét a vi-lág-ba hintsd szét, hí-ven is-mer-nünk e hi-tet, se-gíts meg, s_töltsd be szí-vün-ket!  Á----men.}

% Benedictus
\defsheet[tone={8do-c}]{FuitVir}
  {<-7-5u-7--¨-7iI7-uU6u-uU6-5--6u--5--4---,-rR3-4--3--5--uU67iI7-:-5u-6---5---4---3r-4-?,}
  {  É-le-te * tisz-te---let-re-mél-tó volt, ne--ve Be-ne-dek:      ke-gye-lem-mel ál-dott.}

% ············································································
% Die 15: S. Bonaventurae
% ············································································

% Benedictus
\defsheet[tone={2-f}]{SpirituSapientiae}
  {<-1-1äe--1---¨-1---1----qQ0-1--1--2--3---2---0--0-:-2--4---3---1--3r-eE2-1-?,}
  {  A hasz-nos * böl-cses-ség és ér-te-lem lel-ké-vel be-töl-töt-te őt az  Úr.}

% ============================================================================
% AUGUSTUS
% ============================================================================

% ············································································
% Die 8: S. Dominici
% ············································································

% --- Laudes ---

% Hymnus
\defsheet{NovusAthletaDomini}
  {<ô-2--2---2--1---2---4---3--2í-:-2--3---4----5--4---5---5--6x-,-6--6--5---6---4--5--rR3-2í-:-2--1----2-4--5---2---1--2í-?,-1e4R3E2W1-2í-.}
  {   Di-cső ne-vet mél-tán vi-selt Do-mon-kos, Is-ten baj-no-ka,  az Úr-hoz fűz-te őt ne--ve,  és lett a jó hír har-co-sa.   Á---------men.}

% Benedictus
\defsheet[tone={1re-a}]{Beatus}
  {<-0q--1tz-5--5z-4tT4---:-4-5--6--5--rR3-3rR3-1---:-4--5--4--3---eE2--1---2---0-,-0--1--3-:-3-----3--4t--wW1-3--2---1--1-?,}
  {  Bol-dog az a  szent, * a-ki az Úr-ban-bí---zott, az Úr pa-ran-csát hir-det-te, és az ő   szent he-gyé-re  ál-lít-ta-tott.}

% ············································································
% Die 24: S. Bartholomaei
% ············································································

% --- Laudes ---

% Hymnus
\defsheet{RelucensInterPrincipes}
  {<-5--5--5--5--1---3--2--1í--:-5--5--6---4----5-----7----6--5x-,-5-----6---7--8----7----5---7--zZ5T4x-:-5---5--6---4--2---4---5--5x-?,-5zZ5-4\\t\\-.}
  {  Is-ten nagy há-za né-pe közt ki-rá-lyi fény-ben tün-dö-kölsz, Szent Ber-ta-lan, fi-gyelj re-ánk, ké-rünk, míg így ma-gasz-ta-lunk:  Á----men.}

% ============================================================================
% BEATEAE MARIAE VIRGINIS
% ============================================================================

% ············································································
% Officium Lectionis
% ············································································

% Antiphona 1
\defsheet[tone={1la-a}]{QuisAscendet}
  {<-0--1--1tu-5---:-[-----------tT45uj5T4-:-5--4--3---3r--3-:-2e-4-eE2---1--1---,-eE2-1--2---3---4tT4-3-:-2e-4----eE2-1---1-?,-------3--3----3---2---3---4--eE2-1--1-?,}
  {  Ki me-het fel * az_Úr_hegyé-re,         és ki áll-hat meg az ő szent he-lyén? Az  ár-tat-lan ke---zű  és tisz-ta  szí-vű. <T.P.> és tisz-ta  szí-vű, al-le--lu-ja.}

% Antiphona 2
\defsheet[tone={6-a}]{Sanctificavit}
  {<-3---4----5---4--¨-4t-rR3-3--4-5---4--3-?,--------4--4t-3---3-?,}
  {  Meg-szen-tel-te * az Úr  az ő haj-lé-kát. <T.P.> Al-le-lu--ja.}

% Antiphona 3
\defsheet[tone={5-c}]{Gloriosa}
  {<-¦----------7iI7U6u-:-5---7--6---4t-rR3-4t---4--3--3-?,------4--4t-3--3-?,}
  {  Dicsőséges-nek       mon-da-nak té-ged Szűz Má-ri-a. <T.P.> Al-le-lu-ja.}

% ============================================================================
% APOSTOLORUM
% ============================================================================

% ············································································
% Invitatorium
% ············································································

\defsheet{RegemApostolorum}
  {<-0-0ãwe-2--:-2--2-2---2--3r--rf2-3-1wW1Q0\\-,-0wår-4--4t---5-tT4R3rR3E2-2-.}
  {  U-run--kat, az a-pos-to-lok Ki-rá-lyát,    * jöj--je-tek, i-mád--------juk!}

% ············································································
% Officium Lectionis
% ············································································

% Hymnus
\defsheet{OSempiternaeGloriae}
  {<ô-4--4-4--4--4---4--3---2í-:-4-4---4--4--1-----2----1--0í-,-2--2---3--4----3--2----1--2eE2í-:-2---2---1--0---1----2----1--1í-?,--1wW1-0\\q\\-.}
  {   Az é-gi ud-var dí-sze-i,   ö-rök Ki-rá-lyunk nagy-ja-i,   ki-ket ne-künk ne-velt az Úr,     s_a-pos-to-lok-ként ránk ha-gyott, Á----men.}

% Antiphona 1
\defsheet[tone={2-f}]{InOmnemTerram}
  {<-1e--2---3r---4--:-4--3---1--3-4---1--1--,-1-3---2--3---4---4---3--3--1e-2e-4-1--1-?,--------0--0w-1--1-?,}
  {  Min-den föld-re * el-jut az ő zen-gé-sük, a föl-ke-rek-ség szé-lé-ig az ő  i-gé-jük. <T.P.> Al-le-lu-ja.}

% Antiphona 2
\defsheet[tone={8szo-c}]{AnnuntiaveruntOpera}
  {<-3---5u-zh4-:-6--7---tT4-tT4-3-,-4--3--5-7--6--4--6u--5--4-?,--------6u--5--4t---6--5z-4--4-?,}
  {  Hir-de-tik * Is-ten te--te--it  és az ő te-te-it meg-ér-tik. <T.P.> meg-ér-tik, al-le-lu-ja.}

% Antiphona 3
\defsheet[tone={8szo-c}]{AnnuntiaveruntIustitiam}
  {<-rR3d1-34t-4---:-4--6i-7-5---77-6--,-88-7---8oO8-uU67i7-:-5---5--4--5--5èuj5-zZ5-4-?,--------5--5èuj5-zZ5-4t---6--5--4--4-?,}
  {  Hir---de--tik * az ő  i-gaz-sá-gát, és min-den  nép      meg-lá-ta di-cső---sé--gét. <T.P.> di-cső---sé--gét, al-le-lu-ja.}

% ············································································
% Laudes
% ············································································

% Antiphona 1
\defsheet[tone={1la-a}]{HocEstPraeceptumMeum}
  {<-1--0--1--3--wW1-qQ0---:-3---4---5--4---tT4-3---:-3-3r---33-2---3---4--eE2-1--1--1-?,}
  {  Ez az én pa-ran-csom: * sze-res-sé-tek egy-mást, a-mint én sze-ret-te-lek ti-te-ket.}

% Antiphona 2
\defsheet[tone={1la-a}]{MaioremCaritatem}
  {<X-1--0-----3---4--3t-5--:-6---5--4---4-----3t-4--,-4----4--4-2--4---tT4-3-:-2--3-4--eE2-1--0q-1-?,}
  {   Na-gyobb sze-re-te-te * sen-ki-nek sincs an-nál, mint ki é-le-tét ad--ja  az ő ba-rá--ta-i--ért.}

% Antiphona 3
\defsheet[tone={1la-a}]{VosAmiciMeiEstis}
  {<X-1--1--3--1--0--34t--5----:-4--6---5--4---4---3t---4--,-4-4---4--4---2----4---tT4-3--:-2----3--4--eE2-1-----1-?,}
  {   Ti ba-rá-ta-im vagy-tok, * ha meg-te-szi-tek mind-azt, a-mit pa-ran-csol-tam nek-tek, mond-ja az Úr  Krisz-tus.}

% Responsorium breve
\defsheet{ConstituesEosPrincipes}
  {<X-¢------------4--3----4t-tT4-,-5--4-3----rR3-3-?,-5--5--5----5z-5--:-5---4---4-,-5---4---3--4t--tT4-?,-[-------------5z--5-:-5-5--4-4-:-5--5-4-----3--4t--tT4-.}
  {   Fejedelmekké te-szed ő--ket * Az e-gész föl-dön. Ne-ve-det, U--ram, min-den-kor hir-det-ni fog-ják.   Dicsőség_az_A-tyá-nak a Fi-ú-nak és a Szent-lé-lek-nek.}

% Benedictus
\defsheet[tone={2-c}]{DumSteteris}
  {<-5-5--tG2-4t-5z-5----¨-5--5-5---6uU6-55T4-5-7----5z-6---;-6--8---8---9--iI7-uU6-5--5-:-7---uU6-tT4-5--6u-5----5-,-5----7---7-6--tT4-55--4-:-4-5---77-6u-8o--uU6-5----4t-5-?,}
  {  A-mi-kor ki-rá-lyok * és e-löl-já---rók  e-lőtt ál-ltok, ne ter-vez-zé-tek e---lő-re, mit fog-tok vá-la-szol-ni: mert meg-a-da-tik nék-tek a-zon ó-rá-ban, mit mond-ja-tok.}

% ············································································
% Vesperas II.
% ············································································

% Responsorium breve
\defsheet{AnnuntiateInterGentes}
  {<X-¢-----------------4--3---4t-tT4--,-5--5--4--3---rR3-3-?,-[-----------5z-4-,-4---5---rR3-4t-tT4-?,-[-------------5z--5-:-5-5--4-4-:-5--5-4-----3--4t--tT4-.}
  {   Hirdessétek_a_nem-ze-tek kö-zött * Az Úr di-cső-sé--gét. És_csodatet-te-it  min-den nép kö-zött!  Dicsőség_az_A-tyá-nak a Fi-ú-nak és a Szent-lé-lek-nek.}

% ============================================================================
% UNIUS MARTYRIS
% ============================================================================

% ············································································
% Officium Lectionis
% ············································································

% Antiphona 1
\defsheet[tone={1la-a}]{SiQuisPreseveraverit}
  {<-qQ0-3--rR3-4t--5---:-6----5--4---3t-4---:-4-2---4---tT4-3-:-2--3r--eE2-1-?,--------0--0q-1--1-?,}
  {  A---ki a---zon-ban * mind-vé-gig ki-tart, é-let-ben ma--rad és nem hal meg. <T.P.> Al-le-lu-ja.}

% Antiphona 2
\defsheet[tone={8szo-g}]{Existimo}
  {<ô-1r-2---wW1Q0-1w-2---1---:-5----4--2--4--4-:-[----------------------------6---5---rR34tT4-:-¡------------------4--4---1---1-?,---------1--1--]ñ-ð-:-0--1w-1--1-?,}
  {   En-nek az    é--let-nek * szen-ve-dé-se-i   nem_mérhetők_az_eljövendő_di-cső-ség-hez,      amely_majd_megnyil-vá-nul raj-tunk. <T.P.> Al-le-lu--ja, al-le-lu-ja.}

% Antiphona 3
\defsheet[tone={7do-a}]{TamquamAurumInFornace}
  {<ô-5----5-5--3----5-6--5--5----:-[-------------------3---4--5--wW1-1-:-2--4-4--4--4--2r--3-:--¡--------------4----3-1---1-?,}
  {   Mint a-ra-nyat a ko-hó-ban, * megvizsgálta_válasz-tot-ta-it az  Úr, és é-gő-ál-do-zat-ként elfogadta_őket mind-ö-rök-ké. }

% ············································································
% Laudes
% ············································································

% Antiphona 1
\defsheet[tone={1re-a}]{QuiaMeliorEst}
  {<-1--qQ0-3r-3---4t--:-5u---5--4--5----4--rR3-4t-,-1----0--1----eE2-4---3--1---0q-1-?,}
  {  Mi-vel ir-gal-mad * több-et ér mint az é---let  hadd ma-gasz-tal-jon té-ged aj-kam.}

% Antiphona 2
\defsheet[tone={tper-g}]{MartyresDomini}
  {<-1-0----1---3--3r-4--:-5---5--4---3--2-1---3----4t-4---4--?,}
  {  U-runk vér-ta-nú-i, * áld-já-tok az U-rat mind-ö--rök-ké.}

% Antiphona 3
\defsheet[tone={1re-a}]{QuiVicerit}
  {<-0q-1tz--5-----:-56u-5--XtT45zZ5T4-3r-4t-1-:--2----3--4t--4tT4R3-2e---3rR3E2-0q-1-?,}
  {  A  győz-test, * osz-lo-pom--------má te-szem temp-lo-mom-ban    mond-ja     az Úr.}

% Benedictus
\defsheet[tone={3-c}]{QuiOdit}
  {<-4-4--456u-5--5--:-7-uU6-5z--5-4---5-uj5-6u-6--,-6u--8--7--uU6-tT4-:-tT4R3-4---rR3-2---2-?,}
  {  A-ki gyű--lö-li * é-le--tét e-zen a vi--lá-gon, meg-őr-zi azt az    ö-----rök é---let-re.}

% ============================================================================
% SANCTORUM VIRORUM
% ============================================================================

% ············································································
% Invitatorium
% ············································································

\defsheet{MirabilemInSanctisSuisDominum}
  {<-0ã2e-2---:-2-2--23r-rR3-4---4t--rR3---1w-2e-2----,-1t--5z-5----5-tT4R3rR3E2-2-?,}
  {  Is---tent, a-ki cso-dá--la--tos szent-je-i--ben, * jöj-je-tek, i-mád--------juk!}

% ············································································
% Officium Lectionis
% ············································································

% Antiphona 1
\defsheet[tone={8szo-c}]{Vitam}
  {<-3-4t-4---¨-4t---7--5--:-7--7---7--tT4-5z-5---tT4-4--,-4--7---7i-6---7--5----4-4t--rR3-3---3---4t--uj5-4-,-4--eE2-1--3--4--5---uU6-5z--7---5--4--4-?,}
  {  É-le-tet * kért tő-led, és meg-ad-tad ne-ki, U---ram, di-cső-sé-get és nagy é-kes-sé--get tet-tél rá--ja, fe-jé--re ko-ro-nát rak-tál drá-ga-kő-ből.}

% Antiphona 2
\defsheet[tone={7do-d}]{IustumDeduxit}
  {<-8--8-6--8---9-8---8---8-7i-8---:-8--6---7--uU6-4-,-¦-------------5u-6-:-5--4---5u-6---4-?,}
  {  Az i-ga-zat e-gye-nes u-ta-kon * ve-zet-te az Úr, és_megmutatta ne-ki  Is-ten or-szá-gát.}

% Antiphona 3
\defsheet[tone={8szo-c}]{IustusUtPalma}
  {<-[-------------4-tT4-3--:-3tè7i-uU6-7iiI7-,-7----7-7--6--5---uU6-tT4-tT4-:-3r---tT4-4-?,}
  {  Az_igaz,_mint a pál-ma * vi----rág-zik,    mint a Li-ba-non céd-ru--sa    nagy-ra  nő.}

% ············································································
% Laudes
% ············································································

% Antiphona 1
\defsheet[tone={1la-a}]{Iucunditatem}
  {<-2--0---1tu-5---:-[----------tT45uj5T4-:-3r--5---eE2-1-----,-eE2-1-1---1--2---3--4tT4-3---:-3-3---3rR3-0-:-2-3--4--eE2-1-?,}
  {  Vi-gas-sá--got * és_örvende-zést        hal-moz rá--ja/juk, és  ö-rök ne-vet ad ne---ki(k) ö-rök-sé---gül a mi Is-te--nünk.}

% Antiphona 2
\defsheet[tone={tper-g}]{ServiDomini}
  {<-0--1--3----3r-4-:--5--4----5--4---3--2-1---3----4t-4---4-?,}
  {  Az Úr szol-gá-i, * di-csér-jé-tek az U-rat mind-ö--rök-ké.}

% Antiphona 3
\defsheet[tone={4-a}]{Exsultabunt}
  {<-2--2----2--1äeE2s0-:-2-2e---4---3-2--2e--4---4--,-4--4--3--eE2-1--1---3----1--3--4---2-?,}
  {  Uj-jong-ja-nak     * a szen-tek a di-cső-ség-ben, és vi-ga-doz-za-nak nyug-vó-he-lyü-kön!}

% Responsorium breve
\defsheet{LexDeiEius}
  {<X-33r-3-:-4---4t-tT4-,-5--4-3---rR3-3-?,-5z--4-,-5--5--5---4---3---4t-tT4-?,-[-------------5z--5-:-5-5--4-4-:-5--5-4-----3--4t--tT4-.}
  {   Is--ten tör-vé-nye * Él a szí-vé--ben. Nem is  té-to-váz-nak lép-te-i. Dicsőség_az_A-tyá-nak a Fi-ú-nak és a Szent-lé-lek-nek.}

% ============================================================================
% SANCTORUM MULIERUM
% ============================================================================

% ············································································
% Invitatorium
% ············································································

% Ut supra, in Communi sanctorum virorum

% ············································································
% Officium Lectionis
% ············································································

% Hynus (Pro una sancta muliere)
\defsheet{HaeFeminaeLaudabilis}
  {<-1-----3---eE2---4---ed1--2-4--3\\-:-3--1--2--3---2e--1---2--1í-,-1----3----eE2-4-----tT4--2-4--3\\-:-3---1--2---3---2e---1--2--1í-?,-2eE2-1\\q\\-.}
  {  Szent asz-szony-ról szól é-ne-künk, ki ér-de-mek-ben gaz-da-gon, mert tisz-ta, szent volt é-le-te,   már an-gya-lok közt ün-ne-pel.  Á----men.}

% ············································································
% Laudes
% ············································································

% Hymnus (Pro una sancta muliere)
\defsheet{NobilemChristiFamulam}
  {<X-5-----4---5---4-1í-:-4--4---3----4--5x-5x-,-5----7--8---7-5x-:-5----4--3-4--5x--5x-,-5-----4-5---4---1í----:-4--4--3--4-tT4-3y-,-3--3--2--1í-1í-?,-1wW1-0\\q\\-.}
  {   Krisz-tus-nak é-kes, ne-mes szol-gá-ló-ját  zeng-ze-tes é-nek  hang-ja-i da-lol-ják, szent e-rős asz-szonyt, ki-re az Í-rás ád   is-te-ni ál-dást. Á----men.}

% Responsorium breve (Pro una sancta, Extra tempore paschali)
\defsheet{AdiuvabitEam}
  {<X-3--3---3--4---3---4t-5--:-4--5--rR3-4--3---3---,-[--------------5z--4-:-5----5--5---4--3---4t--tT4-?,-[-------------5z--5-:-5-5--4-4-:-5--5-4-----3--4t--tT4-.}
  {   Is-ten ol-tal-maz-ta őt * És re-á   te-kin-tett. Nem_fog_megren-dül-ni, mert Is-ten la-kik ben-ne.    Dicsőség_az_A-tyá-nak a Fi-ú-nak és a Szent-lé-lek-nek.}

% Benedictus
\defsheet[tone={8szo-c}]{SimileEst}
  {<-4--5z--6--¨-7-7----5----6--5---4-,-4-4i-8---8--6--7--6---6--:-5-4--5-6---7----5----6--5--4---,-rF1-45z--6u--8--7--tT4-4--,-4--4i-8--7---5--uU6-:-4--5---5èuj5-zZ5-4-?,}
  {  Ha-son-ló * a meny-nyek or-szá-ga  a ke-res-ke-dő em-ber-hez, a-ki i-gaz-gyön-gyöt ke-re-sett; ta--lált egy ér-té-ke--set, el-ad-ta min-de-nét,  és meg-vet---te  azt.}

\defsheet{DeusInAduitorium}
  {<ô---£---------------------3r--4-,------£----------------3r-4-?,-£------------------------------------------3r--4--,-£---------------------------------------3r-4--,-4--5--rR3-3-.}
  {<V.> Istenem,_jöjj_segítsé-gem-re! <R.> Uram,_segíts_meg en-gem! Dicsőség_az_Atyának,_a_Fiúnak_és_a_Szentlé-lek-nek, miképpen_kezdetben,_most_és_mindörökké. Á--men. Al-le-lu--ja.}

\defsheet{BenedicamusDomino}
  {<----£---------------2--1w-2--,------£-----------2----1w-2-.}
  {<V.> Mondjunk_áldást az Úr-nak! <R.> Istennek_le-gyen há-la.}

% ============================================================================
% TONI INVITATORII
% ============================================================================

\defsheet{InvitatoriumMi}
  {<-5---4--5---:-[------------4--6--5--:-4--[----------------------4--4z-5--5---,-5---4----[---------4--3--4t--5--:-5--rR3--23r-eE2-1-:-3---3---4-5---rR3E2-2----?,-----5----4--[----------4----6--5--:-4--[--------------------4-----6--5---,-4-[-------------------------4---6---5--:-4--[----------------4---6-----5-,-5-4--3-4t--5--:-2e-4-eE2-1--:-1-¢-----------------------4--5--rR3E2-2e--?,-5---4--[------------------------4----6-5---5-:-4----[-----------------------4----6-5---,-4----5-5-4--3--4t-5---:-5--rR3-2e-4--eE2-1-:-1--3--4--5---rR3E2-2----?,----5---4----[------------------4-6---5--5-,--4--[-----------4----6---5--5-:-t4---[---------------------4---3-4----5--5--:-5-4---[--------4---3--2---3---4---eE2-1-:-1--¢--------------------------4---5---4---rR3E2-2e--?,-5----4---[---------------------4-6---5--5--:-5---4----5----[---------------4--6---5-,-5---4--5---5---5---4-3--4t-5--:-5-4----[----------4--3----2e--4--eE2-1--:-3---1---¢---------------4---5---rR3E2-2-----?,----5--4---[-----------------------------------4---6--5---5--:-3--4---[----------------rR3-2----3--4---eE2-1-:-1--¢-----------3-4---5---rR3E2-2e-?,-.}
  {  Jöj-je-tek,  örvendezzünk az Úr-nak, és ujjongjunk_üdvünk_szik-lá-ja e--lőtt! Lép-jünk színe_elé há-la-dal-lal, ma-gasz-tal-juk őt  han-gos é-nek-szó---val! ANT. Mert az Úr_valóban nagy Is-ten, az összes_isteneknél_di-csőbb ki-rály. A föld_mélységei_az_ő_kezé-ben van-nak,  és övé_a_hegyek_min-den or----ma. Ö-vé a ten-ger, az ő mű--ve,  a szárazföld_is_az_ő_keze al-ko-tá----sa. *  Jöj-je-tek,_boruljunk_le,_hódol-junk e-lőt-te, hull-junk_térdre_Urunk,_Alko-tónk e-lőtt, mert ő a mi Is-te-nünk, mi meg há-za né--pe  és ke-zé-nek nyá---ja. ANT. Bár-csak hallgatnátok_ma_az ő sza-vá-ra: „Ne legyetek_ke-mény szí-vű-ek, mint Meribánál,_Massza_nap-ján a pusz-tá-ban, a-hol atyáitok meg-kí-sér-tet-tek en-gem, és próbára_tettek,_bár_látták szá-mos cso-dá----mat. * Negy-ven évig_terhemre_volt_ez a nem-ze-dék, azt mond-tam: Ez_a_nép_tévely-gő szí-vű. Nem is-mer-ték meg u-ta-i--mat, e-zért haragomban ki-mond-tam es-kü--met: Nyu-gal-mam_helyére_nem jut-hat-nak   el!” ANT. Di-cső-ség_az_Atyának_és_Fiúnak_és_Szentlé-lek Is-ten-nek, mi-kép-pen_kezdetben_va-la, most és min-den-kor és mindörökkön ö-rök-ké! Á-----men.     * ANT.}

\defsheet{InvitatoriumFa}
  {<X-5---6--¦-----------------5--6--5--:-3r-[----------------------6--5--4-3---,-5---6----¦------------5--6---5--:-3--4----[----------------6-5---4---3------?,-5z---¦-------------5----6--5--:-3r-[--------------5---6--5-----4--3---,-5z-¦------------------------5---6---5--:-3r-[------------6---5---4--3-,-5-6--¦------------5--6-5-:-3r-[-----------------------4--5--6--5--?,-5---6--¦---------------------7---5----6-6---5-:-3----4----[----------------6--5----4-3---,-5z---7-7-5--6--6--5---:-3--4---[----------------6-5---4---3---?,---5---6----¦------------------5-6---6--5-:--3--4--5---5---6--5----4---3--3-,-5z---¦---------------------------------5--:-¦-------------------------5---6--5--:-3r-[--------------------------4---5---6---5--5---?,-5----6---¦---------------------5-6---6--5--:-3r--[--------------------6----5--4---3-,-5z--¦-----------------5--:-¦-----------------------5---6--6--5--:-3---4---[---------------6---5---4---3----?,---5--6---¦--------------5--6--5-5-:-3r-5-----5--6---5--4---3--,-5--6---¦------------------------5--6---6---5-:-3r-[-------------4---5---6-5----?,-.}
  {   Jöj-je-tek,_örvendezzünk az Úr-nak, és ujjongjunk_üdvünk_szik-lá-ja e-lőtt. Lép-jünk színe_elé_há-la-dal-lal, ma-gasz-taljuk_őt_hangos é-nek-szó-val! ANT. Mert az_Úr_valóban nagy Is-ten, az összes_istenek-nél di-csőbb ki-rály. A  föld_mélységei_az_ő_kezé-ben van-nak, és övé_a_hegyek min-den or-ma. Ö-vé a_tenger,_az ő mű-ve, a  szárazföld_is_az_ő_keze al-ko-tá-sa. * Jöj-je-tek,_boruljunk_le,_hó-dol-junk e-lőt-te, hull-junk térdre_Urunk,_Al-ko-tónk e-lőtt, mert ő a mi Is-te-nünk, mi meg háza_népe_és_kez-é-nek nyá-ja. ANT. Bár-csak hallgatnátok_ma_az ő sza-vá-ra: „Ne le-gye-tek ke-mény szí-vű-ek, mint Meribánál,_Massza_napján_a_pusztá-ban, ahol_atyáitok_megkísértet-tek en-gem, és próbára_tettek,_bár_látták szá-mos cso-dá-mat. * Negy-ven évig_terhemre_volt_ez a nem-ze-dék, azt mondtam:_Ez_a_nép_té-vely-gő szí-vű. Nem ismerték_meg_utai-mat, ezért_haragomban_kimond-tam es-kü-met: Nyu-gal-mam_helyére_nem jut-hat-nak el!” ANT. Di-cső-ség_az_Atyának és Fi-ú-nak és Szent-lé-lek Is-ten-nek, mi-kép-pen_kezdetben_vala,_most és min-den-kor és mindörökkön_ö-rök-ké! Á-men. * ANT.}

\defsheet{InvitatoriumSol}
  {<ô-5---4--5----[------------4--4t-5--:-4--[-----------------4----4z-6--5-5---,-5---4----[------------4--4t--5--:-4--[---------------------6-5---4---eE2---?,--5----4--[----------4----4t-5--:-4--[-----------------4--4z----5--5---,-5--4----[-------------------4---4t--5-:--4--[--------4----4z--6---5--5-,-5-4--[------------4-4t-5-:-4--[-----------------------6--5--4--eE2-?,-5---4--[------------------------4----4t-5---5-:-4----[------------------4--4z-6----5-5---,-5----4-5-4--4t-5--5---:-4--[--------------------6-5---4---eE2-?,---5---4----[------------------4-4t--5--5-:--4--5--5---5---4--4z---6---5--5-,-5----4--[-------------------------------4--:-4-[---------------------4---4t--5--5--:-4--[--------------------------6---5---4---4--eE2-?,-5----4---[---------------------4-4t--5--5--:-4---[-----------------4--4z---6--5---5-,-5---4--[---------------4--:-4-[----------------------4---4t-5--5--:-4---[------------------6---5---4---eE2--?,---5--4---[-----------------4--4t-5-:-5--5-----5--4z--6--5---5--,-5--4---5---5---4---4t--5--5-:-4----4z-6---5---5-,-4--[-------------6---5---4-eE2--,-.}
  {   Jöj-je-tek, örvendezzünk az Úr-nak, és ujjongjunk_üdvünk szik-lá-ja e-lőtt. Lép-jünk színe_elé_há-la-dal-lal, ma-gasz_taljuk_őt_hangos é-nek-szó-val! ANT. Mert az Úr_valóban nagy Is-ten, az összes_isteneknél di-csőbb ki-rály. A  föld mélységei_az_ő_kezé-ben van-nak, és övé_a_he-gyek min-den or-ma. Ö-vé a_tenger,_az ő mű-ve, a  szárazföld_is_az_ő_keze al-ko-tá-sa. *  Jöj-je-tek,_boruljunk_le,_hódol-junk e--lőt-te, hull-junk_térdre_Urunk, Al-ko-tónk e-lőtt, mert ő a mi Is-te-nünk, mi meg_háza_népe_és_kez-é-nek nyá-ja. ANT. Bár-csak hallgatnátok_ma_az ő sza-vá-ra: „Ne le-gye-tek ke-mény szí-vű-ek, mint Me-ribánál,_Massza_napján_a_pusztá-ban, a-hol_atyáitok_megkísér-tet-tek en-gem, és próbára_tettek,_bár_látták szá-mos cso-dá-mat. * Negy-ven évig_terhemre_volt_ez a nem-ze-dék, azt mondtam:_Ez_a_nép té-vely-gő szí-vű. Nem is-merték_meg_utai-mat, e-zért_haragomban_kimond-tam es-kü-met: Nyu-galmam_helyére_nem jut-hat-nak el!” ANT. Di-cső-ség_az_Atyának_és Fi-ú--nak és Szent-lé-lek Is-ten-nek, mi-kép-pen kez-det-ben va-la, most és min-den-kor és mindörökkön_ö-rök-ké! Á-men. * ANT.}

% ============================================================================
% ALLELUIA (IN TONIS DIVERSIS)
% ============================================================================

\defsheet[tone={1re-a}]{AlleluiaDorian}
  {<X-3--3--4t-5---:-5--4--tT4-3r--4--eE2-1--1-?,}
  {   Al-le-lu-ja, * al-le-lu--ja, al-le--lu-ja!}

\defsheet[tone={2-g}]{AlleluiaHypodorian}
  {<ô-2--2--3--wW1-2--rR3-2--ñ---:-3--4--3--2-?,}
  {   Al-le-lu-ja, al-le--lu-ja, * al-le-lu-ja!}

\defsheet[tone={3-c}]{AlleluiaPhrygian}
  {<-4--5--7--7---:-uU6-5--4t-4-:-rR3-4tT4-2--2-?,}
  {  Al-le-lu-ja, * al--le-lu-ja, al--le---lu-ja!}

\defsheet[tone={4-a}]{AlleluiaHypophrygian}
  {<-2--1--2--1---:-0q-eE2-1w-2-?,}
  {  Al-le-lu-ja, * al-le--lu-ja!}

\defsheet[tone={5-c}]{AlleluiaLydian}
  {<-3--4--4--3---:-5--7--7i-uU6-:-5--4t-3--3-?,}
  {  Al-le-lu-ja, * al-le-lu-ja,   al-le-lu-ja!}

\defsheet[tone={6-a}]{AlleluiaHypolydian}
  {<X-3--4--5--3rf2ed1-:-3r-zZ5-rR3-3-?,}
  {   Al-le-lu-ja,     * al-le--lu--ja!}

\defsheet[tone={7do-d}]{AlleluiaMixolydian}
  {<-6--7--8o-8---:-iI7-5u-zZ5-4-:-5-5uU6-4--4-?,}
  {  Al-le-lu-ja, * al--le-lu--ja, al-le--lu-ja!}

\defsheet[tone={8szo-c}]{AlleluiaHypomixolydian}
  {<-4--4--5--rR3-:-5--7--6u-tT4R3-4t-5--4--4-?,}
  {  Al-le-lu-ja, * al-le-lu-ja,   al-le-lu-ja!}

\defsheet[tone={tper-g}]{TonusPeregrinus}
  {<X-0--1--3r-4---:-5--6--5--rR3-4t-5--4y-4y-?,}
  {   Al-le-lu-ja, * al-le-lu-ja, al-le-lu-ja!}

% ============================================================================
% TONUS DE CANTU LECTIONUM
% ============================================================================

%\defsheet{TonusDeCantuLectionum}
%  {<-------¦------------------------------7--5-.}
%  {<flexa> Boldog_az_ember,_ha_fenyíti_az Is-ten; ne vesd meg azért a Mindenható feddő szavát! Ha sebet üt rajtad, majd be is kötözi, ha szétzúz is, meggyógyít a keze. Ki látta át Isten gondolatait?}

% ============================================================================
% TE DEUM
% ============================================================================

\defsheet{TeDeum}
  {<-2--4t---5--5---5--56u-zZ5---:-2--4t--5--5z--4--5z-uU6Z5-4--?,-7--7----7-5---zZ5T4¨5z-5--:-2---5---5---5z-4--5z--uU6Z5-4--?,-¦--------------5-zZ5T4-5zZ5-5-:-2--5---5---5z-4--5z-uU6Z5-4-?,-¦------------6---5--uU6Z5T4-5zZ5-5-:-2---5---5--5z--4-5z-uU6Z5-4-?,-2tT4--5z-5-?,-2tT4--6-?,-8-----5--6--5z---4--5z--uU6Z5-4-?,--¦---------------5----zZ5T4¨5z-5-:--2---5--5z--4--5z--uU6-5---4-?,-7--7---7-5---zZ5T4¨5z-5-:-5---5z-4---5z--uU6-5----4-?,-7--7---7-----5---zZ5T4¨5z-5-:-5--5z--4---5z--uU6-5----4-?,-¦------------5----zZ5T4-5zZ5-5-:--2---5-5---5z--uU6-5---4-?,-¦-----------5---zZ5T4¨5z-5-:--5-5---5z----4---5z--uU6Z5-4-?,--2tT4-5z-5-:-7--7---7---uU6-4-5z-uU6Z5-4-?,-¦------------5--zZ5T4¨5z-5--:-2-5z---4---5z--uU6-5-4-?,-2tT4--5z--5--:-7---7---4---5z-4----3--2-?,-5---5-----5----4----5--6---4--2--?,-[---------56u-zZ5-:-2----5-5---4--5z--4--2-?,-[----------------------6---7---6---5-:-2--[------------4--5---6---4-3--2-?,-[----------6----7---6----5-:-2---[------------------4--5----6---4---2-?,-[-------------6---7----6--5-:-2-5---5---4--6---4---3---2-?,-7--5-:---6--5--5---4--6--4--2-?,-[--------------6----7--6--5---:-2----[--------------¨:-[--------------4---5---6---4---2-?,-3----2---3--1-:-2-3---£-----------5---3---rR3E2-2-?,--0--1e-¢------------------1-3---wW1Q0-0--:-£------------------3--4t-rR3E2-2-?,-0--1e----3---3---3--:-¢--------------------2--1---4---5-?,-5---7i--8--8---8----89p-oO8-:-8--7--8----9--7---7i-?,-5---7i-©---------------9-pP9-8-:-8----8-8---7--9--7----7i-?,-5---7i----©--------------9---ö-9---8-:-8--8--7---9---7---7i-?,-5----7i-©-------------------9--ö--9---8--:-8----7--8---9---7--7i-?,-5--7i---8--88(---8--8---8---8--9---ö--9---8---:-5----©-----------7----9--7---5-?,-X6---6--4---6--6-----7i-uU6Z5-5--:-3-X4z---6-4---6--7iI7-zZ5T4-5-?,-X4uU667iI7U6Z5T4-5-.}
  {  Is-ten, té-ged di-csé-rünk, * té-ged Úr-nak mi is val---lunk. Té-ged, ö-rök Is-------ten, min-den föl-di ál-lat di----csér. Tenéked_mennye-i an----gya--lok és men-nye-i  fe-je-del---mek, tenéked_kéru-bim és szé-----ra---fim min-den-ko-ron é-ne-kel---nek: szent Is-ten! Szent Úr!  Szent Úr Is-ten, di-cső Ki----rály! Teljes_menynyor-szág és       föld fel-sé-ges di-cső-sé--ged-del. Té-ged a-pos-to-------lok tár-sa-ság-gal di--csér-nek. Té-ged szent pró-fé-------ták vi-gas-ság-gal di--csér-nek. Téged_kínval-lott már---tí---rok, kik é-ret-ted meg-hol-tak, téged_ez_vi-lág sze------rint A-nya-szent-egy-ház di----csér, fel--sé-ges ma-gas-ság-nak ő Is-te----nét, mi_tisztelen-dő U--------runk e-gyet-len egy Fi--a-dat, Szent-lel-ket, min-ket meg-vi-gasz-ta-lót. Te, Krisz-tus, vagy di-cső ki-rály. Te_Atyais-ten-nek   vagy ö-rök-sé-ges Fi-a.   Te_megváltásra_embersé-get fel-ven-ni  Má-riának_szent mé-hét nem u-tá-lád. Te_örök_ha-lált meg-győz-vén men-nyei_szent_kapukat ne-künk meg-nyi-tál. Te_ülsz_Isten-nek jobb-já-ra, A-tyá-nak di-cső-sé-gé-ben.   Hi-szünk el-jö-ven-dő bí-ró-nak. Téged_azért,_U-ram, mi ké-rünk, hogy minket_segélj, @  kiket_szent_vé-red-del meg-vál-tál. Szen-tid kö-zé  ö-rök dicsőségben min-ket szám--lálj. Üd-vö-zítsed_te_népedet, U-ram Is----ten, és_megáldjad_te_ma-ra-dé-ki----dat. És tarts meg min-ket, és_elfelvígy_az_örök di-cső-ség-be!  Min-den na-pon-ként té--ged,  Úr Is-ten, di-csé-rünk. És, U--ram,_te_nevedet ö-rök-ké, mind-ö-rök-ké di-csér-jük.  Mél-tólj, Uram_Isten_min-ket ő-riz-ned ez na-pon bűn nél-kül!  Légy ir-galmas,_Uram_Isten, bű-nö-sök-nek, légy ir-gal-mas ne-künk! Le-gyen te szent ir-gal-mas-sá-god mi-hoz-zánk, mert minekünk_te vagy ol-tal-munk. Te-ben-ned bí-zunk, Úr Is----ten, a--zért ö-rök-ké mi   é-----lünk. Á---------------men.}

% ============================================================================
% MELODIAE HYMNORUM
% ============================================================================

% Hymnus
\defsheet{NepekMegvalto}
  {<ô-2--2--2---1---2---4---3--2í--:-2--3----4--5------4--5---5--6x--,-6----6----5---6--4----5--rR3-2í-:-2--1----2--4--5----2-1--2í-?,-1e4R3E2W1-2í-.}
  {   El-ső nap min-den nap kö-zött, me-lyen ég s_föld al-kot-ta-tott, s_me-lyen fel-tá-madt Al-ko--tónk le-győz-te ér-tünk a ha-lált, Á---------men.}

% Hymnus
\defsheet{IttAzAlkalmas}
  {<ô-4--4----4---4---4---4--3--2í-:-4--4-----4-4-----1--2---1--0í-,-2--2--3---4---3----2---1--2eE2í-:-2---2---1--0--1--2---1--1í-?,-1wW1-0\\q\\-.}
  {   Fé-nyes-ség-osz-tó, bő-ke-zű,  ki fényt á-raszt-va gaz-da-gon  az éj-nek gyá-szát osz-la-tod,    s_a nap vi-lá-ga tűz le ránk. Á----men.}

% Hymnus
\defsheet{ANapvialg}
  {<-5--5--5--5--1---3--2--1í--:-5--5--6---4----5-----7----6--5x-,-5-----6---7--8----7----5---7--zZ5T4x-:-5---5--6---4--2---4---5--5x-?,-5zZ5-4\\t\\-.}
  {  Ég Al-ko-tó-ja, Is-te-nünk, ki éj-jel hold-fényt adsz ne-künk s_nap-pal na-pot, mely kör-be fut,     s_e-lő-írt út-ján cél-ba jut.  Á----men.}

% Hymnus
\defsheet{IttJelenVagyon}
  {<X-5--4----5--4---1í-:-4--4--3---4--5x-5x-,-5---7---8-7---5x-:-5--4---3---4--5x---5x-,-5---4---5----4---1í-:-4-4--3---4-tT4-3y-,-3---3---2-1í---1í-?,-1wW1-0\\q\\-.}
  {   Sá-pad, el-osz-lik  kö-de már az éj-nek, pir-kad a haj-nal, ró-zsa-szí-nű fény kél; for-dul-junk buz-gón  a vi-lág U-rá--hoz  han-gos i-mánk-kal,  Á----men.}

% Hymnus
\defsheet{IttAzAlkalmasSzentIdo}
  {<-7---7---7-7----7---7----6--5-:-7-7----7--7---4---5--4-3--,-5---5---6--7---6--5----4--5zZ5-:-5--5----4--3--4---5-4--4-?,--4tT4-3\\r\\-.}
  {  Már kél a fény-nek csil-la-ga: e-seng-ve kér-jük az U-rat, jár-jon ma min-de-nütt ve-lünk,  ne ront-sa ár-tás é-le-tünk. Á----men.}

% Hymnus
\defsheet{KiMindenFenynek}
  {<-3--3---4---5---4---4-----3--2---:-4--2-----3----1---0--1----3--3--,-3----3---4--5----4---4--3-2---:-4--2---3--1---0---3----1--1-?,-02eE2W1Q0-1-.}
  {  Üd-vös-ség nap-ja, Krisz-tu-sunk, ha-sítsd szét vak-sá-gunk kö-dét, hogy fel-ra-gyog-jon az e-rény, mi-dőn re-ánk nap-palt ho-zol. Á---------men.}

% Hymnus
\defsheet{IstenTeMindentAlkoto}
  {<-5--5---5--5---5--5--4--4-,-5--5--5----4-3---4--5--5-,--5-5---5-4--3---4---3--2----,-3-------4--4--3--4---5-4--4-?,-4tT4-3\\r\\-.}
  {  Vi-lág te-rem-tő Is-te-ne, ki Úr vagy é-jen és na-pon, i-dőt i-dő-vel vál-to-gatsz, s_nincs mű-ve-id-ben u-na-lom, Á----men.}

% ============================================================================
% DEFTONES
% ============================================================================

\deftone{1la-a}{<-[-4-3-rR3-1-.}{1. lá tónus}
\deftone{1re-a}{<-[-4-3-4t-4-.}{1. ré tónus}
\deftone{1re-b}{<-¥-5-4-5z-5-.}{1. ré tónus}

\deftone{2-f}{<-¢-0q-1-.}{2. tónus}
\deftone{2-g}{<-£-1w-2-.}{2. tónus}
\deftone{2-b}{<-¥-3r-4-.}{2. tónus}
\deftone{2-c}{<-¦-4t-5-.}{2. tónus}
\deftone{2-d}{<-©-5z-6-.}{2. tónus}

\deftone{3-b}{<-¥-4z-4-3-.}{3. tónus}
\deftone{3-c}{<-¦-5u-5-4-.}{3. tónus}

\deftone{4-g}{<-£-3-4-3-1-.}{4. tónus}
\deftone{4-a}{<-[-4-5-4-2-.}{4. tónus}

\deftone{5-c}{<-¦-8-6-7-5-.}{5. tónus}
\deftone{5-g}{<-£-5-3-4-2-.}{5. tónus}

\deftone{6-a}{<-[-4-3-4-3-.}{6. tónus}

\deftone{7do-a}{<-[-6-5-rR3-4-.}{7. dó tónus}
\deftone{7do-d}{<-©-9-8-uU6-7-.}{7. dó tónus}
\deftone{7la-a}{<-[-6-5-rR3-2-.}{7. lá tónus}
\deftone{7la-d}{<-©-9-8-uU6-5-.}{7. lá tónus}

\deftone{8szo-g}{<-£-3-4-2-1-.}{8. szó tónus}
\deftone{8szo-b}{<-¥-5-6-4-3-.}{8. szó tónus}
\deftone{8szo-c}{<-¦-6-7-5-4-.}{8. szó tónus}
\deftone{8do-c}{<-¦-5-7-8-7-.}{8. dó tónus}

\deftone{8szo-a-c}{<-2e-[-…-[-4-5-6-‡-5-,-[-4-5-3-„-2-.}{8. szó tónus}

\deftone{tirr-e}{<-¡-3-1-3-2-.}{ton.irreg.}
\deftone{tirr-a}{<-[-6-4-6-5-.}{ton.irreg.}
\deftone{tirr-b}{<-¥-7-5-7-6-.}{ton.irreg.}

\deftone{tper-g}{<-£-5-1-eE2-1-.}{ton.per.}

\title{Tábori lectionarium}
\subtitle{902. Kucsera Ferenc cserkészcsapat - 2026. nyári tábor}

\begin{document}

\maketitle

\section*{Évközi 14. vasárnap - Előtábor 2. nap}

\parttitle[12, 1-25]{Első olvasmány}
Sámuel második könyvéből

\begin{lectio}{Dávid bűnbánata}
\noindent Azokban a napokban az Úr elküldte Nátán prófétát Dávidhoz. El is ment hozzá, és így beszélt: „Egy városban élt két ember. Az egyik gazdag volt, a másik szegény. A gazdagnak rengeteg juha, kecskéje és barma volt. A szegénynek azonban nem volt egyebe, csak egy kisbáránya, vette magának. Etette, úgyhogy megnőtt, vele magával és gyermekeivel. Evett a kenyeréből, és ivott a poharából. Az ölében aludt, s olyan volt, mintha a lánya lett volna. Egyszer vendég érkezett a gazdag emberhez. De nem tudta rászánni magát, hogy a juhai vagy barmai közül egyet is elvegyen, és a vendégnek, aki érkezett, elkészítse. Ezért elvette a szegény ember báránykáját, és azt készítette el vendégének.” Dávid nagyon megharagudott az emberre, és azt mondta Nátánnak: „Amint igaz, hogy az Úr él: az az ember, aki ilyet tesz, az a halál fia. A bárányt négyszeresen vissza kell fizetnie, amiért így tett, és nem volt kímélettel.”

Erre Nátán így szólt Dávidhoz: „Magad vagy ez az ember! Ezt mondja hát az Úr, Izrael Istene: Fölkentelek Izrael királyává, és kiszabadítottalak Saul kezéből. Neked adtam urad házát és kebledre urad asszonyait, rád bíztam Izrael házát és Júda házát. S ha ez kevés lett volna, még megtetéztem volna ezzel meg azzal. Miért vetetted hát meg az Urat, s tettél olyat, ami gonosznak számít az Úr szemében? Kard által elveszítetted a hetita Úriját, hogy feleségül vedd a feleségét. Igen, elveszítetted az ammoniták kardja által! Nos, a kard sose fordul el házadtól, amiért megvetettél, és elvetted a hetita Úrijától a feleségét, hogy a te asszonyod legyen. Ezt mondja az Úr: Nos, a saját házadból hozok rád romlást. Elveszem a szemed láttára asszonyaidat, és odaadom őket más valakinek, hogy napvilágnál asszonyaiddal háljon. Te titokban tetted, de én egész Izrael szeme láttára, napvilágnál viszem ezt végbe!”

Erre Dávid így szólt Nátánhoz: „Vétkeztem az Úr ellen!” Nátán ezt válaszolta Dávidnak: „Így az Úr is megbocsátja bűnödet, s nem halsz meg. De mert ezzel a tetteddel káromoltad az Urat, a fiú, aki született neked, meghal.” Ezzel Nátán elment haza.

Az Úr megverte a gyermeket, akit Úrija felesége szült Dávidnak, úgyhogy súlyosan megbetegedett. Dávid az Úrhoz fordult a fiú miatt; szigorú böjtöt tartott, s amikor hazament, éjszaka a puszta földön aludt zsákkal takarózva. Udvarából a főemberek eléje járultak, és igyekeztek rábeszélni, hogy keljen föl a földről, de nem akart, s nem is evett velük. A hetedik nap a gyermek meghalt. Dávid udvari emberei nem merték neki jelenteni, hogy a gyerek meghalt. Azt mondták magukban: „Már akkor sem hallgatott ránk, amikor még élt a gyerek, s a lelkére beszéltünk. Hát hogy mondhatnánk akkor meg neki: a fiú halott? Még kárt tenne magában!” De Dávid észrevette, hogy udvari emberei suttogtak egymás közt. Ebből megértette, hogy a gyerek meghalt. Ezért Dávid megkérdezte udvari embereitől: „Meghalt a gyerek?” „Meg” – felelték.

Erre Dávid fölkelt a földről, megmosdott, bekente magát olajjal, és más ruhát öltött. Aztán elment az Úr házába, és leborult. Ezután ételt tettek eléje, és evett. Erre udvari emberei így szóltak hozzá: „Mit csinálsz? Amikor még élt a gyerek, böjtöltél, és siránkoztál. Most meg, hogy a gyerek meghalt, fölkelsz, és ételt veszel magadhoz.” Ezt válaszolta: „Amíg a gyermek még élt, böjtöltem és sírtam, mert azt gondoltam, ki tudja, hátha megkönyörül rajtam az Úr, és életben marad a gyerek. De most, hogy meghalt, minek böjtöljek tovább? Hát visszahozhatom? Én utána megyek, de ő nem tér hozzám vissza.”

Dávid Batsebát is vigasztalta. Elment hozzá, és együtt hált vele. (Batseba) fogant, és fiút szült. A Salamon nevet adta neki. Az Úr kedvét találta benne, s ezt tudtul is adta Nátán próféta által, aki az Úr szavára Jedidjának nevezte el.
\end{lectio}

\parttitle[Manassze imája 9. 10. 12; Zsolt 50, 5. 6]{Válaszos ének}

\versicle{%
    Megsokasodtak vétkeim, és bűneimnek száma, mint a teng\underline{e}r hom\underline{o}kja; † bűneim sokasága miatt nem vagyok méltó, hogy lássam a magas eget, mert haragra inger\underline{e}lt\underline{e}lek, * És ami színed előtt gonosz, \underline{o}lyat t\underline{e}ttem.
}{%
    Gonoszságomat bel\underline{á}tom, † bűnöm előttem lebeg szünt\underline{e}len, * egyedül csak ellened v\underline{é}tk\underline{e}ztem.
}

\parttitle{Második olvasmány}
Szent Ágoston beszédeiből

{\centering\redtext{(Serm. 19, 2-3: CCL 41, 252-254)}}

\begin{lectio}{Istennek tetsző áldozat a megtört szív}
\noindent \textit{Gonoszságomat belátom} (Zsolt 50, 5), mondja Dávid. Ha én belátom, akkor te ne ródd föl azt nekem, Istenem. Még ha becsületesen éltünk is, ne gondoljuk, hogy nincs bűnünk. Életünk annyira lesz dicséretes, amennyire Isten bocsánatát kérjük. Kétségbeejtő az olyan ember, aki nem annyira saját bűneit nézi, de annál inkább vájkál a másokéban. Szimatolja, de nem azért, hogy javítson rajta, hanem mert marni akarja. Mivel magát nem tudja tisztára mosni, ezért inkább másokat igyekszik bemocskolni. Mi ne így tegyünk, amikor Dávid példájára imádkozunk, és elégtételt akarunk nyújtani Istennek. Dávid azt mondta: \textit{Gonoszságomat belátom, bűnöm előttem lebeg szüntelen} (Zsolt 50, 5). Nem figyelt ő más bűnére, hanem saját magát vádolta. Nemcsak a felszínen tapogatódzott, hanem mélyen magába szállt. Saját magának nem kegyelmezett, hogy a maga számára ne meggondolatlanul kérje a megkegyelmezést.

Ki akarsz engesztelődni Istennel? Gondolkozz: mit kell tenned, hogy Isten megbékéljen veled. Figyeld meg, mit olvasunk ugyanebben a zsoltárban: \textit{Te a véres áldozatot nem kedveled, égőáldozataimban nem leled kedved} (Zsolt 50, 18). Hát akkor ne is legyen már áldozatod? Ne ajánlj fel neki semmit? Hát semmiféle áldozattal sem lehet Istent kiengesztelned? Hogy is mondtad a zsoltárban? Te a véres áldozatot nem kedveled, égőáldozataimban nem leled kedved. Folytasd csak tovább, és halld: \textit{Isten előtt kedves áldozat a megtört lélek, te nem veted meg a töredelmes és alázatos szívet} (Zsolt 50, 19). Vesd el, amit hajdan áldoztál, hiszen megtaláltad, hogy ezután mit áldozzál. Hajdan atyáid állatokat öltek le az oltáron, és azt tartották áldozatnak. Te a véres áldozatot nem kedveled. Tehát azokat már ne keresd, hanem keresd a te áldozatodat.

Égőáldozataimban – mondja – nem leled kedved. Tehát ha égőáldozataimra rá sem nézel, vajon akkor semmiféle áldozat sem jó neked? Dehogynem! Isten előtt kedves áldozat a megtört lélek, te nem veted meg a töredelmes és alázatos szívet. Van tehát, mit áldozzál. Hát akkor ne nyájak körül nézelődj, de hajóútra se készülj, hogy távoli országokból drága illatszereket hozzál. Saját szívedben keressed, mi volna kedves Istennek. Szívedet kell megtörnöd. Mit félsz attól, hogy elvész, ha összetöröd? Hisz ott áll a zsoltárban megírva: \textit{Tiszta szívet teremts bennem, Istenem!} (Zsolt 50, 12) Hogy tiszta szívet lehessen benned teremteni, ahhoz a tisztátalant össze kell törnöd.

Ne legyen tetszésünkre az, ha vétkezünk, mert a bűn nem tetszik Istennek. Mivel nem vagyunk bűn nélkül, legalább abban hasonlítsunk Istenhez, hogy nekünk se legyen a tetszésünkre az, ami neki nem tetszik. Legalább annyira egyesülj Isten akaratával, hogy neked sem tetszik tenmagadban, amit utál az, aki téged teremtett.
\end{lectio}

\parttitle[Zsolt 50, 12]{Válaszos ének}

\versicle{%
    Bűneim, Uram, mint a nyilak átszög\underline{e}ztek \underline{e}ngem, † de mielőtt sebeim elmérges\underline{e}d\-n\underline{é}nek, * Uram, Istenem gyógyíts meg engem a bűnbocsánat gy\underline{ó}gyszer\underline{é}vel.
}{%
    Tiszta szívet teremts bennem, Ist\underline{e}nem; * új és erős lelket \underline{ö}nts b\underline{e}lém.
}

\section*{Hétfő, 14. hét - Előtábor 3. nap}

\parttitle[15, 7-14. 24-30; 16, 5-13]{Első olvasmány}
Sámuel második könyvéből

\begin{lectio}{Absalom lázadása és Dávid menekülése}
\noindent Azokban a napokban Absalom így szólt Dávid királyhoz: „El szeretnék menni Hebronba, hogy fogadalmamat, amelyet az Úrnak tettem, teljesítsem. Ezt a fogadalmat tette ugyanis szolgád, amikor az arámok földjén, Gesurban voltam: Ha az Úr visszavezérel Jeruzsálembe, akkor az Úr tiszteletére istentiszteletet tartok Hebronban.” A király ezt mondta neki: „Menj békével!” Útra kelt hát, és elment Hebronba. Ám Absalom titokban követeket küldött Izrael törzseihez ezzel az üzenettel: „Ha meghalljátok a harsonák szavát, kiáltsátok ki: Absalom lett Hebron királya!” Absalommal kétszáz férfi elhagyta Jeruzsálemet. Meghívta őket, így gyanútlanul elmentek vele, és mit sem tudtak az egész dologról. A gilonita Achitofelt, Dávid tanácsadóját is odahívatta Absalom a városból, Gilóból, és együtt mutatták be az áldozatot. Az összeesküvés tehát nagy volt, s egyre többen gyűltek Absalom köré.

Hírhozó érkezett Dávidhoz, és jelentette neki: „Izrael népének szíve Absalomhoz fordult.” Erre Dávid így szólt tisztjeihez, akik vele voltak Jeruzsálemben: „Rajta, meneküljünk! Különben nem kerüljük el Absalom kezét. El innét, igyekezzetek, nehogy siessen, és rajtunk üssön, romlást hozzon ránk, és kardélre hányja a várost!”

Ott volt Cádok pap is az összes levitával együtt; ők vitték az Isten ládáját. Míg a nép ki nem vonult a városból, addig az Isten ládáját letették Ebjatár közelében. Akkor a király azt mondta Cádoknak: „Vidd vissza az Isten ládáját a városba! Ha tetszésre találok az Úr szemében, visszavezérel, és viszontláthatom hajlékát. Ha azonban azt mondja: Nem vagy tetszésemre, hát akkor jó, itt vagyok, tegyen velem, amint jónak látja.” Aztán azt mondta a király Cádok papnak: „Nézzétek, te és Ebjatár térjetek vissza nyugodtan a városba, ugyanígy a fiaitok is veletek együtt, a te fiad, Achimaac és Ebjatár fia, Jonatán. Magam a pusztában a gázlónál várok addig, amíg hírt adtok nekem.” Így hát Cádok és Ebjatár visszavitték az Isten ládáját Jeruzsálembe, és ott maradtak. Akkor Dávid elindult fölfelé az Olajfák-hegyén. Egyre sírt, mezítláb volt, s a feje befödve. A nép, amely vele tartott, szintén befödte a fejét, és folytonos nagy sírás közepette haladt fölfelé.

Amikor Dávid király Bachurim közelébe ért, szembejött vele egy ember, ugyanabból a nemzetségből, amelyből Saul családja is származott. Siminek hívták, és Gera fia volt. Egyfolytában átkozódott, ahogy jött, sőt köveket dobált Dávid után és Dávid király tisztjei után, noha a harcosok és a testőrök jobbról és balról körülfogták a királyt. Így hangzottak Simi átkai: „Ki innét, ki, te vérengző, elvetemült ember! Most rád hárította az Úr Saul egész házának vérét, akitől elvetted a királyságot, és fiadnak, Absalomnak kezére juttatta. Látod, most aztán bajba jutottál, amiért olyan vérengző vagy.”

Ceruja fia, Abisáj megkérdezte a királytól: „Mit átkozza ez a döglött kutya uramat, királyomat? Engedd meg, hogy odamenjek, és fejét vegyem!” De a király így válaszolt: „Mit akartok tőlem, Ceruja fia? Hagyjátok, hadd átkozódjék! Ha az Úr meghagyta neki: Átkozd meg Dávidot, akkor ki kérheti számon: miért teszel ilyet?” Aztán azt mondta Dávid Abisájnak és tisztjeinek: „Ha a saját fiam, aki ágyékomból származik, az életemre tör, mennyivel inkább Benjaminnak ez a fia itt. Hagyjátok, hadd átkozzon! Mert az Úr hagyta meg neki. Talán majd megszán az Úr nyomorúságomban, és jót ad ahelyett az átok helyett, amely ma ér.” Ezzel Dávid ment embereivel együtt tovább az útján. Simi azonban szintén lement a lejtőn, s folyvást átkozta, kővel dobálta, és fölkavarta előtte a port.
\end{lectio}

\parttitle[Zsolt 40,10; Mk 14,18]{Válaszos ének}

\versicle{%
    Még b\underline{a}rát\underline{o}m is, † akiben ped\underline{i}g b\underline{í}ztam, * Aki
    együtt ette kenyerét velem, sarkát emelt\underline{e} ell\underline{e}nem.
}{%
    Egyiketek elárul \underline{e}ngem, * egy, aki esz\underline{i}k v\underline{e}lem.
}

\parttitle{Második olvasmány}
Szent I. Kelemen pápának a korintusiakhoz írt leveléből

{\centering\redtext{(Nn 46, 2 – 47, 4; 48, 1-6: Funk 1, 119-123)}}

\begin{lectio}{Mindegyikünk azt keresse, ami mindenkinek hasznos, és nemcsak saját magának}
\noindent Írva van: \textit{Legyetek a szentek társai, mert akik hozzájuk csatlakoznak, maguk is szentek lesznek} (vö. Ef 2, 19). És egy más helyen ismét: \textit{Az ártatlannal ártatlan leszel, a választottal kiválasztott, a gonosszal elvetemülsz} (vö. Zsolt 17, 26). Ezért az ártatlanokkal és igazakkal társuljunk; ők ugyanis Isten választottai. Miért van köztetek viszály, harag, pártoskodás, szakadás és háborúság? Nemde, egy a mi Istenünk, egy a Krisztusunk, egy a kegyelemnek ránk áradó Lelke, és egy a mi hivatásunk Krisztusban? Miért tépjük, szaggatjuk szét Krisztus tagjait, és miért lázadunk saját testünk ellen? Elvakultságunkban megfeledkezünk arról, hogy egymásnak tagjai vagyunk?

Emlékezzetek Urunk, Jézus szavaira. Azt mondotta: \textit{Jaj annak az embernek! Jobb lett volna neki, ha nem született volna, minthogy megbotránkoztasson egyet is választottaim közül; jobb lett volna, ha malomkövet kötöttek volna a nyakára, és a tengerbe vetették volna, mint hogy egyet is megrontson választottaim közül} (vö. Lk 17, 1-2). A köztetek levő szakadás sokakat megront, sokakat lelkileg elcsüggeszt, sokakat megingat, mindnyájunknak szomorúságot okoz, mindnyájunkat fájdalommal tölt el. Mégis megmarad köztetek ez a lázongás.

Vegyétek csak kezetekbe Szent Pál apostol levelét. Mit ír nektek üzenetének elején? Bizonyosan isteni sugalmazásra magáról, Kéfásról és Apollóról ír hozzátok, mivel akkor is szakadás és pártoskodás volt közöttetek. De akkori pártoskodástok kisebb bűn volt, mert a legkiválóbb apostolok és az őáltaluk ajánlott férfiú között oszoltatok meg.

Hamar szüntessük ezt meg; boruljunk le az Úr lábához, és esdeklőn, sírva könyörögjünk: legyen irgalmas hozzánk, engesztelődjön meg, és állítsa vissza köztünk a régi dicső és tiszta testvéries életmódot. Mert ez az élethez vezető igazság kapuja, amint írva van: \textit{Nyissátok meg az igazság kapuit; belépvén rajta magasztalom az Urat. Ez az igazság kapuja; az igazak mennek be rajta} (vö. Zsolt 117, 19). Sok kapu van nyitva, de az, amely az igazság kapuja, ugyanaz, mint amely Krisztushoz vezet. Boldogok mindazok, akik rajta bemennek, és útjukat szentségben és igazságban irányították, mindent zavartalanul teljesítve. Legyen bár valaki hűséges, legyen bár tehetsége az isteni tanítás hirdetésében, legyen bár bölcs a beszédek megítélésében, legyen bár tiszta tetteiben: annál alázatosabbnak kell lennie, minél nagyobbnak látszik. Azt keresse, ami mindenkinek hasznos, és nemcsak saját magának.
\end{lectio}

\parttitle[1 Kor 9,19. 22; Jób 29,15-16]{Válaszos ének}

\versicle{%
    Bár mindenkitől független voltam, mégis mindenkinek szolg\underline{á}ja l\underline{e}ttem. † A
    gyöngék közt gyöng\underline{e} l\underline{e}ttem. * Mindenkinek mindene lettem, hogy mindenkit
    \underline{ü}dvöz\underline{í}tsek.
}{%
    A vaknak úgy szolgáltam, hogy a szeme v\underline{o}ltam;; * a szegény embereknek
    atyj\underline{u}k v\underline{o}ltam.
}

\section*{Kedd, 14. hét - Előtábor 4. nap}

\parttitle[18, 6-17. 24 – 19, 4]{Első olvasmány}
Sámuel második könyvéből

\begin{lectio}{Absalom halála és Dávid gyásza}
\noindent Azokban a napokban: Dávid serege harcba vonult Izrael ellen. Efraim erdeiben ütköztek meg. Izrael fiai vereséget szenvedtek Dávidtól. Óriási veszteségük volt azon a napon: húszezer ember. A harc az egész vidékre kiterjedt. S az erdő még több áldozatot szedett a seregből, mint a kard azon a napon.

Absalom véletlenül útjába került Dávid embereinek. Absalom ugyanis egy öszvéren ült, és az öszvér egy nagy tölgyfa koronája alá tévedt. Absalom feje beleakadt a tölgyfába, s ott maradt ég és föld között himbálózva, az öszvér meg, amely alatta volt, tovább ment. Egy ember meglátta, és jelentette Joábnak: „Láttam Absalomot – mondta –, egy tölgyfán lógott.” Joáb így felelt az embernek, aki a hírt vitte neki: „Hogyan, hát ha láttad, akkor miért nem ölted meg ott helyben, és fektetted a földre? Aztán az én dolgom lett volna, hogy adjak neked tíz ezüstöt meg egy övet.” De az ember ezt válaszolta Joábnak: „Ha ezer ezüst ütötte volna is markomat, akkor se merészeltem volna kezet vetni a király fiára. Hát nem parancsba adta neked, Abisájnak és Ittainak a fülünk hallatára a király: »Vigyázzatok nekem a fiúra, Absalomra!« Ha meg alattomban elbántam volna vele, akkor téged állított volna félre, hiszen a király előtt nem marad semmi titokban.” Joáb azt mondta rá: „Eh, mit! Miért is töltöm tovább az időt veled?”

Ezzel kezébe fogott három lándzsát, és Absalom mellébe döfte őket. Mivel azonban még élve lógott a tölgyfán, odalépett tíz ember, Joáb fegyverhordozói, és agyonütötték. Ekkor Joáb megfúvatta a kürtöket, és az emberek abbahagyták Izrael üldözését, hiszen Joáb megállást parancsolt a seregnek. Aztán fogták Absalomot, s az erdőben egy gödörbe dobták, és egy óriási kőhalmot raktak fölé. Közben egész Izrael menekült, ki-ki a maga sátorába.

Dávid (Machanajimban) a két kapu közt ült. Az őrszem fölment a kapu tetejére. Amikor fölemelte tekintetét, azt látta, hogy egy ember fut felé. Az őrszem közölte a hírt a királlyal. A király azt mondta: „Ha egyedül van, jó hírt hoz a szája.” Amikor azonban az – tovább futva – közelebb ért, az őrszem egy másik embert is látott, amint utána futott. Az őrszem tehát leszólt a kapuba: „Még egy embert látok, egyedül fut.” A király azt mondta: „Az is jó hírt hoz.” Az őrszem újra jelentette: „Úgy látom az első úgy fut, ahogy Cádok fia, Achimaac fut.” A király visszaszólt: „Az derék ember, biztosan jó hírrel jön.”

Achimaac megérkezett, és odakiáltotta a királynak: „Üdv!” Aztán arcra borult a földön a király előtt, és így szólt: „Legyen áldott az Úr, a te Istened, aki cserbenhagyta azokat az embereket, akik uram, királyomra kezet emeltek!” A király megkérdezte: „Jól van a fiú, Absalom?” „Csak egy nagy zűrzavart láttam – válaszolta Achimaac –, amikor a király szolgája, Joáb elküldte szolgádat. Nem tudom, mi történt.” A király megparancsolta: „Állj félre, és maradj ott!” S ő félreállt, és ott maradt.

Most megérkezett a kusita is, és jelentette: „Örömhírt hall uram és királyom! Az Úr ma igazságot szolgáltatott neked azokkal szemben, akik fellázadtak ellened.” A király azonban megkérdezte a kusitától: „Jól van a fiú, Absalom?” „Bárcsak az történne uram és királyom ellenségeivel és azokkal is mind, akik a vesztedre törnek, ami a fiatalemberrel!” – válaszolta a kusita.

Dávid megrendült. Fölment a kapu fölötti szobába, és sírt. Elcsukló hangon így siránkozott: „Fiam, Absalom! Fiam, fiam, Absalom! Ó, bár én haltam volna meg helyetted!” Jelentették Joábnak: „A király sír és jajgat Absalom miatt.” Így azon a napon gyászra fordult a győzelem az egész sereg számára, mert a sereg megtudta azon a napon, hogy a király bánkódik a fia miatt. S ezen a napon a sereg lopva tért vissza a városba, ahogyan az a sereg lopakodik, amelyik szégyenkezik, mivel megfutamodott a csatából. A király befödte az arcát, és hangosan siránkozott: „Fiam, Absalom! Absalom, fiam, fiam!”
\end{lectio}

\parttitle[Zsolt 54, 13. 14. 15; vö. 40, 10; 2 Sám 18, 33]{Válaszos ének}

\versicle{%
    Hogyha ellenségem átk\underline{o}zna \underline{e}ngem, † azt még
    elv\underline{i}s\underline{e}lném; * De te, olyan ember, aki egyenlő velem, akivel oly édes
    volt a barátságunk, sarkadat emelt\underline{e}d ell\underline{e}nem.
}{%
    Megrendült a k\underline{i}rály, † fölment a kapu fölötti szobába, \underline{é}s sírt, * s
    járva-kelve ezt mond\underline{o}g\underline{a}tta:
}

\parttitle{Második olvasmány}
Szent Ágoston püspök zsoltármagyarázataiból

{\centering\redtext{(Ps 32, 29: CCL 38, 272-273)}}

\begin{lectio}{Akik távol vannak, akár akarják, akár nem, a mi testvéreink}
\noindent Testvéreim! Leginkább arra a szeretetre buzdítunk titeket, amely nemcsak egymásra, hanem azokra is kiterjed, akik távol vannak. Akár a pogányokra, akik még nem hisznek Krisztusban, akár a tőlünk elszakadtakra, akik Főnket elismerik, de testétől távol vannak. Érezzünk fájdalmat miattuk, testvérek, úgy mint testvéreink felett. Akár akarják, akár nem akarják, a mi testvéreink ők. Csak akkor szűnnének meg testvéreink lenni, ha felhagynának azt mondani: \textit{Mi Atyánk} (Mt 6, 9).

A Próféta egyesekről így szólt: \textit{Azoknak, akik azt mondják nektek: „Nem vagytok a mi testvéreink” mondjátok: – Testvéreink vagytok”} (Iz 66, 5 a LXX szerint). Nézzetek körül, vajon kikről mondhatta ezt? Vajon a pogányokról? Nem, hiszen a Szentírás és az egyházi szóhasználat szerint ezeket nem nevezzük testvéreknek. Vajon a zsidókról, akik nem hittek Krisztusban?

Olvassátok az Apostolt és lássátok, hogy amikor „testvért” mond minden hozzáadás nélkül, csak a keresztényeket érti ez alatt: \textit{Miért ítéled el tehát testvéredet, vagy miért nézed le testvéredet?} (Róm 14, 10) És más helyen így szól: \textit{Ti igazságtalanságot követtek el, és kárt okoztok, méghozzá testvéreiteknek} (vö. 1 Kor 6, 8).

Azok tehát, akik ezt mondják: „Nem vagytok a mi testvéreink”, pogányoknak mondanak minket. Ezért akarnak tehát minket újrakeresztelni, mondván, hogy amit ők adnak, az nincs meg nekünk. Ebből nyilvánvaló a tévedésük, midőn tagadják, hogy mi az ő testvéreik lennénk. De miért mondja nekünk a Próféta: \textit{Mondjátok nekik: „Testvéreink vagytok”}, hacsak nem azért, mivel mi elismerjük bennük a keresztséget, amelyet nem kívánunk megismételni náluk? Ők ugyanis amikor nem ismerik el a mi keresztségünket, tagadják, hogy testvérek vagyunk. Mi viszont elismerve a mienket, nem követeljük az ő keresztségük megismétlését, mert azt mondjuk nekik: \textit{Testvéreink vagytok}.

Mondják csak ők: „Mit kerestek minket, mit akartok tőlünk?” Azt feleljük erre: \textit{Testvéreink vagytok}. Ők azt mondják: „Menjetek tőlünk, nincs veletek közös ügyünk.” Nekünk viszont van veletek közös ügyünk: egy Krisztusban hiszünk, tehát egy testben egy Fő alatt kell lennünk.

Testvérek! Esküvel kérünk titeket a szeretet belső érzése által, a mi Urunk, Jézus Krisztus által, akinek kenyere megerősít minket, akinek teje táplál minket, a szelídsége által, esküvel kérünk titeket, hogy szeretetetek belső érzéseivel könyörögjetek Istennél érettük. (Itt van ugyanis az ideje, hogy igen nagy szeretetet tanúsítsunk velük szemben, túláradó irgalommal könyörögjünk értük Istennél, hogy végre megadja nekik a tisztánlátást, hogy észre térjenek és meglássák, hogy semmijük sincs, amikor az igazság ellen beszélnek. Nem maradt nekik más, mint csak a heveskedés, amely annál erőtlenebb, minél erősebbnek érzi magát.) Könyörögjetek a betegekért, a testileg bölcsekért, az állatiasságban, testiségben levőkért, de mégis csak testvéreinkért, akik ugyanazon szentségeket ünneplik, még ha nem is velünk együtt, de mégis ugyanazokat. Akik ugyanazt az \textit{Ámen}t válaszolják, még ha nem is velünk együtt, de mégis ugyanazt.
\end{lectio}

\parttitle[Ef 4, 1. 3. 4]{Válaszos ének}

\versicle{%
    Kérlek bennetek\underline{e}t az \underline{Ú}rban, † hogy éljetek méltón ahhoz a hivatáshoz,
    amelyet k\underline{a}pt\underline{a}tok. * Törekedjetek rá, hogy a béke kötelékével
    fenntartsátok a lelk\underline{i} egys\underline{é}get.
}{%
    Egy a Test és egy a L\underline{é}lek, * mint ahogy hivatástok is egy
    rem\underline{é}nyr\underline{e} szól.
}

\section*{Szerda, 14. hét - Előtábor 5. nap}

\parttitle[24, 1-4. 10-18. 24b-25]{Első olvasmány}
Sámuel második könyvéből

\begin{lectio}{Népszámlálás és oltárépítés}
\noindent Azokban a napokban: Újra fellángolt az Úr haragja Izrael ellen, úgyhogy ismét felingerelte ellenük Dávidot és azt mondta neki: „Menj s vedd számba Izraelt és Júdát!” A király ezért megparancsolta Joábnak és a sereg vezéreinek: „Járjátok sorra Izrael törzseit Dántól egészen Beersebáig, és vegyétek számba a népet, hogy megtudjam, mennyi a nép!” Joáb azonban azt mondta a királynak: „Az Úr, a te Istened, bármily népes is a nép, tegye százszor népesebbé, és uram, királyom ezt lássa meg a saját szemével! De miért kíván uram és királyom ilyesmit?” Ám a király parancsa sürgette Joábot, valamint a sereg vezéreit, így hát Joáb, valamint a sereg vezérei távoztak a királytól, hogy számba vegyék Izrael népét.

Miután számba vétette a népet, Dávidnak megdobbant a szíve, így szólt hát az Úrhoz: „Nagy bűnt követtem el! De bocsásd meg, Uram, szolgádnak ezt a bűnt, mert ostobaságot csináltam.”

Amikor Dávid reggel fölkelt, az Úr ezt a szózatot intézte Dávid látóemberéhez, Gád prófétához: „Menj és add tudtára Dávidnak: Ezt mondja az Úr: Hármat mondok, aztán amit választasz, azt teszem veled!” Gád elment Dávidhoz, és tudtára adta. Így szólt: „Jöjjön-e három évig tartó éhínség büntetésül az országra? Vagy azt akarod, hogy három hónapon át menekülj ellenséged elől, ő meg üldözzön? Vagy legyen három napig pestis országodban? Gondolkozz rajta, és mondd meg, milyen választ vigyek annak, aki küldött!” Erre Dávid azt felelte Gádnak: „Nagy szorultságban vagyok. Inkább kerüljünk az Úr kezébe, hiszen nagy az irgalma! De emberek kezétől nem szeretnék elveszni.”

Akkor az Úr pestist küldött Izraelre, reggeltől a megszabott időre. A csapás lesújtott a népre, és hetvenezren meghaltak a népből Dántól egészen Beersebáig. Amikor azonban az angyal Jeruzsálem felé is kinyújtotta a kezét, hogy elpusztítsa, az Úr megbánta ezt a csapást, és megparancsolta az angyalnak, aki lesújtott a népre: „Most aztán elég! Húzd vissza a kezed!” Az Úr angyala éppen a jebuzita Arauna szérűjénél volt.

Amikor Dávid látta az angyalt, aki lesújtott a népre, így könyörgött az Úrhoz: „Én vétkeztem, én követtem el gonoszságot! Ezek, a nyáj, mit vétettek? Ezért hát forduljon kezed ellenem és házam ellen!” Azon a napon Gád elment Dávidhoz, és azt mondta neki: „Menj föl, és a jebuzita Arauna szérűjénél építs oltárt az Úrnak.” Így hát Dávid megvette a szérűt és a marhákat 50 ezüst sékelért. Aztán Dávid oltárt emelt az Úrnak, és égőáldozatot, valamint közösségi áldozatot mutatott be rajta. Erre az Úr irgalommal fordult az ország felé, és elmúlt Izraelről a csapás.
\end{lectio}

\parttitle[Vö. Judit 9, 18; 1 Krón 21, 15; Sám 24, 17]{Válaszos ének}

\versicle{%
    Emlékezzél, Uram, szövetségedre, és mondd pusztító \underline{a}ngyal\underline{o}dnak: † Vond
    vissza \underline{a} k\underline{e}zed, * Hogy ne pusztuljon el a föld, és ne veszítsd el az
    él\underline{ő}lény\underline{e}ket.
}{%
    Én vétkeztem, én követtem el gonoszs\underline{á}got! † Ezek, a nyáj mit vét\underline{e}ttek? * Ezért
    hát kérlek, forduljon el haragod népedt\underline{ő}l, \underline{U}ram.
}

\parttitle{Második olvasmány}
„A tizenkét apostol tanításá”-nak nevezett őskeresztény írásból

{\centering\redtext{(Nn. 9, 1 – 10, 6; 14, 1-3: Funk 2, 19-22. 26)}}

\begin{lectio}{Az Eukarisztia}
\noindent Így adjatok hálát: Először a kehely fölött mondjátok: „Hálát adunk neked, Atyánk, szolgádnak, Dávidnak szent szőlőtövéért, amelyet nekünk a te Fiad, Jézus által nyilatkoztattál ki; dicsőség neked mindörökké!”

A kenyértörés után pedig ekképpen: „Hálát adunk neked, Atyánk, azért az életért és tanításért, amelyet nekünk a te Fiad, Jézus által nyilatkoztattál ki; dicsőség neked mindörökké! Miként ez a kenyér szét volt hintve a földeken, és egybegyűjtve egy kenyérré vált, úgy gyűljön össze Egyházad is a föld határairól a te országodba; mert tiéd a hatalom és a dicsőség Jézus Krisztus által mindörökké!”

Eukarisztiátokból pedig senki más ne egyék és igyék, csak az, aki meg van keresztelve az Úr nevében; erről mondotta ugyanis az Úr: \textit{Ne adjátok a szent dolgokat a kutyáknak!} (Mt 7, 6)

Miután pedig magatokhoz vettétek, így adjatok hálát: „Hálát adunk neked, szent Atyánk, szent nevedért, amelynek szívünkben trónust készítettél; hálát adunk a tanításért, a hitért és a halhatatlanságért, amelyet a te Fiad, Jézus által kinyilvánítottál nekünk; dicsőség neked mindörökké!

Mindenható Urunk, te mindent nevedért teremtettél, és ételt meg italt adtál az embernek, hogy éljenek vele és neked hálát adjanak. Nekünk pedig lelki ételt és italt adtál és örök életet Fiad által. Mindenekelőtt hálát adunk neked, mert hatalmas vagy; dicsőség neked mindörökké!

Emlékezzél meg, Urunk, Egyházadról! Védd meg minden bajtól, és tedd tökéletessé szeretetedben! Megszentelve gyűjtsd össze a négy égtáj felől országodba, amelyet számára készítettél; mert tiéd a hatalom és a dicsőség mindörökké!

Jöjjön el a te kegyelmed, és szűnjék meg a bűnös világ! Hozsanna Dávid Istenének! Ha valaki szent, járuljon ide; ha nem az, tartson bűnbánatot; Marana tha. (Jöjj el, Urunk!) Ámen.”

Az Úr napján összegyülekezve végezzétek a kenyértörést, és adjatok hálát, miután bűneiteket megvallottátok, hogy tiszta legyen áldozatotok. Akinek pedig viszálya van testvérével, ne vegyen részt a gyülekezetben, amíg ki nem békül, hogy csorbát ne szenvedjen áldozatotok szentsége. Mert ezt mondta az Úr: \textit{Tiszta áldozatot mutassanak be nekem minden helyen és minden időben, mert én vagyok a nagy Király, úgymond az Úr, és nevem csodálatos a nemzetek között} (vö. Mal 1, 11).
\end{lectio}

\parttitle[1 Kor 10, 16-17]{Válaszos ének}

\versicle{%
    Az áldás kelyhe, amely\underline{e}t meg\underline{á}ldunk, † nemde a Krisztus vérében való
    rész\underline{e}s\underline{e}dés? * S a kenyér, amelyet megtörünk, nemde a Krisztus testében
    való r\underline{é}szesedés?
}{%
    M\underline{i} ugyanis sokan egy kenyér és egy test v\underline{a}gyunk, * mivel mindnyájan egy
    kenyérből rész\underline{e}s\underline{ü}lünk.
}

\section*{Csütörtök, 14. hét - 1. nap}

\parttitle[22, 5-19]{Első olvasmány}
A Krónikák első könyvéből

\begin{lectio}{Dávid előkészületei a templomépítéshez}
\noindent Azokban a napokban Dávid így gondolkozott: „A fiam, Salamon fiatal és gyenge, a háznak azonban, amelyet az Úrnak építünk, egészen nagynak kell lennie, hogy dicsőséget és hírnevet szerezzen minden nép előtt. Ezért előkészületeket kell tennem.” Dávid halála előtt valóban megtette az előkészületeket, azután hívatta a fiát, Salamont, és megparancsolta neki, hogy építsen házat az Úrnak, Izrael Istenének. Dávid ezt mondta Salamonnak: „Fiam, nekem magamnak volt szándékomban, hogy házat építsek az Úr, az én Istenem nevének, de az Úr szózatot intézett hozzám: »Sok vért ontottál és nagy háborúkat vívtál, ne építs hát házat a nevemnek, mert sok vért ontottál színem előtt a földre. Nézd, született egy fiad, ő a béke embere lesz. Békét adok neki körös-körül minden ellenségétől. Salamon lesz a neve. Az ő napjaiban békét és nyugalmat adok Izraelnek, ő építsen házat a nevemnek. A fiam lesz, én meg atyja leszek. Örökre megerősítem Izrael fölötti uralmának trónját.«

Nos hát, fiam, legyen veled az Úr, sikerüljön házat építened az Úrnak, a te Istenednek, ahogy előre megmondta rólad. Az Úr adjon neked bölcsességet és okosságot, és nyilvánítsa ki neked Izraelre vonatkozó akaratát. Te pedig tartsd meg az Úrnak, a te Istenednek a törvényét. Szerencsés leszel, ha igyekszel megtartani a törvényt és a parancsokat, amelyeket az Úr Izrael számára Mózesnek adott. Légy erős és bátor! Ne félj és ne rettegj! Nézd, saját szegénységemből félre tudtam tenni az Úr háza számára 100.000 arany talentumot, egy millió ezüst talentumot, s annyi vasat és bronzot, hogy meg sem lehet mérni. Fát és követ is szereztem, de még te is tehetsz hozzájuk. Mesteremberek is nagy számban állnak rendelkezésedre: kőművesek, ácsok és mindenféle művészetben jártas számtalan ember, akik meg tudják munkálni az aranyat, az ezüstöt, a bronzot és a vasat. Rajta hát, láss hozzá, és legyen veled az Úr.”

Dávid megparancsolta Izrael minden tisztviselőjének, hogy segítse Salamont: „Az Isten, az Úr nemde veletek volt, és nyugalmat adott nektek körös-körül. Az ország népét a kezembe adta, s a föld meghódolt az Úr és az ő népe előtt. Most tehát egész szívvel és lélekkel keressétek az Urat, a ti Isteneteket! Rajta, építsetek szentélyt Istennek, az Úrnak, hogy a szövetség ládáját és az Isten szent dolgait bevigyétek abba a házba, amelyet az Úr nevének építetek.”
\end{lectio}

\parttitle[1 Krón 22, 19; Zsolt 131, 7; Iz 56, 7]{Válaszos ének}

\versicle{%
    R\underline{a}jta, † egész szívvel és lélekkel építsetek szentélyt Istennek, \underline{a}z \underline{Ú}rnak; * Induljunk hát sátorába, boruljunk le lábának zsám\underline{o}lya \underline{e}lőtt!
}{%
    \underline{E}zt mondja \underline{a}z Úr: * házam minden nép számára az imádság h\underline{á}z\underline{a} lesz. 
}

\parttitle{Második olvasmány}
Szent Ambrus püspöknek a 118. zsoltárhoz írt magyarázatából

{\centering\redtext{(Nn. 12. 13-14: CSEL 62, 258-259)}}

\begin{lectio}{Isten temploma szent, és ti vagytok az}
\noindent \textit{Én és az Atya hozzá megyünk, és benne fogunk lakni} (Jn 14, 23). Nyisd hát ki ajtódat az érkező előtt, tárd elé a lelked, mutasd meg neki szíved mélyét, hogy lássa nagy őszinteségedet, békességed gazdagságát és kegyelme iránti szeretetedet. Tárd ki szívedet, siess az örök fény napja elé, \textit{amely minden embert megvilágít} (Jn 1, 9). Ez az igaz fény mindenkit megvilágít; ha pedig valaki bezárja előtte ablakait, akkor megfosztja magát az örök világosságtól. Kizárod Krisztust is, ha becsukod lelked ajtaját. Jóllehet be tudna menni akkor is, de nem akar erőszakkal betörni, sem pedig kényszeríteni azt, aki nem kívánja, hogy betérjen hozzá.

A Szűz méhéből született, és beragyogja az egész földet, hogy mindenkit megvilágítson. Befogadják azok, akik arra az örök fénysugárra vágyakoznak, amelyet semmiféle éjszaka nem szakít meg. A napot, amelyet mindennap látunk, sötét éjszaka követi; az igazság Napja azonban soha le nem nyugszik, mert a bölcsességet nem követheti rosszaság.

Tehát boldog az, akinek ajtaján Krisztus zörget. A mi ajtónk a hit, amely ha erős, megőrzi az egész házat. Ezen az ajtón át lép be Krisztus. Ezért mondja az Egyház is az Énekek énekében: \textit{Testvérem hangja zörget az ajtómon} (Én 5, 2). Halld meg, ki zörget, halld meg azt, aki be akar jönni: \textit{Nyisd ki az ajtót, húgocskám, kedvesem, galambom, gyöngyöm, mert a fejemet belepte a harmat, az éjszaka párája a fürtjeimet} (Én 5, 2).

Fontold meg: mikor kopogtat az isteni Ige ilyen nagyon az ajtódon, mikor lepte be a fejét az éjszaka harmata? A szorongattatásban és a kísértésben levőkét látogatja meg ilyen jóságosan, nehogy legyőzze őket ez a megpróbáltatás. A feje akkor lesz harmatos és párás, amikor teste szenved. Tehát ekkor virrasztani kell, nehogy amikor jön a Jegyes, zárt ajtóra találva továbbmenjen. Ha ugyanis alszol és nem virraszt a szíved, eltávozik, még mielőtt kopogtatott volna; ha viszont virraszt a szíved, kopogtat, és kéri, hogy nyiss neki ajtót.

Lelkünknek tehát van ajtaja, de kapui is vannak, ezekről mondja a Szentírás: \textit{Emeljétek föl fejeteket, kapuk, táruljatok föl örök kapuk, a dicsőség királya bevonul} (Zsolt 23, 7). Ha tehát ki akarod tárni hited kapuit, akkor belép hozzád a dicsőség Királya, és elhozza neked az ő szenvedéséből fakadó megdicsőülést. Az igazságnak is van kapuja. Hiszen erről is olvasunk a Szentírásban, amikor az Úr Jézus a próféta által ezt mondja: \textit{Tárjátok ki az igazság kapuit} (vö. Zsolt 117, 19).

Van tehát lélek, amelyiknek ajtaja van, és van, amelyiknek kapui vannak. Krisztus eljön ehhez az ajtóhoz és kopogtat, kopogtat a kapunál is. Nyiss tehát neki ajtót: be akar jönni, virrasztva akarja találni Jegyesét.
\end{lectio}

\parttitle[Jel 3, 20; Mt 24, 46]{Válaszos ének}

\versicle{%
    Nézd, az ajtóban állok, \underline{é}s kop\underline{o}gok. † Aki meghallja szavam, és \underline{a}jt\underline{ó}t nyit, * Bemegyek hozzá, vele eszem, ő m\underline{e}g énv\underline{e}lem.
}{%
    Boldog a sz\underline{o}lga, * ha ura hazatérve ilyen munkában t\underline{a}l\underline{á}lja.
}

\section*{Péntek, 14. hét - 2. nap}

\parttitle[1, 11-35; 2, 10-12]{Első olvasmány}
A Krónikák első könyvéből

\begin{lectio}{Dávid Salamont jelöli ki utódjául}
\noindent Azokban a napokban Nátán próféta így szólt Batsebához, Salamon anyjához: „Bizonyára hallottad, hogy Adonija, Haggit fia király lett anélkül, hogy Urunk, Dávid tudna róla. Nos, hadd adjak neked tanácsot, hogyan mentheted meg az életedet és fiad, Salamon életét. Fogd magad, menj be a királyhoz s mondd neki: Uram és királyom, megesküdtél rá szolgálódnak: Fiad, Salamon lesz utánam a király, ő fog trónomra ülni. Miért lett hát mégis Adonija a király? Mialatt még beszélsz a királlyal, magam is bemegyek és megerősítem szavaidat.”

Batseba tehát bement a királyhoz a szobába. A király már nagyon öreg volt, s a sunemi Abiság viselte gondját. Amikor Batseba meghajolt és leborult a király előtt, a király megkérdezte: „Mit kívánsz?” Így felelt neki: „Uram, megesküdtél az Úrra, a te Istenedre szolgálódnak: Fiad, Salamon lesz utánam a király, ő fog trónomra ülni. Most mégis Adonija lett a király, s te, uram és királyom, mit sem tudsz róla. Rengeteg marhát, hizlalt borjút és juhot levágatott, s mind meghívta a király fiait, Ebjatár papot és Joábot, a sereg vezérét. Szolgádat, Salamont azonban nem hívta meg. Egész Izraelnek rajtad a szeme, uram és királyom, add tudtukra, ki legyen, aki majd utódként uram és királyom trónjára ül. Különben megtörténhet, hogy ha uram és királyom majd atyáihoz megtér, engem és fiamat, Salamont bűnösnek nyilvánítanak.”

Mialatt még beszélt a királlyal, megérkezett Nátán próféta. Jelentették a királynak: „Nátán próféta van itt.” Erre az a király elé lépett és arcra borult a földön a király előtt. S így beszélt Nátán: „Uram és királyom, bizonyára magad jelentetted ki: Adonija legyen utánam a király, ő üljön trónomra. Mert ma elment, és rengeteg marhát, hízott borjút és juhot levágatott. Meghívta mind a király fiait s a sereg vezérét meg Ebjatár papot is. S most ott esznek-isznak nála, és köszöntgetik: »Éljen Adonija király!« Engemet azonban, a te szolgádat, aztán Cádok papot, Jojada fiát, Bénáját és szolgádat, Salamont nem hívta meg. Ha ez a dolog mégis uramtól és királyomtól indult volna ki, akkor nem adtad tudtára szolgáidnak, ki üljön utódnak uram és királyom trónjára.”

Ekkor Dávid király vette át a szót és így rendelkezett: „Hívjátok ide Batsebát!” Az bement a királyhoz, és a király elé járult. A király megesküdött: – Úgy igaz, ahogy az Úr él, aki minden szorongatásomból kiszabadított, még ma úgy lesz, amint megesküdtem neked az Úrra, Izrael Istenére: fiad, Salamon lesz utánam a király, ő ül helyettem a trónra.” Batseba földig hajolt, leborult a király előtt és felkiáltott: „Örökké éljen az én királyom, Dávid!”

Aztán Dávid király megparancsolta: „Hívjátok ide Cádok papot, Nátán prófétát és Jojada fiát, Bénáját.” Azok megjelentek a király előtt, s a király ezt mondta nekik: „Vigyétek magatokkal uratok szolgáit, azután ültessétek fiamat, Salamont a tulajdon öszvéremre, és vezessétek le Gichonba. Ott Cádok pap és Nátán próféta kenje föl Izrael királyává. Fúvassátok meg a harsonát, és kiáltsátok: »Éljen Salamon király!« Azután az ő nyomában haladva vonuljatok ide vissza, ő pedig jöjjön, és foglalja el trónomat, legyen helyettem a király. Őt választottam ki arra, hogy Izrael és Júda fejedelme legyen.”

Aztán Dávid megtért atyáihoz; Dávid városában temették el. Az az idő, amíg Dávid királyként uralkodott Izrael fölött, negyven esztendőt tett ki: Hebronban hét évig volt király, Jeruzsálemben pedig harminchárom esztendeig. Salamon elfoglalta atyja trónját, és megerősítette királyságát.
\end{lectio}

\parttitle[Én 3, 11; Zsolt 71, 1. 2]{Válaszos ének}

\versicle{%
    Sion leányai, gyert\underline{e}k, nézz\underline{é}tek: † Salamon király, fején a korona, amellyel szülőanyja megkor\underline{o}n\underline{á}zta * Szíve öröm\underline{é}nek n\underline{a}pján!
}{%
    Istenem, ítéletedet add a kir\underline{á}lynak, * hogy méltányosan ítélkezzék szegény\underline{ei}den. 
}

\parttitle{Második olvasmány}
Szent I. Kelemen pápának a korintusiakhoz írt leveléből

{\centering\redtext{(Nn. 50, 1 – 51, 3; 55, 1-4: Funk 1, 125-127. 129)}}

\begin{lectio}{Boldogok leszünk, ha az Úr akaratát szeretetből megtesszük}
\noindent Látjátok, testvéreim, milyen nagy és csodálatos dolog a szeretet. Tökéletes voltát nem lehet szóval elmondani. Ki alkalmas rá, hogy a szeretetben éljen? Csak azok, akiket Isten méltóvá akart tenni erre. Imádkozzunk tehát és kérjük irgalmát, hogy minden bűnös hajlamtól mentesen megmaradjunk a szeretetben. Ádámtól mindmáig elmúltak az összes nemzedékek, de akik Isten kegyelméből az ő szeretetében haltak meg, elnyerték a szentek helyét, és amikor megjelenik Krisztus országa, megdicsőülnek. Mert írva van: \textit{Menj be kamrádba rövid időre, míg haragos indulatom el nem múlik, és visszagondolok a jó napra, és feltámasztalak benneteket sírjaitokból} (vö. Iz 26, 20).

Boldogok vagyunk, szeretteim, ha az Úr parancsait egyetértő szeretettel teljesítjük, hogy a szeretet által bűneink bocsánatát elnyerjük. Mert meg van írva: \textit{Boldogok, akiknek gonoszsága megbocsátva, és akinek bűnét eltörölték. Boldog az, akinek az Úr nem számítja be vétkét, és akinek lelke hamisság nélkül való} (Zsolt 31, 1). Ez a boldogság azokra vonatkozik, akiket Isten a mi Urunk, Jézus Krisztus által kiválasztott. Neki legyen dicsőség mindörökkön-örökké. Ámen.

Tehát ellenségünk kísértései által elcsábítva bármit vétettünk és bármi rosszat tettünk, könyörögjünk bűneink bocsánatáért. Azok pedig, akik a lázadásban és a pártoskodásban vezetők voltak, tekintsenek a közös reményre. Mert akik istenfélelemben és szeretetben töltik életüket, azok a kínzatást inkább maguknak kívánjak, mint embertársaiknak, és azt szeretnék, hogy inkább őket érje a feddés, mintsem hogy károsodás érje a szépen és igazságosan megőrzött egyetértést. Jobb ugyanis az embernek megvallania vétket, mint megkeményítenie a szívét.

Ki van hát köztetek nemes lelkű, irgalmas, akiben teljes a szeretet? Mondja azt: „Ha miattam van a zendülés, a széthúzás, a szakadás, inkább elmegyek, odébbállok, ahová akarjátok, és azt teszem, amit a gyülekezet parancsolt, csak Krisztus nyája éljen békében a vezetőkül rendelt presbiterekkel.” Aki így cselekszik, Krisztusnál nagy dicsőséget szerez magának, és bárhol befogadják őt, \textit{mert az Úré a föld, és mind, ami betölti} (Zsolt 23, 1). Ezt teszik és fogják is tenni azok, akik az isteni életet élik, ezt ugyanis sohasem fogják megbánni.
\end{lectio}

\parttitle[1 Jn 4, 21; Mt 22, 40]{Válaszos ének}

\versicle{%
    Ezt \underline{a} par\underline{a}ncsot † kaptuk \underline{I}st\underline{e}ntől: * Aki Istent szereti, szeresse t\underline{e}stvér\underline{é}t is.
}{%
    Ezen a két parancson al\underline{a}pszik * az egész törvény és a pr\underline{ó}f\underline{é}ták.
}

\section*{Szent Benedek apát, Európa fővédőszentje - 3. nap}

\parttitle[3, 7 – 4, 1. 4-9]{Első olvasmány}
Szent Pál apostolnak a filippiekhez írt leveléből

\begin{lectio}{Örüljetek az Úrban szüntelenül!}
\noindent Testvéreim! Amit akkor előnynek tartottam, azt Krisztusért hátránynak tekintem. Sőt, Uramnak, Krisztus Jézusnak fönséges ismeretéhez mérten mindent hátránynak tartok. Érte mindent elvetettem, sőt szemétnek tekintettem, csakhogy Krisztust elnyerhessem, és hozzá tartozzam. Hiszen nem a törvény útján váltam igazzá, hanem a Jézus Krisztusba vetett hit révén. Isten ugyanis a hit által tett igazzá, hogy megismerjem őt és feltámadásának erejét, de a szenvedésben is vállaljam vele a közösséget, így hozzá hasonulok a halálban, hogy ezáltal eljussak a halálból a feltámadásra is.

Nem mintha már elértem volna, vagy már célba értem volna, de futok utána, hogy magamhoz ragadjam, ahogy Krisztus is magával ragadott engem. Testvérek, nem gondolom, hogy máris magamhoz ragadtam, de azt igen, hogy elfelejtem, ami mögöttem van, és nekilendülök annak, ami előttem van. Futok a kitűzött cél felé, az égi hivatás jutalmáért, amelyre Isten meghívott Krisztus (Jézusban).

Mi, tökéletesek, gondolkodjunk így! Ha valamiben még másképpen éreztek, Isten majd megvilágosít benneteket, de amit már elértünk, abban tartsunk ki.

Testvérek, kövessétek a példámat mindnyájan! Nézzétek azokat, akik úgy élnek, ahogy példámon látjátok. Hiszen – mint már többször mondtam, most meg könnyek közt mondom – sokan úgy élnek, mint Krisztus keresztjének ellenségei. Végük a pusztulás, istenük a hasuk, azzal dicsekszenek, ami gyalázatukra válik, s eszüket földi dolgokon járatják.

A mi hazánk azonban a mennyben van. Onnan várjuk az Üdvözítőt is, Urunkat, Jézus Krisztust. Ő azzal az erővel, amellyel mindent hatalma alá vethet, átalakítja gyarló testünket, és hasonlóvá teszi megdicsőült testéhez.

Ezért, szívből szeretett testvéreim, örömöm és koronám: így álljatok helyt az Úrban, szeretteim! Örüljetek az Úrban szüntelenül! Újra csak azt mondom, örüljetek. Jóságos emberségeteket ismerje meg mindenki! Az Úr közel van. Ne aggódjatok semmi miatt, hanem minden imádságotokban és könyörgésetekben terjesszétek kéréseteket az Úr elé, hálaadásotokkal együtt. Akkor Isten békéje, amely minden értelmet meghalad, megőrzi szíveteket és értelmeteket Krisztus Jézusban.

Egyébként, testvéreim, arra irányuljanak gondolataitok, ami igaz, tisztességes, igazságos, ami ártatlan, kedves, dicséretre méltó, ami erényes és magasztos.

Amit tanultatok és elfogadtatok, amit hallottatok, és példámon láttatok, azt váltsátok tettekre, s veletek lesz a béke Istene.
\end{lectio}

\parttitle[Lk 12, 35-36; Mt 24, 42]{Válaszos ének}

\versicle{%
    Csípőtök legyen f\underline{e}löv\underline{e}zve, † és kezetekben égjen a lámp\underline{á}s\underline{o}tok. * Hasonlítsatok azokhoz az emberekhez, akik Urukra várnak, hogy megérkezzék a m\underline{e}nyegz\underline{ő}ről.
}{%
    Legyetek hát éb\underline{e}rek, † mert nem tudj\underline{á}tok, * melyik órában jön el \underline{U}r\underline{a}tok.
}

\parttitle{Második olvasmány}
Szent Benedek apát Regulájából

{\centering\redtext{(Prologus, 4-22; cap. 72, 1-12: CSEL 75, 2-5. 162-163)}}

\begin{lectio}{Krisztusnak semmit elébe ne tegyenek!}
\noindent Először is: bármi jóba kezdesz, igen állhatatosan imádsággal kérjed, hogy ő vigye azt végbe, hogy ő, aki bennünket már fiai sorába méltóztatott számítani, ne legyen kénytelen valaha is szomorkodni rossz tetteink miatt. Úgy kell őt a bennünk levő adományaival mindenkor szolgálnunk, hogy mint haragvó atya örökségéből ki ne tagadjon minket, fiait, de úgy se, mint félelmetes úr, bűneinktől felingerelve, gonosz szolgák gyanánt át ne adja az örök büntetésre azokat, akik nem akarták őt követni a dicsőségbe.

Keljünk föl tehát végre valahára a Szentírás serkentő szavára: \textit{Itt az óra, hogy felébredjünk az álomból} (Róm 13, 11). Nyissuk meg szemünket a megistenítő fénynek, és megdöbbent füllel halljuk, mire int bennünket a mindennap felénk kiáltó isteni szózat: \textit{Ma, ha az Úr szavát halljátok, meg ne keményítsétek szíveteket!} (Zsolt 94, 8) És ismét: \textit{Akinek füle van a hallásra, hallja meg, mit mond a Lélek az egyházaknak!} (Jel 2, 7) És mit mond? \textit{Jöjjetek, fiaim, hallgassatok rám! Az Úr félelmére tanítalak titeket} (Zsolt 33, 12). \textit{Siessetek, míg tiétek az élet világossága, hogy a halál sötétsége meg ne lepjen benneteket!} (Jn 12, 35)

Ezeket kiáltja oda az Úr a népsokaságnak, amikor keresi munkását, majd így folytatja: \textit{Ki az az ember, aki életet óhajt, és jó napokat kíván látni?} (Zsolt 33, 13) Ha ennek hallatára azt feleled: „Én”, akkor Isten ezt mondja neked: Ha igazi és örök életet akarsz, \textit{tiltsd el nyelvedet a gonosztól, és ajkad ne szóljon csalárdságot, fordulj el a rossztól, és tedd a jót! Keresd a békét, és járj utána!} (Zsolt 33, 14-15) Ha ezt megteszitek, szemem rajtatok lesz, és fülem meghallja könyörgéseteket. És még mielőtt segítségül hívnátok, azt mondom nektek: \textit{Íme, itt vagyok} (Iz 52, 6).

Mi lehetne édesebb számunkra, szeretett testvéreim, az Úrnak e minket hívó szavánál? Íme, így mutatja meg jóságosan az Úr az élet útját.

Övezzük fel tehát derekunkat hittel és a jó cselekedetek gyakorlásával, és az evangélium vezetésével járjuk az ő útjait, hogy méltók legyünk meglátni azt, \textit{aki országába hívott minket} (1 Tessz 2, 13). Ám ha országának sátorában akarunk lakni, hacsak jó cselekedetekkel nem igyekszünk, soha oda el nem jutunk.

Amint van keserűségből fakadó rossz buzgóság, amely elszakít Istentől, és a pokolra vezet, úgy van jó buzgóság is, amely a vétkektől választ el, és Istenhez és az örök életre vezet. Ezt a buzgóságot tehát a legizzóbb szeretettel gyakorolják a szerzetesek, azaz: \textit{Egymást a tiszteletadásban előzzék meg} (Róm 12, 10). Egymásnak mind testi, mind lelki fogyatkozásait a legtürelmesebben viseljék el. Egymásnak versengve engedelmeskedjenek. Senki se kövesse azt, amit maga számára hasznosnak ítél, hanem ami másnak az. Egymást tisztán, testvéri szeretettel szeressék. Istent szeretve féljék. Apátjukat őszinte és alázatos szeretettel szeressék. Krisztusnak semmit elébe ne tegyenek, aki bennünket, mindnyájunkat az örök életre vezessen.
\end{lectio}

\parttitle{Válaszos ének}

\versicle{%
    Szent Benedek elhagyta szülei ház\underline{á}t, vagy\underline{o}nát, † és egyedül Istennek akarván tetszeni, a szent szerzetesi életet vál\underline{a}szt\underline{o}tta. * A magányban élt, Isten sz\underline{í}ne \underline{e}lőtt.
}{%
    Visszavonult a vil\underline{á}gtól, † hogy ne tudja, amit már meg\underline{i}smert, * és megtapasztalja, amit még n\underline{e}m t\underline{u}dott.
}

\section*{Évközi 15. vasárnap - 4. nap}

\parttitle[16, 29 – 17, 16]{Első olvasmány}
A királyok első könyvéből

\begin{lectio}{Illés próféta fellépése Acháb izraeli király idejében}
\noindent Omri fia, Acháb – Júda királyának, Azának 38. esztendejében lett király Izraelben. Omri fia, Acháb huszonkét évig uralkodott Izrael fölött Szamariában. S Acháb, Omri fia azt tette, ami gonosznak számít az Úr szemében, még sokkal inkább, mint elődei.

Nem volt neki elég, hogy Nebat fia, Jerobeám nyomán járt a bűn útján, hanem még a szidoniak királyának, Etbaalnak a lányát, Izebelt is elvette feleségül, aztán elment Baalnak szolgálni, őt imádni. Baal templomában, amelyet Szamariában épített, oltárt emelt Baalnak. Visszaállította a szent fákat, sőt hogy ingerelje az Urat, Izrael Istenét, még több gonoszságot művelt, mint Izrael összes többi királya, aki előtte volt. Az ő idejében a beteli Hiel fölépítette újra Jerikót. Elsőszülötte, Abirám árán vetette meg alapjait, s legkisebb fia, Segub árán csinált rá kapukat, az Úr szava szerint, amelyet Nun fia, Józsue által hallatott.

A gileádi Tisbéből való Illés, a tisbei így szólt Achábhoz: „Amint igaz, hogy él az Úr, Izrael Istene, akinek a szolgálatában állok: ezekben az években ne hulljon se harmat, se eső, csak az én szavamra.” Azt mondta neki az Úr: „Menj el innen, vedd utadat kelet felé, és rejtőzz el a Kerit pataknál, amely a Jordántól keletre folyik. Igyál a patakból, a hollóknak pedig megparancsolom, hogy gondoskodjanak rólad.” Tehát elment és úgy tett, amint az Úr parancsolta. Elment és a Kerit pataknál maradt, amely a Jordántól keletre folyik. S a hollók vittek neki reggel kenyeret, este húst, a patakból meg ivott. Egy bizonyos idő elteltével történt, hogy a patak kiszáradt, mert nem esett eső az országban.

Akkor az Úr így szólt hozzá: „Kelj útra és menj el a Szidonhoz tartozó Careftába, és maradj ott. Nézd, ott megparancsoltam egy özvegyasszonynak, hogy gondoskodjék rólad.” Erre útra kelt és elment Careftába. Amikor a város kapujához ért, épp ott volt egy özvegyasszony – rőzsét szedegetett. Megszólította és azt mondta neki: „Hozz nekem egy kis vizet korsóban, hadd igyam!” Amikor elment hozni, utána szólt: „Hozz egy harapás kenyeret is!” Azt felelte: „Amint igaz, hogy a te Istened él: nincs sütve semmim, csak egy marék lisztem van a szakajtóban meg egy kis olajam a korsóban. Épp azon vagyok, hogy rőzsét szedjek, aztán megyek és elkészítem magamnak és fiamnak. Megesszük, aztán meghalunk.” Illés azonban így válaszolt neki: „Ne félj! Menj s tedd, amit mondtál; csak előbb csinálj belőle egy kis lángost, aztán hozd ki nekem; magadnak és fiadnak csak utána készíts. Mert azt mondja az Úr, Izrael Istene: A szakajtó ne ürüljön ki, a korsó ne apadjon el addig, amíg az Úr esőt nem hullat a földre.”

Elment hát és úgy tett, amint Illés mondta. S volt mit ennie, neki is, fiának is. A szakajtó nem ürült ki, és a korsó nem apadt ki az Úr szava szerint, amelyet Illés által hallatott.
\end{lectio}

\parttitle[Jak 5, 17. 18; Sir 48, 1. 3]{Válaszos ének}

\versicle{%
    Illés esedezve imádkozott, hogy ne ess\underline{é}k az \underline{e}ső, † s nem \underline{i}s \underline{e}sett. * Majd ismét imádkozott, és esőt \underline{a}dott \underline{a}z ég.
}{%
    Mint a tűzvész, Illés prófét\underline{a} jött, † szava lángolt, mint az égő f\underline{á}klya, * az Úr szavával elzárta \underline{a}z \underline{e}get.
}

\parttitle{Második olvasmány}
Kezdődik Szent Ambrus püspök „A misztériumok” című értekezése

{\centering\redtext{(Nn. 1-7: SCh 25 bis, 156-158)}}

\begin{lectio}{Oktatás a keresztség előtti szertartásokról}
\noindent Naponként szóltam azelőtt hozzátok erkölcsi kérdésekről, amikor a pátriárkák történetét vagy a Példabeszédek könyvének parancsait olvastuk. Az volt a célom, hogy ezek ismeretével és elsajátításával hozzászokjatok az ősatyák példáján való járáshoz; hogy az ő útjukat követve szót fogadjatok az isteni kinyilatkoztatásnak; hogy újjászületve a keresztség által olyan legyen életetek, amilyen a megkereszteltekhez méltó.

Most az idő arra int, hogy a misztériumokról beszéljek, és a szent jeleknek magyarázatát adjam. Ha ezt a keresztség előtt szándékoztam volna közölni olyanokkal, akik még nem voltak beavatva, inkább árulónak, mint magyarázónak tűntem volna fel. Aztán meg így – előkészítés nélkül – jobban belétek hatol a misztériumok fénye, mintha már valami tanítás megelőzte volna.

Nyissátok meg hát fületeket, és élvezzétek az örök életnek a szentségek által rátok lehelt jó illatát. Ezt mutattuk be nektek jelképesen, midőn a fül megnyitásának szent szertartását végeztük és mondottuk: \textit{Effeta}, azaz \textit{nyílj meg} (Mk 7, 34), hogy akik a kegyelem forrásához jöttetek, megértsétek a feladott kérdéseket és emlékezetben tartsátok a feleleteket. Ezt a misztériumot maga Krisztus jelezte az evangéliumban, amikor a süketnémát meggyógyította.

Ezután megnyílt számodra a legszentebb szentély, beléptél az újjászületés szentélyébe. Most idézd emlékezetedbe, mit kérdeztek tőled, és tartsd emlékezetedben, hogy mit feleltél. Ellene mondottál a sátánnak és cselekedeteinek, a világnak, a világi gyönyöröknek és élvezeteknek, ígéretedet nem a holtak sírja őrzi, hanem az élet könyve.

Láttál ott levitát, láttál áldozópapot és láttad a püspököt is. Ne tekintsd testük alakját, hanem szolgálatuk kegyelmét. Angyalok színe előtt szóltál, miként írva van: \textit{A pap ajkának kell a tudást őriznie, és az ő szájából várják a tanítást, mert ő a Seregek Urának követe} (Mal 2, 7). Nincs helye a tévedésnek, nincs helye a tagadásnak: Isten követe az, aki Krisztus országát és az örök életet hirdeti. Ne külseje szerint értékeld, hanem küldetése szerint. Azt nézd, amit neked juttatott; mérlegeld, mit tett, és ismerd fel méltóságát.

Beléptél tehát, hogy ellenséged elé állj; elhatároztad, hogy nyíltan ellene mondasz neki; kelet felé fordulsz: aki ugyanis ellene fordult az ördögnek, az Krisztushoz fordul, ránéz bizakodó tekintettel.
\end{lectio}

\parttitle[Tit 3, 3. 5; Ef 2, 3]{Válaszos ének}

\versicle{%
    Egykor mi magunk is balgák, engedetlenek és tév\underline{e}lygők v\underline{o}ltunk; † gonoszságunk, irigységünk miatt utálatra méltók voltunk, és gyűlölt\underline{ü}k \underline{e}gymást. * De Isten megmentett minket irgalmasságból, s a Szentlélekben való újjászületés és megújulás f\underline{ü}rdőj\underline{é}ben.
}{%
    Valamikor testi vágyainkban \underline{é}ltünk, * és születésünknél fogva a harag gyermeke\underline{i} v\underline{o}ltunk.
}

\section*{Hétfő, 15. hét - 5. nap}

\parttitle[18, 16b-40]{Első olvasmány}
A királyok első könyvéből

\begin{lectio}{Illés Kármel hegyén legyőzi Baal papjait}
\noindent Azokban a napokban: Acháb elindult Illés felé. Amikor Acháb megpillantotta Illést, azt kérdezte: „Te vagy az, Izrael megrontója?” De ő így felelt: „Nem én sodortam pusztulásba Izraelt, hanem te és atyádnak háza, amikor elhagytátok az Urat, és amikor Baalhoz szegődtél. Most azonban küldj el s hívd nekem össze egész Izraelt Kármel hegyén, aztán Baal négyszázötven prófétáját is, akik Izebel asztaláról esznek!” Acháb tehát sorra üzent Izrael fiainak és összehívta a prófétákat is Kármel hegyén. Illés odalépett az egész nép színe elé és megkérdezte: „Meddig akartok még kétfelé sántikálni? Ha az Úr az Isten, akkor őt kövessétek; ha meg Baal, akkor kövessétek azt!” A nép nem válaszolt egyetlen szót sem.

Akkor Illés így szólt a néphez: „Az Úr prófétái közül egyedül én maradtam meg, Baalnak azonban négyszázötven prófétája van. Adjatok két bikát! Válasszák ki az egyiket, aztán vágják darabokra és tegyék a máglyára, de tüzet ne gyújtsanak! Én majd előkészítem a másik bikát, de tüzet én sem gyújtok. Akkor hívják segítségül ők a ti istenetek nevét, én meg majd segítségül hívom az Úr nevét! Az az Isten, aki tűzzel válaszol, az legyen az Isten!” Az egész nép azt felelte rá: „Rendben van!” Illés tehát így szólt Baal prófétáihoz: „Válasszátok ki magatoknak az egyik bikát! Ti vagytok többségben. Aztán hívjátok segítségül istenetek nevét, de tüzet nem szabad gyújtanotok!” Fogták hát az egyik bikát, előkészítették, és reggeltől délig szólítgatták Baal nevét ezekkel a szavakkal: „Baal, hallgass meg!” De semmi hang nem hallatszott, nem adott senki feleletet. Közben ott ugrabugráltak oltáruk körül, amit csináltak.

Délben Illés így csúfolta őket: „Szólítsátok hangosabban, hiszen isten! Hátha belemélyedt gondolataiba vagy kiment, esetleg úton van, vagy éppen elaludt, és föl kell ébreszteni.” Erre egyre emeltebb hangon szólították és szokásukat követve karddal és lándzsával addig vagdosták magukat, míg ki nem buggyant a vérük. Dél elmúltával odáig fajult a dolog, hogy őrjöngeni kezdtek egészen addig, míg el nem jött az esti ételáldozat bemutatásának ideje. De nem hallatszott se hang, se felelet, semmi jele a meghallgatásnak. Akkor Illés így szólt a néphez: „Lépjetek elém!” Az egész nép eléje lépett. Ezután helyreállította az Úr oltárát, amelyet leromboltak. Mégpedig úgy, hogy fogott tizenkét szikladarabot Jákob tizenkét fia törzsének megfelelően, akinek az Úr azt mondta: „Izrael legyen a neved!” és oltárt épített a szikladarabokból az Úr nevének.

Az oltárt körülvette árokkal, akkorával, amekkora két mérő gabonának elegendő. Aztán máglyát rakott, feldarabolta a bikát, rátette a máglyára és azt mondta: „Töltsetek meg négy korsót vízzel, és öntsétek rá az égőáldozatra meg a máglyára.” Megtették. Erre így szólt: „Ismételjétek meg!” Megismételték. Akkor azt mondta: „Harmadszor is tegyétek meg!” Harmadszor is megtették, úgyhogy víz folyt az oltár körül. Az árkokat is megtöltötte vízzel. Amikor aztán elérkezett az esti ételáldozat bemutatásának az ideje, Illés próféta előlépett és felkiáltott: „Uram, Ábrahám, Izsák és Izrael Istene! Nyilvánítsd ki a mai napon, hogy te vagy az Isten Izraelben, én a te szolgád vagyok s ezeket mind a te szavadra teszem! Hallgass meg, Uram, hallgass meg! Engedd, hogy ez a nép fölismerje: te, az Úr vagy az Isten, te téríted meg a szívét.”

Erre tűz hullott az Úrtól, megemésztette az égőáldozatot és a máglyát, még az árokban lévő vizet is elnyelte. Ennek láttára az egész nép arcra borult és így szólt: „Az Úr az Isten, az Úr az Isten!” Illés erre meghagyta nekik: „Ragadjátok meg Baal prófétáit! Ne meneküljön meg egy se közülük!” Megragadták, és Illés levitette őket a Kison patakhoz, s ott megölette őket.
\end{lectio}

\parttitle[1 Kir 18, 21; Mt 6, 24]{Válaszos ének}

\versicle{%
    Illés odalépett az egész nép színe elé és m\underline{e}gkérd\underline{e}zte: † Meddig akartok még kétfelé sánt\underline{i}k\underline{á}lni? * Ha az Úr az Isten, akkor őt k\underline{ö}vess\underline{é}tek!
}{%
    Senki sem szolgálhat két Úrn\underline{a}k; † nem szolgálhattok az Istenn\underline{e}k is, * a Mamm\underline{o}nn\underline{a}k is.
}

\parttitle{Második olvasmány}
Szent Ambrus püspök „A misztériumok” című értekezéséből

{\centering\redtext{(Nn. 8-11: SCh 25 bis, 158-160)}}

\begin{lectio}{Újjászülettünk vízből és Szentlélekből}
\noindent Mit láttál a keresztelő kápolnában? Láttál ott vizet, de nemcsak ezt, láttál ott levitákat, akik szolgáltak, és láttad a püspököt, aki kérdéseket tett fel, és végezte a beavatást. Mindenkinél előbb maga az Apostol tanít arra, hogy \textit{ne a láthatóra, hanem a láthatatlanra fordítsuk figyelmünket, mert a látható mulandó, a láthatatlan azonban örök} (2 Kor 4, 18). Másutt meg ezt olvashatod: \textit{Ami Istenben láthatatlan: örök ereje és isteni mivolta, arra a világ teremtése óta műveiből következtethetünk} (Róm 1, 20). Ezért mondja maga az Úr is: \textit{Ha nekem nem hisztek, higgyetek a tetteknek!} (Jn 10, 38) Hidd el tehát, hogy az Istenség van ott jelen. Hiszel abban, hogy ott működik, és nem hiszed azt, hogy jelen van? Honnan is származnék a működése, ha meg nem előzné a jelenléte?

Fontold meg, mily ősrégi ez a misztérium, hiszen már a világ kezdetén megvan az előképe. Kezdetben, amikor Isten az eget és a földet teremtette – úgymond a Szentírás – \textit{Isten lelke lebegett a vizek fölött} (Ter 1, 2). Vajon az, aki ott lebegett a vizek fölött, nem volt tevékeny? Vedd tudomásul, hogy működött már magában a világ teremtésében is, hiszen azt mondja a Próféta: \textit{Az egek az Úr szavára lettek, csillagseregüket az ő ajka lehelte} (Zsolt 32, 6). Mindkét tény a Szentírás bizonyságán nyugszik: az is, hogy lebegett, az is, hogy működött. Hogy ott lebegett, azt Mózes mondja; hogy ott működött, azt Dávid bizonyítja.

Halld a másik bizonyítékot is. Minden test megromlott a gonoszságtól. \textit{Nem marad éltető Lelkem az emberben örökké, mivel test} (Ter 6, 3) – mondja az Úr. Ezzel Isten arra mutat rá, hogy a test tisztátalansága és a súlyosabb bűn szennye a lelki kegyelmet elfordítja tőlünk. Azért, hogy az Úristen helyrehozza, ami megromlott, vízözönt bocsátott a földre, és az igaz Noénak megparancsolta, hogy szálljon a bárkába. Noé pedig a víz apadtával először hollót bocsátott ki, majd azután galambot, ez – mint olvassuk – olajággal tért vissza a bárkába. Látod, itt szó van vízről, fáról, galambról, és kétségbe vonod, hogy mindez titokzatos előkép?

A vízbe merül az ember teste, hogy megtisztuljon minden testi bűntől. Ez a víz minden gonosztettnek a temetője. A fa pedig annak a fának az előképe, amelyre az Úr Jézust szegezték fel, amikor értünk kínhalált szenvedett. A galamb azt a galambot jelzi, amelynek képében a Szentlélek szállt le, mint ahogy tanultad az Újszövetségben. A Szentlélek leheli szívedbe a békét és lelkedbe a megnyugvást.
\end{lectio}

\parttitle[Iz 44, 3. 4; Jn 4, 14]{Válaszos ének}

\versicle{%
    Elárasztom vízzel a t\underline{i}kkadt m\underline{e}zőt, † és bővizű patakokkal a kiasz\underline{o}tt f\underline{ö}ldet. * Kiárasztom Lelkemet, úgy sarjadzanak majd, mint a fűzfák a vízfoly\underline{á}sok m\underline{e}llett.
}{%
    Örök életre szök\underline{e}llő * vízf\underline{o}rr\underline{á}s lesz.
}

\section*{Kedd, 15. hét - 6. nap}

\parttitle[19, 1-9a. 11-21]{Első olvasmány}
A királyok első könyvéből

\begin{lectio}{Isten kinyilatkoztatja magát Illésnek a Hóreben}
\noindent Azokban a napokban: Acháb elbeszélte Izebelnek, amit Illés végbevitt, az egész történetét annak, ahogyan sorra megölette a prófétákat. Erre Izebel hírvivőt küldött Illéshez. Ezt üzente neki: „Ezt és ezt tegyék velem az istenek, ha holnap ebben az órában hasonlóvá nem teszem életedet ezek egyikének életéhez!”

Illés megrettent, útra kelt és elment, hogy megmentse életét. Amikor a Júdához tartozó Beersebába ért, otthagyta szolgáját, maga pedig behúzódott egy napi járásnyira a pusztába. Amikor odaért, leült egy borókabokor alá, és a halálát kívánta. Azt mondta: „Most már elég, Uram! Vedd magadhoz lelkemet! Én sem vagyok különb atyáimnál.” Ezzel lefeküdt és elaludt. Egyszer csak angyal érintette meg, és így szólt hozzá: „Kelj föl, és egyél!” Ahogy odapillantott, lám, a fejénél egy sült cipó meg egy korsó víz volt. Evett is, ivott is, de aztán újra lefeküdt aludni.

Ám az Úr angyala másodszor is megjelent, megérintette, és azt mondta: „Kelj föl, és egyél! Különben túl hosszú lesz neked az út.” Fölkelt, evett, ivott, aztán negyven nap és negyven éjjel vándorolt ennek az ételnek az erejéből egészen az Isten hegyéig, a Hórebig.

Bement egy barlangba, és ott töltötte az éjszakát. S lám, az Úr hallatta szavát, így szólt hozzá: „Menj, és a hegyen járulj az Úr színe elé!” S lám, az Úr elvonult arra. Hegyeket tépő, sziklákat sodró, hatalmas szélvész haladt az Úr előtt, de az Úr nem volt a szélviharban. A szélvésznek földrengés lépett a nyomába, de az Úr nem volt a földrengésben. A földrengés után tűz következett, de az Úr nem volt a tűzben. A tüzet enyhe szellő kísérte. Amikor Illés észrevette, befödte arcát köntösével, kiment, és a barlang szája elé állt. Egy hang megszólította ezekkel a szavakkal: „Mit csinálsz itt, Illés?” Azt felelte: „Emészt a buzgalom az Úrért, a Seregek Istenéért. Mert Izrael fiai elhagytak, oltáraidat lerombolták, prófétáidat meg kardélre hányták, és most nekem is az életemre törnek.” Az Úr azonban azt mondta neki: „Menj, fordulj vissza, és vedd utadat Damaszkusz pusztasága felé! Aztán menj, és kend föl Hazaelt Arám királyává, Jehut, Nimsi fiát pedig Izrael királyává. Végül Elizeust, Safát fiát Abel-Mecholából kend fel magad helyett prófétává. Aztán ez történik: Aki Hazael kardjától megmenekül, azt megöli Jehu. És aki Jehu kardjától megmenekül, azt megöli Elizeus. De hétezret életben hagyok Izraelben: minden térdet, amely nem hajolt meg Baal előtt, és minden szájat, amely nem illette csókkal (a kezét).”

Amikor onnét elment, találkozott Safát fiával, Elizeussal. Éppen szántott, tizenkét pár ökör haladt előtte, maga pedig a tizenkettedik pár mögött haladt. Illés odament hozzá, és rávetette köntösét. Az otthagyta az ökröket, Illés után szaladt, és így szólt: „Engedd, hadd adjak előbb búcsúcsókot apámnak meg anyámnak, aztán követlek.” Azt mondta neki: „Menj csak, de térj vissza! De hát mit tettem veled?” Azzal Elizeus elment, fogta a pár ökröt, feláldozta, majd az ekét felhasználva megsütötte az ökröket, és odaadta az embereknek, egyék meg őket. Aztán elindult, Illés nyomába szegődött, és a szolgája lett.
\end{lectio}

\parttitle[Vö. Kiv 33, 22. 20; Jn 1, 18]{Válaszos ének}

\versicle{%
    Így szólt az \underline{Ú}r Móz\underline{e}shez: † Ha majd elvonul előtted dicsőségem, a szikla mélyedésébe teszlek és kezemmel befödlek, amíg elvonulok \underline{e}l\underline{ő}tted. * Mert nem láthatja ember Istent úgy, hogy életb\underline{e}n mar\underline{a}djon.
}{%
    Istent nem látta soha s\underline{e}nki, † az egyszülött Fiú nyilatkoztatt\underline{a} ki, * aki az Atya \underline{ö}l\underline{é}n van.
}

\parttitle{Második olvasmány}
Szent Ambrus püspök „A misztériumok” című értekezéséből

{\centering\redtext{(Nn. 12-16. 19: SCh 25 bis, 162-164)}}

\begin{lectio}{Mindez előkép a számunkra}
\noindent Arra tanít téged az Apostol: \textit{Atyáink mind ott voltak a felhő alatt, mind átkeltek a tengeren, s mind megkeresztelkedtek Mózesre, a felhőben és a tengerben} (1 Kor 10, 1-2). De maga Mózes is azt mondja énekében: \textit{Elküldted Lelkedet és a tenger elborította őket} (Kiv 15, 10). Észrevetted ugye, hogy a keresztség szentségének előképe rejtőzik a zsidóknak abban az átvonulásában, amelyben az egyiptomi elveszett, a zsidó pedig megmenekült. Hisz mi másra tanít minket naponként ez a szentség, mint hogy benne a bűn elmerül és a vétek elpusztul, de jámborság és teljes ártatlanság lép ki belőle.

Hallottad azt is, hogy atyáink a felhő alatt voltak, mégpedig jó felhő alatt, amely a testi szenvedélyek lobogását lehűtötte. Jó felhő árnyékozza be azokat, akiket a Szentlélek meglátogat. Végül Szűz Máriára szállt rá, és a Magasságbeli ereje borította be őt árnyékával, amikor az emberi nem számára megszülte a Megváltót. Az a felhőcsoda is előkép szerint történt Mózes által. Ha tehát akkor az előképben jelen volt a Szentlélek, vajon nincs-e jelen most a jelzett igazságban, hiszen a Szentírás is mondja neked: \textit{A törvényt Mózes közvetítette, a kegyelem és igazság azonban Jézus Krisztus által valósult meg} (Jn 1, 17).

A „Mara” forrás vize keserű volt, Mózes fát dobott bele, és édes lett. Amíg az Úr keresztjét nem hirdették, a víznek a jövendő üdvösség szempontjából nem volt semmi jelentősége, amikor azonban a kereszt misztériuma már üdvösségszerzővé szentelte, azóta lelki fürdőnek és életadó italnak használható. Miként tehát abba a vízbe Mózes, a próféta fát bocsájtott, ugyanúgy ebbe a forrásba is a pap az Úr keresztjének hirdetését bocsátja, és ez a víz édessé válik a kegyelem számára.

Ne higgy tehát tisztán testi szemeidnek: itt inkább azt kell nézni, ami láthatatlan, mert a látható ideigvaló, amaz pedig örökkévaló. Jobban lehet itt azt látni, amit a testi szem nem vesz észre, de a szív és a lélek meglát.

Azután okulj a Királyok könyvének olvasmányán is (2 Kir 5, 1-27). Naamán szír ember volt, leprás lett, és senki nem tudta abból kigyógyítani. Akkor az egyik rableány azt mondta neki, hogy van Izraelben egy próféta, aki a leprától meg tudná őt tisztítani. Naamán aranyat és ezüstöt vett magához, elment Izrael királyához. Amikor ez meghallotta érkezésének okát, megszaggatta ruháját, mondván, hogy istenkísértés az, amikor olyat kérnek tőle, amit nem tud megadni királyi hatalma. Elizeus azonban titokban megüzente a királynak, küldje el hozzá ezt a szír embert, hadd tudja meg, hogy van Isten Izraelben. Amikor pedig az odaérkezett, meghagyta neki, hogy merüljön alá hétszer a Jordán vizében. Erre Naaman azon kezdett gondolkodni, hogy jobb az ő hazája folyóinak a vize, és azokban már sokszor megfürdött, de mégsem szabadult meg leprájától. Ez a meggondolása visszatartotta attól, hogy szót fogadjon a prófétának. De amikor engedett szolgái biztatásának és rábeszélésének, mégis alámerült a Jordánban. Amikor pedig abban megtisztult, menten megértette, hogy nem a víz, hanem a kegyelem ereje tisztítja meg az embert.

Az a pogány kételkedett, mielőtt meggyógyult volna; te már meggyógyultál, ne kételkedj tehát.
\end{lectio}

\parttitle[Zsolt 77, 52. 53; 1 Kor 10, 2]{Válaszos ének}

\versicle{%
    Kihozta népét az Úr, mint \underline{a} juh\underline{o}kat, † biztosan vezette őket, hogy ne f\underline{é}lj\underline{e}nek, * S ellenségeiket elnyelt\underline{e} a t\underline{e}nger.
}{%
    M\underline{i}nd megkeresztelkedtek Móz\underline{e}sre * a felhőben és a t\underline{e}ng\underline{e}rben.
}

\section*{Szent Bonaventura püspök és egyháztanító - 7. nap}

\parttitle[21, 1-21. 27-29]{Első olvasmány}
A királyok első könyvéből

\begin{lectio}{Illés a szegények igazának a védelmezője}
\noindent Abban az időben: Nabotnak, a jiszreelitának volt egy szőlőskertje Achábnak, Szamaria királyának a palotája mellett. Acháb beszélt Nabottal. Így szólt hozzá: „Engedd át nekem a szőlőskertedet, hogy veteményeskertül szolgáljon nekem; egész közel fekszik házamhoz. Adok érte egy jobb szőlőskertet, vagy ha úgy tetszik jobban, kifizetem az árát pénzben.” Nabot azonban azt válaszolta Achábnak: „Az Úr óvjon attól, hogy atyáim örökségét átengedjem neked!” Acháb hazament. Kedvetlen volt és haragos a válasz miatt, amit Nabot, a jiszreelita adott neki, amikor azt mondta: „Nem engedem át neked atyáim örökségét.” Leheveredett, befordult, és nem evett semmit.

Odament hozzá a felesége, Izebel, és megkérdezte: „Mi van veled, hogy olyan kedvetlen vagy, és nem eszel semmit?” „Megbeszélésem volt a jiszreeli Nabottal – felelte. Azt mondtam neki: Engedd át nekem szőlőskertedet arany ellenében, vagy ha úgy tetszik jobban, adok neked érte egy másik szőlőskertet. De kijelentette: »Nem engedem át neked atyáim örökségét!«” Felesége, Izebel erre azt mondta neki: „Hát nem vagy Izrael királya? Kelj föl, egyél, és légy jókedvű. Majd én megkerítem a jiszreeli Nabot szőlejét.”

Aztán leveleket írt, lepecsételte az ő pecsétjével, és elküldte őket a Nabottal együtt lakó véneknek és előkelőknek. Ezt írta a levelekben: „Hirdessetek böjtöt, és Nabotot ültessétek az asztalfőre az emberek közé. Aztán vele szemben két semmirekellő ember kapjon helyet, akik tanúskodnak majd ellene, és kijelentik: »Káromoltad Istent és a királyt.« Aztán vezessétek ki, és kövezzétek agyon!” S városának lakói, a vének és előkelők, akik együtt laktak vele a városban, úgy tettek, amint Izebel parancsolta, ahogy a nekik küldött levelekben írva volt. Böjtöt hirdettek, és Nabotot az asztalfőre ültették az emberek közé. Aztán előkerült két ember, semmirekellő ember, s leült vele szemben. Tanúságot tettek, és kijelentették: „Nabot káromolta Istent és a királyt.” Erre kivezették a város elé, és agyonkövezték. Majd jelentették Izebelnek: „Nabotot megköveztük, halott.”

Mihelyt Izebel meghallotta, hogy Nabotot agyonkövezték, azt mondta Izebel Achábnak: „Nos, vedd birtokba Nabot szőlejét, amelyet nem akart neked pénzért átengedni. Mert Nabot nem él többé, halott.” Amikor Acháb meghallotta, hogy Nabot halott, fogta magát Acháb, elment a jiszreelita Nabot szőlejébe, és birtokba vette.

Az Úr hallatta szavát, és azt mondta a tisbei Illésnek: „Indulj, menj el Achábhoz, Izrael királyához, aki Szamariában székel; épp Nabot szőlejében van, lement, hogy birtokba vegye. És mondd meg neki: Így szól az Úr: Gyilkoltál, és most az örökséget is eltulajdonítod? Ezért azt mondja az Úr: Ott fogják a kutyák a te véredet is felnyalni, ahol Nabot vérét felnyalták a kutyák!” Acháb azt mondta: „Hát rám találtál, ellenségem?” „Rád találtam – felelte. Mert arra adtad magad, hogy azt tedd, ami gonosznak számít az Úr szemében. Lásd, romlást hozok rád, elsöpörlek, kiirtom Acháb házából mind aki férfi, a szolgát is, a szabadot is Izraelben.”

Amikor Acháb meghallotta ezeket a szavakat, megszaggatta ruháját, vezeklő ruhát öltött mezítelen testére, és böjtölt. Aludni is vezeklő ruhában tért és kedvetlenül járt-kelt. Az Úr ezért így szólt a tisbei Illéshez: „Látod, hogy megalázkodott előttem Acháb? Mivel megalázkodott előttem, nem hagyom, hogy életében rátörjön a baj házára, majd csak fia idejében sújtom házát romlással.”
\end{lectio}

\parttitle[Jak 4, 8. 9. 10; 5, 6]{Válaszos ének}

\versicle{
    Mossátok tisztára kezetek\underline{e}t, bűn\underline{ö}sök, † és tisztítsátok meg szíveteket, megosztott l\underline{e}lk\underline{ű}ek. * Gyászoljatok és szomorkodjatok, alázkodjatok meg \underline{a}z Úr \underline{e}lőtt.
}{%
    \underline{A}z igazat elítéltétek, † meg\underline{ö}ltétek, * s ő nem tanúsított ell\underline{e}n\underline{á}llást.
}

\parttitle{Második olvasmány}
Szent Bonaventura püspöknek „A lélek útikönyve Istenhez” című könyvecskéjéből

{\centering\redtext{(Cap. 7, 1. 2. 4. 6; Opera omnia, 5, 312-313)}}

\begin{lectio}{A Szentlélek által kinyilatkoztatott misztikus bölcsesség}
\noindent Krisztus az út és az ajtó. Krisztus a lépcső és a segítség az Istenhez jutásban, mivel ő \textit{a szövetség ládájára helyezett kiengesztelési tábla} (Kiv 26, 34), és ő \textit{a kezdettől elrejtett titok} (Ef 3, 6). Aki erre az engesztelő táblára teljesen odafordított arccal néz, nézi a keresztre feszített Krisztust hittel, reménnyel, szeretettel, áhítattal, csodálattal, örvendezéssel, megbecsüléssel, dicsérettel és ujjongással, az ilyen ember pászkát, azaz átvonulást vállal vele, hogy a kereszt vesszejével átmenjen a Vörös-tengeren, Egyiptomból kivonul a pusztába, ahol titokzatos mannát ízlel, és Krisztussal együtt nyugszik a sírban, külsőleg mintegy meghalva, de lélekben megérezve mégis azt – már amennyire e földi életben lehetséges –, amit a hozzá forduló gonosztevőnek mondott Krisztus a kereszten: \textit{Még ma velem leszel a paradicsomban} (Lk 23, 43).

Ebben az átvonulásban, ha ez tökéletes, minden értelmi tevékenységet el kell hagyni, és a szeretet minden hevének Isten felé és Istenbe kell átfordulnia. Ez azonban misztikus és titokzatos dolog, amit senki más nem ismer, csak aki megkapja; senki más nem kapja meg, csak aki vágyódik utána; de vágyódni sem tud utána senki más, csak akinek szívét a Szentlélek felgyújtotta, az a Szentlélek, akit Krisztus küldött a földre. Ezért is mondja az Apostol, hogy ezt a misztikus bölcsességet a Szentlélek nyilatkoztatta ki.

Ha azt kutatod, hogyan történik mindez, a kegyelmet kérdezd, ne a tanítást; a vágyat kérdezd, ne az értelmet; az imádság sóhajtásait, ne a tanulmányokat; a jegyest, ne a tanítót; Istent, ne az embert; a felhőt, ne a ragyogást; ne a fényt, hanem a mindenestől lángra gyújtó tüzet, amely elragadó illatával és lángoló szeretetével visz Istenhez. Ez a tűz pedig az Isten, ez a tüzes kemence Jeruzsálemben van, és Krisztus gyújtotta meg izzó szenvedésének lobogásában, amit egyedül csak az fog fel igazán, aki ezt mondja: A \textit{megfulladást áhította a lelkem, s csontjaim a halált} (vö. Jób 7, 15). Aki a halált szereti, az megláthatja Istent, mert kétségtelenül igaz: \textit{Nem láthat meg engem élő ember} (Kiv 33, 20). Haljunk meg tehát, és lépjünk be a felhőbe, parancsoljunk csendet gondjainknak, vágyainknak és képzeletünknek; vonuljunk át a megfeszített Krisztussal \textit{ebből a világból az Atyához}, hogy miután megmutatta nekünk az Atyát, mondhassuk Fülöppel: \textit{Ez elég nekünk} (Jn 14, 8), és hallhassuk Pállal: \textit{Elég neked az én kegyelmem} (2 Kor 12, 9), és Dáviddal is ujjongva mondhassuk: \textit{Érje bár testemet és szívemet enyészet, szívem Istene és örökrésze mindig az Isten} (vö. Zsolt 72, 26). \textit{Áldott legyen az Úr öröktől fogva mindörökké. Az egész nép mondja: Úgy legyen! Úgy legyen!} (Zsolt 105, 48)
\end{lectio}

\parttitle[1 Jn 3, 24; Sir 1, 9. 10]{Válaszos ének}

\versicle{%
    Aki teljesíti Isten parancsait, benne marad \underline{a}z Ist\underline{e}nben, † és az Isten is \underline{ő}b\underline{e}nne. * Annak, hogy ő bennünk marad-e, a Lélek a próbaköve, akit \underline{a}dott n\underline{e}künk.
}{%
    \underline{A}z Úr teremtette a bölcsességet a Szentlél\underline{e}kben, * és kiárasztotta minden műv\underline{ei}re. 
}

\section*{Csütörtök, 15. hét - 8. nap}

\parttitle[22, 1-9. 15-23. 29. 34-38]{Első olvasmány}
A királyok első könyvéből

\begin{lectio}{Prófétai szó előremondja a gonosz Acháb végzetét}
\noindent Azokban a napokban: Három évig nyugalom volt, nem dúlt háború Arám és Izrael között. A harmadik esztendőben azonban Júda királya, Jehosafát elment látogatóba Izrael királyához, (Achábhoz). Akkor Izrael királya így szólt szolgáihoz: „Tudjátok jól, hogy Gileádban Ramot hozzánk tartozik, s mi ülünk tétlenül, ahelyett, hogy kiragadnánk Arám kezéből.” Majd Jehosafáthoz fordult: „Hadba vonulsz velem a gileádi Ramot ellen?” Jehosafát azt felelte Izrael királyának: „Ott leszek éppen úgy, mint te, népem, mint a te néped, lovaim, mint a te lovaid.” Aztán ezt mondta Jehosafát Izrael királyának: „Kérj előbb útmutatást az Úrtól!” Izrael királya összehívta hát a prófétákat, mintegy négyszáz embert, és így szólt hozzájuk: „Hadba vonuljak a gileádi Ramot ellen, vagy mondjak le róla?” Ezt a választ adták: „Vonulj fel, az Úr a király kezébe adja!” Jehosafát azonban így szólt: „Nincs itt véletlenül az Úrnak valamelyik prófétája, hogy megkérdezhetnénk?” „Van itt egy – válaszolta Izrael királya Jehosafátnak –, aki által megkérdezhetnénk az Urat, bár a magam részéről nem kedvelem, mert sose jövendöl nekem jót, mindig csak rosszat: Jimla fia, Michajehu.” Jehosafát azonban megjegyezte: „Ne beszéljen így a király!” Izrael királya erre hívatta egyik udvari emberét, és azt mondta neki: „Hozd ide gyorsan Jimla fiát, Michajehut!”

Amikor a király elé ért, a király így szólt hozzá: „Michajehu, hadba vonuljunk a gileádi Ramot ellen, vagy mondjunk le róla?” Azt válaszolta: „Vonulj! Győzelmet aratsz, az Úr a király kezébe adja!” A király azonban megkérdezte: „Hányszor eskesselek meg, hogy csak az igazat mondod az Úr nevében?” Erre megszólalt: „Látom egész Izraelt, szétszórva a hegyen, mint a pásztor nélkül maradt nyájat. S az Úr ezt mondja: Nincs gazdájuk, térjen hát vissza ki-ki békében házába!” Izrael királya Jehosafáthoz fordult: „Hát nem megmondtam, hogy sose jövendöl nekem jót, hanem mindig csak rosszat?” De az folytatta: „Halld hát az Úr szavát: Láttam az Urat, a trónján ült, az egész mennyei sereg ott állt a jobbján és a balján. S azt mondta az Úr: »Ki veszi rá Achábot, (Izrael királyát), hogy hadba vonuljon és Ramotnál, Gileádban elessen?« Az egyik ezt mondta, a másik azt mondta, míg végül az Úr elé állt a lélek, és így szólt: »Én veszem rá.« S az Úr megkérdezte: »Milyen úton-módon? Hogyan?« »Én megyek – felelte –, és a hazugság lelke leszek a prófétái ajkán!« Erre azt mondta neki: »Vedd hát rá, hiszen képes vagy rá! Csak menj, és tedd meg!« S lám az Úr valóban a hazugság lelkét adta ezeknek a te prófétáidnak az ajkukra, mert az Úr vesztedet határozta el.”

Izrael királya és Júda királya, Jehosafát így hadba vonult a gileádi Ramot ellen. De egy ember vaktában kifeszítette íját, s eltalálta Izrael királyát a mellvért és az öv között. Erre megparancsolta szekere hajtójának: „Fordulj meg, és vigyél ki az ütközetből, mert megsebesültem!” A harc azonban egyre hevesebbé vált azon a napon, a király kénytelen volt a szekerén állva maradni, szemben az arámokkal. Estére meghalt. Sebéből a vér a kocsi aljára folyt. Naplementekor kiáltás hallatszott a táborban: „Vissza! Ki-ki a maga városába, ki-ki a maga földjére! A király halott!”

Szamariába mentek és Szamariában temették el a királyt. Amikor Szamaria tavában lemosták a kocsit, a kutyák nyalták fel a vérét és a ringyók fürödtek benne az Úr szava szerint.
\end{lectio}

\parttitle[Jer 29, 8. 9. 11; MTörv 18, 18]{Válaszos ének}

\versicle{%
    Nehogy félrevezessenek titeket a köztetek lak\underline{ó} próféták, † mert hazugságot jövendölnek nektek az én n\underline{e}v\underline{e}mben; * Tudom, milyen terveket gondoltam el felőletek, m\underline{o}ndja \underline{a}z Úr.
}{%
    Prófétát tám\underline{a}sztok, * és ajkára adom szav\underline{ai}mat.
}

\parttitle{Második olvasmány}
Szent Ambrus püspök „A misztériumok” című értekezéséből

{\centering\redtext{(Nn. 29-30. 34-35. 37. 42: SCh 25 bis, 172-178)}}

\begin{lectio}{A keresztség utáni szertartások magyarázata.}
\noindent Miután a keresztvízben újjászülettél, odaléptél a pap elé. Fontold meg, mi következett ezután. Nemde az, amiről Dávid így beszél: \textit{Mint fejen a drága kenet, amely végigcsordul a szakállon, Áron szakállán} (Zsolt 132, 2). Ez az az olaj, amelyről Salamon is azt mondja: \textit{Mint a kiöntött olaj, olyan a neved, azért szeretnek a lányok} (Én 1, 3). Hány újjászületett lélek kedvelt meg ma téged, Urunk, Jézus, és így szól hozzád: \textit{Vigyél magaddal, ruhád jó illata után sietünk} (Én 1, 4), hogy a feltámadás illatából meríthessenek.

Értsd meg, hogy miért van ez, \textit{hiszen a bölcs maga elé néz} (Préd 2, 14). Azért csordult végig arcodon az olaj, hogy jelezze újjászületésed örömét; azért mondja: \textit{Áron szakállán}, hogy te is \textit{a kiválasztott nemzetség} (1 Pét 2, 9), a papi és szent nép tagja légy. Hiszen ez a lelki kegyelem mindannyiunkat az Isten országa számára és papságára ken fel.

Ezután fehér ruhát kaptál annak jeléül, hogy levetetted a bűn öltözetét, és magadra öltötted az ártatlanság tiszta ruházatát; erről mondja a Próféta: \textit{Hints meg izsóppal, és tiszta leszek, moss meg, és fehérebb leszek, mint a hó} (Zsolt 50, 9). Mert akit alámerítenek akár a törvény, akár az evangélium szerint, az megtisztultan lép ki onnan: a törvény szerint, mivel Mózes izsóppal hintette szét a népre a bárány vérét, az evangélium szerint pedig, mivel Krisztusnak fehér volt a ruhája, mint a hó, hiszen feltámadásának dicsőségét mutatta be. Hónál fehérebb lesz az, aki bűnbocsánatot nyer. Azért mondja Izajás által az Úr: \textit{Ha olyanok volnának is bűneitek, mint a skarlát, fehérek lesznek, mint a hó} (Iz 1, 18).

Az Egyház ezt a ruhát \textit{a megújulás fürdőjében} (Tit 3, 5) vette magára. Az Énekek énekében azt mondja: \textit{Sötét bőrű vagyok, de azért szép, Jeruzsálem leányai} (Én 1, 4). Feketének mondja magát az emberi természet gyarlósága miatt, ékesnek a kegyelem miatt. Fekete, mivel a bűnösök közül való, de ékes a hit szentsége által. Amikor azonban Jeruzsálem leányai meglátták őt fehér ruhában, álmélkodva kérdezgették: \textit{Kicsoda ő, aki fehér ruhában lép elő?} (Vö. Jel 7, 14). Előbb fekete volt, mitől lett most hirtelen fehérré?

Zakariás prófétánál olvasod, hogy Jézus az Egyházáért \textit{a bűnök szennyes ruházatát} (vö. Zak 3, 3) vette magára. Most viszont látva az Egyházat hófehér ruhában, és ugyanígy látva a megújulás fürdőjében a most megtisztult és fehérre mosott lelket is, Krisztus így szólítja meg: \textit{De szép vagy, kedvesem, igen szép vagy, szemeid, mint a galambok} (Én 4, 1); a Szentlélek is galamb képében szállt le az égből.

Emlékezz csak vissza, hogy lelkedbe kaptad pecsétként \textit{a bölcsesség és értelem, a tanács és erősség, a tudás és jámborság és az Úr félelmének lelkét} (Iz 11, 2). Őrizd meg azt, amit megkaptál. Megjelölt téged az Atyaisten, megerősített az Úr Jézus Krisztus, és \textit{szívedbe árasztotta foglalóul} (2 Kor 1, 22) a Szentlelket, mint azt az Apostol leveléből is hallottad.
\end{lectio}

\parttitle[Ef 1, 13-14; 2 Kor 1, 21-22]{Válaszos ének}

\versicle{%
    Miután hittetek, megkaptátok a megígért Szentlél\underline{e}k pecs\underline{é}tjét, † ő a foglaló öröks\underline{é}g\underline{ü}nkre. * Hogy teljesen megváltva az öv\underline{é}i l\underline{e}gyünk.
}{%
    Isten, aki minket fölkent Kriszt\underline{u}sban, † pecsétjével megjelölt m\underline{i}nket, * és foglalóul a szívünkbe árasztotta \underline{a} L\underline{e}lket.
}

\section*{Péntek, 15. hét - 9. nap}

\parttitle[20, 1-9. 13-24]{Első olvasmány}
A királyok második könyvéből

\begin{lectio}{Isten csodás győzelmet ad a jámbor Jehosafát királynak}
\noindent Azokban a napokban: Moáb és Ammon fiai egy csapatnyi meunitával együtt hadba vonultak Jehosafát ellen. Mentek és jelentették neki: „Hatalmas tömeg jön ellened a tengeren túlról, Edomból. Immár Hacacon-Támárban, vagyis Engediben vannak.” Erre Jehosafát megijedt, és elhatározta, hogy az Úrhoz fordul. Egész Júdának böjtöt hirdetett. Erre Júda össze is jött, hogy az Úrhoz könyörögjön, sőt Júda valamennyi városából egybegyűltek, hogy könyörögjenek az Úrhoz. Ekkor Jehosafát az Úr házában az új udvar előtt odaállt Júda és Jeruzsálem gyülekezete elé. Így imádkozott: „Urunk, atyáink Istene! Te vagy az Isten a mennyben, és te vagy a népek összes királyságának uralkodója. A te kezedben olyan erő és hatalom van, hogy annak senki sem tud ellenállni. Istenünk, te űzted ki e föld lakóit néped, Izrael elől, és te adtad barátod, Ábrahám ivadékának mindörökre, ők le is telepedtek rajta, és szentélyt építettek nevednek azzal a szándékkal: Ha valami nyomorúság tör ránk, kard vagy büntetés, pestis vagy éhínség, akkor megállunk e templom előtt és a te színed előtt, mert ebben a házban van a te neved. Aztán hozzád kiáltunk szorongattatásainkban, és akkor te meghallgatsz és megsegítesz bennünket.”

Ezalatt egész Júda ott állt az Úr előtt gyermekeikkel és feleségeikkel együtt. Ekkor a sokaság közepette az Úr lelke leszállt egy Aszaf fiai közül való levitára, Jahazielre, Zekarjahu fiára, Bénája, Jeuel és Mattanja utódára, így kezdett beszélni: „Ide figyeljetek, egész Júda, Jeruzsálem lakói és te is, Jehosafát király! Ezt üzeni nektek az Úr: Ne féljetek, és ne is rettegjetek ettől a nagy sokaságtól, mert nem nektek kell ellenük harcolnotok, hanem az Istennek. Holnap gyertek le eléjük! A Cic nevű domboldalon fognak felvonulni, így rájuk találtok a völgy végén Jeruel pusztájával szemben. Nektek azonban nem kell harcolnotok. Csak álljatok oda, és maradjatok veszteg. Akkor meg fogjátok látni, hogyan segít meg az Úr benneteket, Júdát és Jeruzsálemet. Ne féljetek és ne rettegjetek! Holnap vonuljatok ki eléjük, és az Úr veletek lesz!”

Ekkor Jehosafát a földig hajolt, egész Júda és Jeruzsálem lakói leborultak az Úr előtt, hogy hódoljanak előtte. Aztán a Kehát fiai és a Korach fiai közül való leviták felálltak, hogy emelt hangon, harsogva magasztalják az Urat, Izrael Istenét. Kora reggel föl is keltek, és kivonultak Tekoa pusztájába. Elindulás előtt Jehosafát eléjük állt, és így szólt: „Hallgassatok rám, Júda és Jeruzsálem lakói! Bízzatok az Úrban, a ti Istenetekben, és akkor életben maradtok. Higgyetek prófétáinak, és boldogultok.” Miután ezt a tanácsot adta a népnek, felállította az Úr énekeseit, hogy kivonuláskor a hadba keltek élén szent díszben magasztalják őt e szavakkal: „Dicsérjétek az Urat, mert irgalma örökkévaló!”

Mihelyt felhangzott a dicsőítő és magasztaló ének, az Úr a lesben álló csapatokat azok ellen fordította, akik Júda ellen vonultak: Ammonnak, Moábnak és Szeirnek fiai ellen, így ezek vereséget szenvedtek. Aztán Ammon és Moáb fiai fölkeltek Szeir hegyének fiai ellen, hogy kiirtsák és elpusztítsák őket. Miután végeztek Szeir lakóival, az egyik férfi a másikat segítette a pusztításban. Amikor Júda a pusztai kilátóhelyhez ért, és a sokaság felé tekintett: lám, azok holtan feküdtek a földön. Senki sem menekült meg.
\end{lectio}

\parttitle[Ef 6, 12. 14; 2 Krón 20, 17]{Válaszos ének}

\versicle{%
    Nem annyira a vér és test ellen k\underline{e}ll küzd\underline{e}nünk, † hanem a fejedelemségek és hatalmasságok és a gonosz szellem\underline{e}k \underline{e}llen. * Így készüljetek föl: csatoljátok derekatokra az ig\underline{a}zság \underline{ö}vét.
}{%
    Álljatok \underline{o}da, † akkor meg fogjátok l\underline{á}tni, * hogyan segít meg az Úr benn\underline{e}teket.
}

\parttitle{Második olvasmány}
Szent Ambrus püspök „A misztériumok” című értekezéséből

{\centering\redtext{(Nn. 43. 47-49: SCh 25 bis, 178-180. 182)}}

\begin{lectio}{Az Eukarisztiáról az újonnan megkereszteltekhez}
\noindent A keresztségben megtisztult nép ezután ragyogó ruházatban Krisztus oltárához vonul, ezt imádkozva: \textit{Odalépek Isten oltárához, Istenhez, aki nékem ujjongó örömöm} (Zsolt 42, 4). Miután ugyanis levetette magáról az ősi tévedés bűnét, \textit{sasként megújult fiatal szívvel} (vö. Zsolt 102, 6) sietve igyekszik a mennyei lakomára. Megérkezik a körmenet, mindenki meglátja a feldíszített szent oltárt, és felkiált szívében: \textit{Asztalt terítesz nekem} (Zsolt 22, 5). Dávid szavait vesszük ajkunkra, amikor így kiáltunk: \textit{Az Úr az én pásztorom, nincs is semmiben hiányom, zöldellő mezőkön terelget engem, és csöndes vizekhez vezet} (Zsolt 22, 1). Aztán így folytatjuk: \textit{Járjak bár halálsötét völgyekben, ott sem félek a bajtól, hiszen te vagy énvelem. Pásztorbotod és terelővessződ az én bizodalmam. Asztalt terítesz nekem ellenségeim szeme előtt. Dús kenettel frissíted fejemet, színültig teli serlegem kicsordul} (Zsolt 22, 1-5).

Valóban csodálatos volt az, hogy az ősatyáknak Isten mannát hullatott, és hogy mindennap megadta nekik ezt a mennyei táplálékot. Azért mondja a Szentírás: \textit{Angyalok kenyerét ette az ember} (Zsolt 77, 25). Mégis, akik azt a kenyeret ették, mind meghaltak a pusztában, ez az étel azonban, amelyet te eszel, \textit{az élő kenyér, amely a mennyből szállt alá} (Jn 6, 51); ez az örök élet táplálékát nyújtja, s mindaz, aki ezt eszi, az \textit{nem hal meg mindörökké} (Jn 11, 26), mert ez Krisztus teste.

Most hasonlítsd össze te magad, vajon az angyalok kenyere, a manna kiválóbb-e vagy Krisztus teste, amely életet adó test. Az a manna az égből jött, ez annál is magasabbról; az mennyei volt, ez a menny Uráé; arra romlás várt, ha másnapra eltették, ez mentes minden romlástól, és bárki áhítattal ízleli ezt, annak nem árt semmiféle romlás. Az ősatyák számára a sziklából víz fakadt, számodra Krisztus vére folyik; az ő szomjukat csak egy ideig csillapította az a víz, téged ez a vér mindörökre megtisztított. A választott nép fia ivott és tovább szomjazott, te pedig, ha ebből iszol, nem maradhatsz szomjas. Az képletesen történt, ez pedig már valóságban.

Ha csodálod azt, ami csak jelkép, milyen nagy lehet az, aminek jelképét is csodálod. Halljad, hogy valóban előképben történt mindaz, ami az ősatyákkal történt, írva van: \textit{Ittak a lelki sziklából, amely kísérte őket, s a szikla Krisztus volt. De legtöbbjükben nem telt kedve az Istennek, ezért odavesztek a pusztában. Ez mind intő példa lett számunkra} (1 Kor 10, 4). Megismerted most, melyik kiválóbb. Mert többet ér a fény, mint az árnyék, többet a valóság, mint az előkép, többet ér a megváltás Szerzőjének teste, mint az égi manna.
\end{lectio}

\parttitle[1 Kor 10, 1-2. 11. 3-4]{Válaszos ének}

\versicle{%
    Atyáink mind átkeltek \underline{a} teng\underline{e}ren † s mind megkeresztelkedtek Mózesre a felhőben és a t\underline{e}ng\underline{e}rben. * Mindez előkép \underline{a} szám\underline{u}nkra.
}{%
    Mindnyájan ugyanazt a lelki ételt \underline{e}tték, * ugyanazt a lelki it\underline{a}lt \underline{i}tták.
}

\section*{Szent Hedvig királynő - 10. nap}

\parttitle[2, 1-15]{Első olvasmány}
A királyok második könyvéből

\begin{lectio}{Illés menybevitele}
\noindent Azokban a napokban, amikor az Úr azt akarta, hogy Illést a forgószél fölvigye az égbe, az történt, hogy Illés és Elizeus éppen elmentek Gilgalból. Illés így szólt Elizeushoz: „Maradj itt, mert engem az Úr Bételbe küld.” Elizeus azonban ezt felelte: „Úgy igaz, ahogy az Úr él, s úgy igaz, amint te élsz: nem hagylak el!”

Így hát elmentek Bételbe. A Bételben lakó prófétatanítványok odamentek Elizeushoz, és megkérdezték: „Tudod, hogy ma az Úr elragadja uradat a fejed fölött?” „Tudom én is, csak hallgassatok!” – felelte. Illés meg így szólt hozzá: „Elizeus, maradj itt, mert engem Jerikóba küld az Úr.” De az ezt mondta neki: „Úgy igaz, ahogy az Úr él, s úgy igaz, amint te élsz: nem hagylak el!”

Így hát elmentek Jerikóba. A Jerikóban élő prófétatanítványok Elizeushoz fordultak, és azt mondták neki: „Tudod, hogy ma az Úr elragadja a fejed fölött uradat?” „Tudom én is, csak hallgassatok!” – válaszolta. Illés pedig így szólt hozzá: „Maradj itt, mert engem az Úr a Jordán mellé küld.” De ő azt mondta: „Úgy igaz, amint az Úr él, és úgy igaz, ahogy te élsz: Nem hagylak el.” Így mind a ketten elmentek. A prófétatanítványok közül is elment ötven, s egy bizonyos távolságra maradtak tőlük, míg ők ketten a Jordán mellett álltak meg.

Illés fogta a köntösét, összegöngyölte, és ráütött vele a vízre. Erre az kettévált, az egyik és a másik irányban, úgyhogy mindketten szárazon átmentek rajta. Amikor átértek, Illés azt mondta Elizeusnak: „Kérj valamit, mit tegyek még meg neked, mielőtt elválok tőled!” Elizeus így felelt: „Hát akkor szálljon kétannyi rész rám lelkedből!” Azt mondta neki: „Olyat kértél, amit nehéz teljesíteni. Ha majd látod, mint ragadtatom el, osztályrészül jut neked, de ha nem, akkor nem kapod meg.”

S történt, amint mentek és beszélgettek, egyszer csak jött egy tüzes szekér tüzes lovakkal, s elválasztotta őket egymástól, aztán Illés a forgószéllel fölment az égbe. Amikor Elizeus ezt látta, fölkiáltott: „Atyám, atyám, Izrael szekere és fogata!” S többé nem látta. Akkor fogta ruháját, és kettérepesztette. Majd fölemelte Illés köntösét, amely leesett róla, visszafordult, és újra megállt a Jordán partján. Aztán fogta Illés köntösét, ráütött vele a vízre, és azt mondta: „Hol az Úr, Illés Istene?” Mihelyt a vízre ütött, az szétvált az egyik és a másik irányban, s Elizeus átment.

Amikor a prófétatanítványok odaátról látták, felkiáltottak: „Illés lelke leszállt Elizeusra!” S eléje mentek, a földre borultak előtte.
\end{lectio}

\parttitle[Mal 4, 5; Lk 1, 15. 17]{Válaszos ének}

\versicle{
    Nézzétek, elküldöm nektek Ill\underline{é}s próf\underline{é}tát, † mielőtt elérkeznék az Úr nagy és félelmet\underline{e}s n\underline{a}pja. * Ő újra fiaik felé fordítja az apák szívét, és apáik felé a fi\underline{a}ik sz\underline{í}vét.
}{%
    János nagy lesz az Úr \underline{e}lőtt, † Illés szellemével és erej\underline{é}vel * előtte fog járni \underline{a}z \underline{Ú}rnak.
}

\parttitle{Második olvasmány}
Jan Dlugosz krakkói kanonok „A lengyel királyság évkönyvei” című művéből

{\centering\redtext{(Históriáé Pol., t. III. p. 531)}}

\begin{lectio}{Nagyszerű erényei miatt hírnévnek és tiszteletnek örvend}
\noindent A mai napon, azaz 1399. július 17-én meghalt Hedvig királyné. Nagyon bájos arcú volt, de szokásait és erényeit tekintve még kedvesebb; a katolikus hit elterjesztője Litvániában. Ő állította fel a zsoltározók kollégiumát a krakkói székesegyházban, és két oltárt ugyanott. Ő alapította a piaszki Szűz Mária-kolostort Krakkó mellett. Ő kezdte el a szláv testvérek kolostorának felépítését. Nagyböjt és Ádvent idején vezeklőövvel és rendkívüli önmegtagadásokkal fegyelmezte testét. Bőkezű volt a szegények, özvegyek, jövevények, zarándokok, mindenféle nyomorgók és szükséget szenvedők iránt.

Nem volt benne könnyelműség, nem volt harag, nem lehetett benne gőgöt, irigységet vagy dühöt találni. Isten iránti nagy áhítatával és mérhetetlen szeretetével tűnt ki; minden világi hiúságot elutasítva magától, lelkét és gondolatvilágát egyedül az imádságra és szent könyvek, nevezetesen az Ó- és Újszövetség, a négy egyháztanító homíliái, az atyák életrajzai, prédikációk, szentek élete, Boldog Bernát, Szent Ambrus elmélkedései és prédikációi, Szent Brigitta jelenései és más, latinról lengyelre fordított könyvek olvasására fordította. Sok tehetséges, magát a tudománynak szentelő ifjú ellátásáról gondoskodott. Prágában kollégiumot alapított a litvánok számára, és fáradozott a krakkói Hittudományi Kar felállításán. Végrendeletében minden ékszerét, ruháját, pénzét és minden királyi öltözetét a szegények megsegítésére és a krakkói egyetem megújítására hagyta.

Ugyancsak a krakkói székesegyházra hagyott egy gyöngyökkel kirakott főpapi melldíszt. Olyan híres volt, és annyira megbecsülték az egész katolikus világban nagyszerű erkölcsi magatartása miatt, hogy életében a szentség példaképeként tisztelte mindenki.
\end{lectio}

\parttitle[Mt 25, 35. 40; Péld 19, 17]{Válaszos ének}

\versicle{%
    Éhes voltam, és adtatok ennem, szomjas voltam, és adt\underline{a}tok \underline{i}nnom, † idegen voltam, és befog\underline{a}dt\underline{a}tok. * Bizony, mondom nektek, amit e legkisebb testvéreim közül eggyel is tettetek, vel\underline{e}m tett\underline{é}tek.
}{%
    Aki megszánja a sz\underline{e}gényt, * az Úrnak \underline{a}d k\underline{ö}lcsönt. 
}

\end{document}